\documentclass[11pt,a4paper]{article}
\usepackage[T1]{fontenc}
\usepackage[utf8]{inputenc}
\usepackage{amsmath,amssymb,amsthm}
\usepackage[margin=25mm]{geometry}
\usepackage[hidelinks]{hyperref}
\newtheorem{definition}{Definition}
\newtheorem{proposition}{Proposition}
\newtheorem{theorem}{Theorem}
\newtheorem{remark}{Remark}
\title{OP5: A Conditional QAM Residual Lift}
\author{Working note for the Quantum Sphaera Companion}
\date{13 October 2026}
\begin{document}
\maketitle
\noindent\textbf{Status.} This note instantiates the three-component fit score
of the Companion v4.7.0, r04, and the scalar update of r08--r09. It proves a
conditional witness-lift result and an explicit residual model. It does not
identify that model with a derived MM/HC witness operator. OP5 is considered
only through this local C5 fit interface; global closure is not asserted.
The added witness-shadow theorems prove exact local residual zero at an
admissible boundary witness under compactness or a summable step bound,
continuity and an independently justified local closure characterization.
The final native-data audit distinguishes the C5 fit score, the HC/MM radial
chart certificate and the older OP5 Theta-balance defect; no identification
between these objects is assumed.

\section{Witness data and codes}
\begin{definition}[Witness domain]
Let $A$ be a specified set of admissible local witness states, with residual map
\[
R:A\to E=H_{\rm bd}\oplus H_{\rm ext}\oplus H_{\rm mix},\qquad
\Phi(x)=\|R(x)\|_\oplus^2=\sum_j\|\rho_j(x)\|^2.
\]
The component spaces are normed spaces with a positive definite sum norm, and
the components are normalized to permit a dimensionless sum. A state also
stores sidedness, scope, and separate $C4$ and $C6$ ledger coordinates.
The additional r02 slots $\rho_{\rm leak}$ and $\rho_{\rm fit}$ are retained
as separate data; their elimination or identification is not assumed.
Membership in $A$ requires independently evaluated leakage, absorption and
scope guards, not merely stored zero flags.
\end{definition}

Set $s_x=\langle0,x\rangle$, $S_A=\{s_x:x\in A\}$ and
$L=\langle1,\ell,S_A\rangle$. For a function $F:A\to A$, its QAM operator
code is
\[
o_F=\langle2,\ell,\ell,G_F\rangle,\qquad
G_F=\{(s_x,s_{Fx}):x\in A\}.
\]
The exact score readout is
\[
r_\Phi=\langle3,\ell,\mathbb R_{\ge0},
\{(s_x,\Phi(x)):x\in A\}\rangle.
\]
All tuples, graphs and values are understood as ZF-coded data. These formulas
provide functional operator and readout codes; they do not by themselves
certify native closure-class preservation.

\section{Conditional lift}
\begin{theorem}[Witness-level contraction]
Suppose a concrete function $F:A\to A$ and a residual map $T$ satisfy
\[
R(Fx)=TR(x),\qquad \|TR(x)\|_\oplus^2=\frac34\|R(x)\|_\oplus^2
\quad(x\in A).
\]
Suppose also that $F$ preserves the $C4$, $C6$, sidedness and scope
coordinates and that the actual guard evaluations remain zero on its orbit.
Then for every $x_0\in A$ with $\Phi(x_0)=1$ and every finite $n\ge0$,
\[
x_n=F^n x_0\in A,\qquad \Phi(x_n)=(3/4)^n.
\]
The ledger coordinates and guard conditions remain as specified.
\end{theorem}
\begin{proof}
The mapping property $F:A\to A$ gives all finite iterates. The commuting
residual relation and norm identity give $\Phi(Fx)=3\Phi(x)/4$.
Induction proves the formula. Applying the separately established preservation
conditions at each step proves the guard and ledger statements.
\end{proof}

\begin{remark}[QAM admissibility]
To assert $\mathrm{QAdmTransport}(o_F,s_x,s_{Fx})$, one additionally needs a
native closure-class invariant $C(Fx)=C(x)$ and a justified coding of native
class equality into $\mathrm{QClEq}$. Neither score contraction nor copying
ledger flags establishes this invariant. A readout need not be
closure-compatible merely to be evaluable as $\mathrm{QRead}$.
\end{remark}

\section{An explicit residual model}
\begin{proposition}[Amplitude lift and invariant ball]
On $E=\mathbb R^3$ with Euclidean norm, set
\[
T=\frac{\sqrt3}{2}I,\quad R_0=(1,0,0),\quad R_n=T^nR_0.
\]
Then $T^*T=3I/4$, the closed unit ball is invariant under $T$, and
\[
\|R_n\|^2=(3/4)^n>0\quad(n<\infty),\qquad R_n\longrightarrow0.
\]
Keeping separate ledger coordinates fixed preserves those model coordinates.
\end{proposition}
\begin{proof}
The scalar matrix gives $T^*T=3I/4$ and
$\|TR\|=\sqrt3\|R\|/2\le\|R\|$. Iteration gives the formula and the
positive finite values. The geometric amplitude tends to zero. An update
that is the identity on ledger coordinates preserves their projections.
\end{proof}

\begin{remark}[Model versus native witness]
Multiplying residual amplitudes by $3/4$ instead would multiply the score by
$9/16$. The displayed amplitude lift realizes the requested score ratio,
but a native witness map with $R(Fx)=TR(x)$ remains to be constructed.
Preservation of model coordinates is not a proof of native leakage absence.
\end{remark}

\begin{proposition}[A norm estimate does not imply domain invariance]
There exist a domain $A\subset\mathbb R^2$ and a linear map $T$ with
$\|TR\|^2=3\|R\|^2/4$ for every $R$, but $T(A)\not\subset A$.
\end{proposition}
\begin{proof}
Take $A=\{(u,0):|u|\le1\}$ and let $T$ be $\sqrt3/2$ times a
rotation through $90$ degrees. The rotation is isometric, but
$T(1,0)=(0,\sqrt3/2)\notin A$.
\end{proof}

\section{Detector and exact readouts}
Define $\theta=1/256$ and
\[
D_\theta(p)=\begin{cases}0_{\rm det},&0\le p\le\theta,\\
1_{\rm det},&p>\theta.\end{cases}
\]
Its QAM readout code is obtained from $r_\Phi$ by replacing the target with
$\{0_{\rm det},1_{\rm det}\}$ and the graph value with $D_\theta(\Phi(x))$.
For the model,
\[
\Phi_{19}=\frac{1162261467}{274877906944}>\theta,\qquad
0<\Phi_{20}=\frac{3486784401}{1099511627776}<\theta.
\]
Thus step $20$ is the first detector-zero step, with
$\|R_{20}\|\le1/16$ but $R_{20}\ne0$.

\begin{proposition}[The missing zero reflection]
For a specified admissible domain $A$ and independently defined target
property $Z_{\rm exact}$, detector-zero promotion is valid precisely when
\[
\{x\in A:D_\theta(\Phi(x))=0_{\rm det}\}
\subseteq\{x\in A:Z_{\rm exact}(x)\}.
\]
For $Z_{\rm exact}(x)\equiv[R(x)=0]$, this condition fails in the residual
model on any domain containing $x_{20}$.
\end{proposition}
\begin{proof}
The inclusion is exactly the quantified implication. At $x_{20}$ its premise
holds and its conclusion fails, by the computed positive score.
\end{proof}

\section{Nonstandard halo and standard-part readout}
Work in a specified proper nonstandard extension of the real superstructure,
with transfer, hypernaturals ${}^*\mathbb N$, and the standard-part map on
finite hyperreals. This is an additional foundation for this section;
it is not inferred from the finite witness ledgers. Write ${}^*\Phi$ for the
natural extension of the model score sequence.

\begin{definition}[Halo and two distinct zero readings]
The halo (monad) of zero is
\[
\operatorname{halo}(0)=
\{p\in{}^*\mathbb R: |p|<\varepsilon
\text{ for every standard real }\varepsilon>0\}.
\]
A number $p$ is an element of this halo, not the halo itself. On finite
hyperreals, $\operatorname{st}(p)=0$ is equivalent to $p\approx0$.
It is not equivalent to $p=0$.
\end{definition}

\begin{proposition}[Positive infinitesimal C5 remainder]
For every infinite $H\in{}^*\mathbb N\setminus\mathbb N$, the extended
model has
\[
p_H={}^*\Phi(H)=(3/4)^H>0,\qquad
p_H\in\operatorname{halo}(0),\qquad \operatorname{st}(p_H)=0.
\]
In the explicit residual model, with coordinatewise standard part,
\[
{}^*R_H=(\sqrt3/2)^H(1,0,0)\ne0,\qquad
\operatorname{st}({}^*R_H)=0\in\mathbb R^3.
\]
\end{proposition}
\begin{proof}
The standard statement $(3/4)^n>0$ for every $n\in\mathbb N$ transfers
to every hypernatural index, so $p_H>0$. Fix a standard $\varepsilon>0$.
Convergence of the geometric sequence gives a standard $N$ such that
$n\ge N$ implies $(3/4)^n<\varepsilon$. Transfer of this tail statement
and $H\ge N$ give $p_H<\varepsilon$. Since this holds for every standard
positive $\varepsilon$, $p_H$ is infinitesimal and its standard part is zero.
The same argument for $(\sqrt3/2)^n$ proves the vector statement.
\end{proof}

The fixed detector threshold $\theta=1/256$ is a positive standard real.
Thus ${}^*D_\theta(p_H)=0_{\rm det}$, but the converse detector-to-halo
implication fails: $p=\theta/2$ is a positive standard real that also reads
as detector-zero. In particular, $\Phi_{20}$ is not infinitesimal and
$\operatorname{st}(\Phi_{20})=\Phi_{20}>0$.

There are consequently three distinct assertions:
\[
{}^*D_\theta(p)=0_{\rm det},\qquad
\operatorname{st}(p)=0_{\rm st},\qquad p=0.
\]
Here $0_{\rm st}$ is the ordinary real zero, labelled only to record the
readout target. On nonnegative finite hyperreals, exact zero implies
standard-part zero, which implies detector-zero at any fixed positive
standard threshold. Neither converse holds in general. Equivalently,
standard-part zero requires smallness at every standard positive scale,
not just at the single scale $1/256$.

\begin{remark}[QAM boundary and transfer]
The standard-part map and the halo of zero are external relative to the
nonstandard universe. An internal hyperfinite iteration can produce
${}^*R_H$, but applying standard part is a separate external readout.
It is not an ordinary finite execution, and an infinite $H$ is not a
standard finite stage or the topological point $\infty$ used in r13.
To represent this in QAM requires a specified nonstandard model and a
separate layer for an external standard-part readout; the existing
standard-coded readout graph is not automatically an internal map by transfer.
External readouts can be described in the surrounding set-theoretic
framework, but their admissibility and native closure interpretation still
require justification.
\end{remark}

For the explicit model this halo statement is an alternative formulation of
the already proved standard limit. It does not establish a new MM/HC witness
operator. Extending a native iteration to hypernatural indices first requires
a well-defined transferable native state domain, update and guard predicates.
Passing its standard part back to an admissible native state requires further
near-standardness and closure results. Coordinatewise standard part used
above is justified in $\mathbb R^3$; it is not asserted for arbitrary
infinite-dimensional witness spaces.

\section{A local witness-to-closure bridge}
The scalar halo argument can be lifted to a genuine standard witness under
explicit geometric hypotheses. In this section the residual target is
$\mathbb R^m$ with a positive definite norm. Complex coordinates or finitely
many matrix entries may be represented by their real and imaginary parts.
Fixed discrete route, sidedness and scope tags are kept in one fixed fibre.
The iteration counter belongs to the run history, not to the geometric
coordinates whose standard part is taken.

\begin{theorem}[Compact witness shadow with exact local residual zero]
\label{thm:compact-witness-shadow}
Let $A\subset\mathbb R^d$ be a nonempty compact set of admissible witnesses,
let $F:A\to A$ be a specified function, and let $R:A\to\mathbb R^m$ be
continuous. Suppose the standard orbit $x_n=F^n(x_0)$ satisfies
\[
\|R(x_n)\|^2\le (3/4)^n\|R(x_0)\|^2.
\]
For each infinite hypernatural $H$, there is then a standard witness
\[
x_H^\circ=\operatorname{st}({}^*x_H)\in A,
\qquad R(x_H^\circ)=0.
\]
If the independently specified local closure property is
$Z_{\rm loc}(x)\Longleftrightarrow R(x)=0$, then $x_H^\circ$ is exactly
locally closed. No finite iterate is thereby declared exactly closed.
\end{theorem}
\begin{proof}
The standard orbit is a function $\mathbb N\to A$. Its natural extension
has ${}^*x_H\in{}^*A$. Compactness gives a finite coordinate vector and a
standard part $x_H^\circ\in A$. Transfer of the displayed residual estimate
gives an infinitesimal upper bound on $\|{}^*R({}^*x_H)\|^2$.
Hence ${}^*R({}^*x_H)\approx0$. By continuity at $x_H^\circ$,
\[
{}^*R({}^*x_H)\approx R(x_H^\circ).
\]
The right-hand vector is standard, so it must be exactly zero. The closure
conclusion uses the separately stated zero-set characterization.
\end{proof}

\begin{remark}[Standard counterpart and scope]
The same existence result follows from a convergent subsequence in compact
$A$ and continuity of $R$. The nonstandard formulation selects a shadow at
$H$; without orbit convergence, different infinite indices can select
different zero-residual witnesses. Neither formulation supplies the
unconstructed native map $F$ or proves that the three-component score
detects every native OP5 obstruction. That zero-set characterization must
account for the additional fit and leakage slots before it is used.
\end{remark}

\subsection{Closed guards pass to the witness shadow}
For a finite list of continuous guard functions, constraints of the form
\[
g_i(x)=0,\qquad h_j(x)\ge0
\]
define a closed feasible set. If these constraints hold on the orbit, they
hold at every witness shadow furnished by
Theorem~\ref{thm:compact-witness-shadow}. This follows either from continuity
and taking standard parts, or from closedness of the feasible set. A strict
guard $h(x)>0$ requires a uniform standard margin $h(x_n)\ge\delta>0$ if
the strict inequality is to survive the shadow operation.
For instance, $x_n=(\sqrt3/2)^n>0$ has shadow zero: positivity alone is
not a preserved strict margin.

Separate ledger projections remain fixed at the shadow when their values
are constant on the orbit and their fibres are closed. This is a conditional
statement about actual continuous guard evaluations and fixed tags, not
about freely copied certificate flags.

\subsection{A canonical shadow from a summable step bound}
\begin{theorem}[Witness convergence controlled by the residual]
\label{thm:canonical-witness-shadow}
Let $A\subset\mathbb R^d$ be closed, $F:A\to A$, and $R:A\to\mathbb R^m$
be continuous. Set $a=\sqrt3/2$. Suppose the orbit has
\[
\|R(x_n)\|\le a^n\|R(x_0)\|,
\qquad \|x_{n+1}-x_n\|\le B\|R(x_n)\|
\]
with a standard constant $B\ge0$. Then it converges to a unique orbit
limit $x_\infty\in A$, with
\[
\|x_n-x_\infty\|
\le\frac{B\|R(x_0)\|}{1-a}\,a^n,
\qquad R(x_\infty)=0.
\]
For every infinite $H$, $\operatorname{st}({}^*x_H)=x_\infty$.
If $F$ is continuous at $x_\infty$, it also satisfies
$F(x_\infty)=x_\infty$.
\end{theorem}
\begin{proof}
For $k>n$, the triangle inequality and the geometric sum give
\[
\|x_k-x_n\|
\le\sum_{j=n}^{k-1} B\|R(x_0)\|a^j
\le\frac{B\|R(x_0)\|}{1-a}a^n.
\]
Thus $(x_n)$ is Cauchy. Completeness of $\mathbb R^d$ and closedness of
$A$ give its unique limit in $A$. Passing $k$ to the limit gives the tail
bound. Continuity of $R$ and residual decay give $R(x_\infty)=0$.
Transfer of the tail bound gives ${}^*x_H\approx x_\infty$ for every
infinite $H$. Finally, continuity of $F$ gives
$F(x_\infty)=\lim_n F(x_n)=\lim_n x_{n+1}=x_\infty$.
\end{proof}

\begin{proposition}[Full verification in the residual model]
Both witness-shadow theorems apply to $A=\{x\in\mathbb R^3:\|x\|\le1\}$,
$F(x)=ax$ and $R(x)=x$. The step bound holds with $B=1-a$.
The unique shadow and fixed point is $x_\infty=0$, the residual vanishes
exactly there, and any additional fixed ledger coordinates remain fixed.
\end{proposition}
\begin{proof}
The ball is compact, $F$ leaves it invariant and is continuous, and
\[
\|R(Fx)\|=a\|R(x)\|,\qquad
\|Fx-x\|=(1-a)\|R(x)\|.
\]
The theorems therefore give the claimed limit; directly $F(0)=0$.
\end{proof}

\subsection{QAM certificate at the boundary}
When the canonical theorem holds for every initial witness in a specified
domain $A_0\subseteq A$, define the standard boundary map
\[
\mathcal B:A_0\to Z_R,\qquad
\mathcal B(x)=\lim_{n\to\infty}F^n(x),\qquad
Z_R=\{y\in A:R(y)=0\}.
\]
It is independent of the infinite index used to realize its nonstandard
shadow. With source and target layers $\ell_A,\ell_Z$, its standard QAM
operator graph is
\[
o_{\mathcal B}=\langle2,\ell_A,\ell_Z,
\{(s_x,s_{\mathcal B(x)}):x\in A_0\}\rangle.
\]
This graph describes a standard boundary function whose existence was proved;
it does not claim finite executability or make the external standard-part
operation internal. A QWitness payload can store the theorem hypotheses,
their verified instantiation, the boundary witness and the local scope.
If $Z_R$ is natively characterized as the exact local closure set, this is
an exact local boundary certificate. A certificate for global/all-readout
closure still requires its own zero-set characterization and coverage proof.

For $\mathrm{QAdmTransport}$ of this boundary graph one further needs
$\mathrm{QClEq}(s_x,s_{\mathcal B(x)})$. One sufficient condition is that
the native closure-class fibre containing the orbit is closed and contains
all its iterates. Finite-step class preservation alone does not imply this
limit property. In the illustrative direction-class model, the nonzero ray
does not contain its zero limit, so that model does not justify such a
QAM admissibility claim at the boundary.

\paragraph{Current native proof obligations.}
The bridge now has an explicit proof rather than an unspecified standard-part
identification. Applying it to the Companion still requires the actual
finite-coordinate witness domain, native residual functions and their
continuity, a uniformly admissible native iteration, and either a compact
orbit closure inside the admissible domain or the summable step estimate.
The guards must define a closed feasible region, with positive margins for
strict exclusions. The native local closure condition must be exactly
characterized by the retained residuals. These are construction-specific
conditions to check, not consequences of the scalar score table.

\section{Structural test of a residual film over spacetime}
The image of a thin film over spacetime can be tested as a geometric
construction. The spacetime base is an input, not a consequence of C5
contraction. Let $M$ be a smooth Lorentzian manifold and let
$\pi:E\to M$ be a finite-rank real vector bundle for residual values.
Use a positive definite auxiliary fibre metric to measure residual size;
the indefinite spacetime metric does not itself define a positive fit norm.
No extra spacetime dimension or material thickness is assumed.

\subsection{Fibrewise halo and chart compatibility}
In local trivializations, residual sections have coordinate representatives
$r_i:U_i\to\mathbb R^m$ with
\[
r_i(x)=g_{ij}(x)r_j(x),\qquad
g_{ij}(x)g_{jk}(x)=g_{ik}(x),
\]
where the transition matrices are smooth and invertible.
At a standard base point $x$, define the residual halo in the extended fibre
as the vectors $v\in{}^*E_x$ with infinitesimal auxiliary norm.
For a standard continuous transition matrix, its entries at a standard
point are finite, so it sends infinitesimal vectors to infinitesimal vectors.
The inverse gives the reverse implication. On compact chart overlaps both
maps are bounded, and the same conclusion holds uniformly there.
Thus halo membership does not depend on the chosen regular local frame.
An unlimited, nonstandard frame rescaling would require a separate check.

The collection of these halos is an external fibrewise neighbourhood of the
zero section. It is not asserted to be a new internal vector bundle or a
material layer. A particular residual field is a section taking values in
these halos, not the collection of halos itself.

\begin{proposition}[A consistent spacetime residual-film model]
Let $s_0$ be a smooth standard section of $E$, set $a=\sqrt3/2$, and define
the sequence of smooth sections $s_n=a^n s_0$. For any infinite hypernatural
$H$, let $s_H={}^*s(H,\cdot)$ be the extended sequence at $H$. Then at every
standard base point $x$,
\[
s_H(x)=a^H s_0(x)\approx0,\qquad
\operatorname{st}(s_H(x))=0_x.
\]
Where $s_0(x)\ne0$, the hyperreal field value is nevertheless nonzero.
On each fixed compact set $K\subset M$,
\[
\sup_{x\in{}^*K}\|s_H(x)\|
\le a^H C_K\approx0,
\qquad C_K=\sup_{x\in K}\|s_0(x)\|<\infty.
\]
The field representatives glue consistently on overlaps and the pointwise
score satisfies $\|s_n(x)\|^2=(3/4)^n\|s_0(x)\|^2$.
\end{proposition}
\begin{proof}
Scalar multiplication is defined fibrewise and commutes with every linear
transition map: $a^n g_{ij}s_{0,j}=g_{ij}a^n s_{0,j}$. Thus every $s_n$ is
a section and the overlap equations transfer to $H$. The scalar $a^H$ is
positive infinitesimal. Multiplication by a fixed finite standard vector
gives the pointwise assertion. Smoothness and compactness bound the standard
initial section on $K$, and transfer gives the displayed uniform estimate.
Homogeneity of the fibre norm proves the score identity.
\end{proof}

In particular, on the trivial bundle $M\times\mathbb R^3$ with
$s_0(x)=(1,0,0)$, this constructs an everywhere nonzero infinitesimal
residual section with identically zero standard-part readout. This is a
concrete structural realization of the film image. On a general nontrivial
bundle, a nowhere-zero section is an additional existence condition;
the construction does not remove topological obstructions.

\subsection{Derivative and observable tests}
Fix a smooth standard connection $\nabla$ on $E$. For each standard finite
derivative order $k$ and each compact $K$, smoothness gives a finite bound
on $\nabla^k s_0$ and
\[
{}^*\nabla^k s_H=a^H\,{}^*(\nabla^k s_0)\approx0
\quad\text{uniformly on }{}^*K.
\]
Consequently regular standard local observables continuous in a fixed finite
jet of this residual field have the same standard part as their zero-residual
values at standard points. An observable Lipschitz in the field, with a fixed
finite bound on $K$, has infinitesimal change uniformly there. Integration
of such a bounded change over a fixed compact region of finite standard
volume also has infinitesimal change. These claims do not cover arbitrary
singular observables, noncompact integration regions or unlimited gains.

Infinitesimal amplitudes alone are insufficient. In a dimensionless local
coordinate $t$, for a positive infinitesimal $\varepsilon$, the internal
smooth field
\[
u_\varepsilon(t)=\varepsilon\sin(t/\varepsilon)
\]
has infinitesimal amplitude, but
\[
\partial_t u_\varepsilon(t)=\cos(t/\varepsilon),\qquad
\partial_t u_\varepsilon(0)=1.
\]
Thus taking standard part and differentiating need not commute for a
general internal family. A derivative-dependent energy term can remain
appreciable. The specific fixed-profile construction above excludes this
counterexample by scaling a single standard smooth profile, including its
derivatives. Likewise, rescaling $s_H$ by the unlimited gain $a^{-H}$
recovers $s_0$; this is a new readout scale, not a standard bounded response.

\subsection{Relation to the Companion and QAM}
The Companion's legacy probe--closure field block, labelled
\texttt{sec:field-eq}, already introduces spacetime-dependent probe,
connection, closure-projector and compensator fields. This gives a possible
host for a residual bundle, but does not identify our new section with any
of those fields. A native identification requires the actual tensor or
representation type of each residual, compatible transition maps, and a
pointwise or locally uniform witness update that realizes the contraction.

In QAM, local state payloads would carry the base chart, residual section,
transition data, guard witnesses and declared observable class. The overlap
equations must be proved before local readouts can be glued. The standard-part
readout then sends the infinitesimal residual section to the zero section
over the same standard base. This local geometric gluing does not certify
global/all-readout closure of the native OP5 architecture.

\paragraph{Structural verdict.}
The film assumption is consistent and constructible as an infinitesimal
residual section over a supplied spacetime, with coordinate-compatible zero
standard readout. It does not yet establish a new physical substance, an
extra spacetime sheet, friction, stress-energy or a changed spacetime metric.
Those interpretations require specified dynamics and observables. The
construction is an extension of the tested residual model, not a derivation
of a native MM/HC spacetime field.

\section{Navier--Stokes test of the residual film}
\label{sec:fluid-test}
We now assign fluid variables to the residual field as an additional model.
The assignment is not inferred from its infinitesimal size. Work in a chosen
nonrelativistic spatial chart with time $t$, on a periodic cube of side $L$.
This uses a spatial velocity field, not a Lorentz-covariant fluid theory on
an arbitrary curved spacetime. The fluid is Newtonian and incompressible,
with standard constant density $\rho_0>0$ and kinematic viscosity $\nu>0$.
The dynamic viscosity is $\mu=\rho_0\nu$.

The residual-to-velocity map must supply a spatial vector with velocity units.
For a normalized three-component residual $r$, a local model assignment is
$u=U\mathcal I(r)$, with standard speed $U$ and a specified linear
identification $\mathcal I$ into the spatial tangent space. Divergence-free
velocity and compatible frame transitions are extra constraints, not
properties of every QAM residual. Pressure is denoted $P$ to distinguish it
from the earlier scalar residual $p$.

\subsection{Equations and scale test}
With body-force acceleration $f$, the equations are
\[
\partial_t u+(u\cdot\nabla)u
=-\frac1{\rho_0}\nabla P+\nu\Delta u+f,
\qquad \nabla\cdot u=0.
\]
For a residual-scale velocity $u=\varepsilon U v$, with positive
infinitesimal $\varepsilon$ and a regular dimensionless profile $v$,
the Reynolds number based on standard length $L$ is
\[
\mathrm{Re}_{\rm res}=\frac{\varepsilon UL}{\nu}\approx0.
\]
On the viscous time scale $L^2/\nu$, the dimensionless advection term is
multiplied by this Reynolds number. Provided derivatives of the profile
remain bounded, the leading equation is unsteady Stokes flow; dropping
advection is not a general exact identity.
Infinitesimal viscosity or gradients on infinitesimal scales need a separate
scaling analysis. Rescaling an arbitrary nonlinear Navier--Stokes solution
by an amplitude alone need not yield another solution, since advection is
quadratic while the time derivative and viscous term are linear.

\subsection{An exact viscous shear solution}
Choose $k=2\pi m/L$ for an integer $m\ge1$. Let $\eta>0$ be an amplitude
parameter, which may be standard or, by transfer, hyperreal. Set
\[
u_\eta(x,y,z,t)
=\eta U e^{-\nu k^2t}\sin(ky)\,e_x,
\qquad P=P_0,\qquad f=0.
\]
Here $P_0$ is spatially constant. The periodic geometry has no free surface
and no separately prescribed film thickness.

\begin{proposition}[Exact incompressible residual flow]
The displayed velocity and pressure solve the incompressible Navier--Stokes
equations for every $t\ge0$. With $\eta=\varepsilon$ positive infinitesimal,
the velocity is infinitesimal and nonzero wherever $\sin(ky)\ne0$, at every
standard finite time. Its standard velocity readout is identically zero.
\end{proposition}
\begin{proof}
The only nonzero component is independent of $x$, hence
\[
\nabla\cdot u_\eta=0,\qquad
(u_\eta\cdot\nabla)u_\eta=u_{\eta,x}\partial_xu_\eta=0.
\]
Moreover,
\[
\partial_t u_\eta=-\nu k^2u_\eta,\qquad
\Delta u_\eta=-k^2u_\eta,\qquad \nabla P=0.
\]
These identities verify the equation exactly, not just to first order.
For standard finite $t$, the exponential and profile factors are finite;
where the sine is nonzero they are positive in magnitude. Multiplication by
positive infinitesimal $\varepsilon$ gives the final assertions.
\end{proof}

\subsection{Energy, stress and the C5 sampling ratio}
Let $\langle\cdot\rangle$ denote the spatial mean on the periodic cube.
The kinetic energy per unit volume and its viscous dissipation satisfy
\[
\mathcal K_\eta(t)=\frac{\rho_0}{2}\langle|u_\eta|^2\rangle
=\frac{\rho_0\eta^2U^2}{4}e^{-2\nu k^2t},
\]
\[
\frac{d\mathcal K_\eta}{dt}
=-\rho_0\nu\langle|\nabla u_\eta|^2\rangle
=-\frac{\rho_0\nu k^2\eta^2U^2}{2}e^{-2\nu k^2t}.
\]
Thus energy and dissipation are of order $\varepsilon^2$ for
$\eta=\varepsilon$. The shear stress is
\[
\tau_{xy}=\mu\partial_yu_{\eta,x}
=\mu k\eta U e^{-\nu k^2t}\cos(ky),
\]
of order $\varepsilon$ at standard fixed $k,\mu,t$. Relative to the
same fluid at rest, the standard parts of these velocity-dependent changes
vanish. The assumed standard density and background pressure themselves
do not vanish; this model describes a residual motion of a supplied fluid,
not the creation of a new fluid substance from the halo.

For the standard-amplitude solution $\eta=1$, choose
\[
\Delta t=\frac{\log(4/3)}{2\nu k^2},\qquad t_n=n\Delta t.
\]
Then the normalized mean-square-speed score obeys
\[
\Phi_{\rm fluid}(t_n)
:=\frac{\langle|u_1(t_n)|^2\rangle}
{\langle|u_1(0)|^2\rangle}
=e^{-2\nu k^2 n\Delta t}=(3/4)^n.
\]
The velocity amplitude ratio is $\sqrt3/2$ per sampled step. This gives
an explicit fluid realization of the tested C5 score recurrence, with
preserved incompressibility and a bounded smooth profile. At infinite
hypernatural $H$, $u_1(t_H)$ has zero standard part but can remain nonzero.
Here $t_H=H\Delta t$ is an unlimited time; it is not a finite-time exact
closure. The matching ratio fixes only the combination $\nu k^2\Delta t$,
not native viscosity, density, wavelength or an OP5 physical time unit.

\subsection{What changes for a literal thin layer}
If ``film'' also means a material thickness $h$, a different geometry and
boundary conditions must be supplied. In a flat channel $0<y<h$ with
stationary no-slip walls, the same shear construction works with
$\sin(\pi y/h)$ and $k=\pi/h$. Its damping time is
\[
t_{\rm damp}=\frac{h^2}{\nu\pi^2}.
\]
For infinitesimal thickness $h=\varepsilon h_0$ but standard $\nu$,
this is an infinitesimal time. A velocity amplitude $\varepsilon U$ has
an initial transverse derivative of order $U/h_0$, which need not be
infinitesimal. Thus an infinitesimal-amplitude field and an infinitesimally
thin material layer are distinct assumptions. A free-surface film would
add moving-boundary and stress-balance equations; they are not included in
the periodic or fixed-channel model.

\paragraph{Fluid-test result and native interface.}
An incompressible viscous residual flow with zero standard velocity readout
exists explicitly, and its sampled energy can reproduce the $3/4$ ratio.
For fixed standard geometric and viscosity scales it is a low-Reynolds
residual motion. The calculation establishes this particular shear model,
not a turbulence classification or a claim about all three-dimensional
Navier--Stokes solutions. To turn it into a native MM/HC fluid model,
the residual-to-velocity map, conservation of mass and momentum, pressure
constraint, constitutive viscosity law, boundary data and native guard
preservation must be derived or independently specified.
QAM can store these quantities and the exact shear certificate; the
certificate scope is this declared fluid model, not native all-readout closure.

\section{General OP5 transport--dissipation candidate}
\label{sec:general-op5-transport}
This section replaces the special shear ansatz by a proposed multichannel
residual evolution and derives a conditional energy certificate. It does
not assert that the native MM/HC witness map already generates this equation.
All spatial and time variables here are normalized using fixed standard
scales. The domain is the periodic cube $\Omega=\mathbb T_L^3$.

\subsection{State, coefficients and evolution}
Write the retained C5 residuals as one real coordinate vector
\[
r=(r_{\rm bd},r_{\rm ext},r_{\rm mix})\in\mathbb R^m,
\qquad m=m_{\rm bd}+m_{\rm ext}+m_{\rm mix}.
\]
The three channel labels are not automatically three spatial velocity
components. Complex or tensor residuals require a specified real coordinate
representation and positive definite fit metric. In the present model the
channel norm is the Euclidean direct-sum norm. On fields, define
\[
\mathcal E(t)=\frac12\int_\Omega |r(x,t)|^2\,dx,
\qquad \Phi_{\rm tr}(t)=2\mathcal E(t).
\]
A claim that this integral equals the native finite witness score requires
an explicit spatial representation and normalization; it is not made solely
because both quantities are sums of squared residuals.

The proposed evolution is
\begin{equation}
\label{eq:general-residual-transport}
\partial_t r+b\cdot\nabla r
=\sum_{j=1}^3\partial_j(D\partial_jr)-Kr-Br
+\mathcal N(r)+s.
\end{equation}
Here $b$ is a smooth divergence-free transport velocity, $D=D^{\mathsf T}$
is a constant channel-diffusion matrix with $D\ge d_0I$, $d_0>0$,
$K^{\mathsf T}=-K$ is a pointwise skew channel coupling, and
$B=B^{\mathsf T}\ge b_0I$, $b_0\ge0$, is a damping matrix.
The source $s$ is external or supplied by other native channels.
The nonlinear contribution is required to satisfy an explicit energy bound
\[
\int_\Omega r\cdot\mathcal N(r)\,dx
\le\sigma\|r\|_{L^2}^2
\]
with a standard constant $\sigma\ge0$. This is a hypothesis on the chosen
nonlinearity, not a property of every possible OP5 interaction.
The coefficients may couple the three retained channels, but any coupling
to excluded or guarded native channels requires the separate guard test below.

\begin{theorem}[Multichannel residual energy balance]
\label{thm:transport-energy}
On any time interval where the proposed equation has a sufficiently smooth
periodic solution and the stated coefficient conditions hold,
\begin{align*}
\frac{d\mathcal E}{dt}
&+\sum_{j=1}^3\int_\Omega(\partial_jr)^{\mathsf T}D(\partial_jr)\,dx
+\int_\Omega r^{\mathsf T}Br\,dx\\
&=\int_\Omega r\cdot\mathcal N(r)\,dx
+\int_\Omega r\cdot s\,dx.
\end{align*}
In particular, divergence-free advection and the skew channel coupling
contribute no net residual energy.
\end{theorem}
\begin{proof}
Multiply Equation~\eqref{eq:general-residual-transport} by $r$ and integrate.
The transport term is the integral of $b\cdot\nabla(|r|^2/2)$, which is zero
by periodicity and $\nabla\cdot b=0$. Integration by parts gives the positive
diffusion term on the left. Pointwise $r^{\mathsf T}Kr=0$. The remaining
terms give the displayed identity.
\end{proof}

\subsection{An exact transported, channel-mixing test}
For a periodic coordinate $x$ and compatible wavenumber $k$, take
$b=(c,0,0)$, $D=\nu I$, $B=\beta I$ with $\beta\ge0$, and
\[
K=\begin{pmatrix}0&\omega&0\\-\omega&0&0\\0&0&0\end{pmatrix},
\qquad \mathcal N=s=0.
\]
Then
\[
r(x,t)=e^{-(\nu k^2+\beta)t}\sin(k(x-ct))
\begin{pmatrix}\cos(\omega t)\\\sin(\omega t)\\0\end{pmatrix}
\]
is an exact solution of Equation~\eqref{eq:general-residual-transport}.
Translation accounts for the transport term, the rotating channel vector
has derivative $-Kr$, and diffusion contributes $-\nu k^2r$.
Thus there is real spatial transport and skew channel mixing, yet the
integrated score decays exactly as $e^{-2(\nu k^2+\beta)t}$.
This is a residual transport example with prescribed $b$, not automatically
a self-advecting Navier--Stokes velocity field. Its equation residual and
norm identity were checked by symbolic differentiation.

\subsection{Decay, zero modes and the sampled C5 bound}
If a Poincare estimate $\|\nabla r\|_{L^2}^2\ge
\lambda\|r\|_{L^2}^2$ holds throughout the evolution, set
\[
\alpha=d_0\lambda+b_0-\sigma.
\]
For mean-zero periodic fields one may take $\lambda=(2\pi/L)^2$.
Mean-zero must itself be invariant under the proposed coefficients,
nonlinearity and source; otherwise this positive gap cannot be used.
Without such a constraint take $\lambda=0$, and obtain decay only if the
reaction damping still gives $b_0>\sigma$.

\begin{proposition}[Unforced decay and finite sampling certificate]
\label{prop:general-sampled-decay}
Assume $s=0$ and $\alpha>0$. Then on the solution interval,
\[
\Phi_{\rm tr}(t)\le e^{-2\alpha t}\Phi_{\rm tr}(0).
\]
For $t_n=n\Delta t$ with $\Delta t=\log(4/3)/(2\alpha)$,
\[
\Phi_{\rm tr}(t_n)\le(3/4)^n\Phi_{\rm tr}(0).
\]
With initial score one, the detector threshold $1/256$ is reached by step
$20$, provided the solution and all guard certificates exist to that time.
Equality and the first hitting step need not match the shear case.
\end{proposition}
\begin{proof}
The energy identity and coefficient bounds give
$\mathcal E'+2\alpha\mathcal E\le0$. Integration and the definition
of $\Phi_{\rm tr}$ give the first inequality; substitution of $\Delta t$
gives the sampled bound. Finally $(3/4)^{20}<1/256$.
\end{proof}

Diffusion alone does not remove the spatially constant mode. For example,
$r(x,t)=c\ne0$, with $B=K=\mathcal N=s=0$, is an exact stationary solution
for any divergence-free $b$. This is a direct obstruction to general
zero-residual closure without a positive damping gap or an excluded mean mode.

\subsection{Sources can leave a nonzero collar}
With $\alpha>0$ and $\|s(t)\|_{L^2}\le S$, the energy balance gives
\[
\mathcal E'+2\alpha\mathcal E
\le\sqrt{2\mathcal E}\,S
\le\alpha\mathcal E+\frac{S^2}{2\alpha}.
\]
Thus
\[
\Phi_{\rm tr}(t)
\le e^{-\alpha t}\Phi_{\rm tr}(0)
+\frac{S^2}{\alpha^2}(1-e^{-\alpha t}).
\]
The upper bound allows a persistent forced collar. A concrete example is
$b=K=\mathcal N=0$, $B=\alpha I$, and a spatially constant source
$s=s_0$: the constant steady residual $r=s_0/\alpha$ solves the equation.
Hence a positive standard forcing can genuinely prevent exact closure.
Infinitesimal forcing with a positive standard gap permits an infinitesimal
energy collar, but neither a fixed detector value nor a generic source bound
proves native pointwise or global closure.

\subsection{Precise Navier--Stokes reduction}
The transport system becomes incompressible Navier--Stokes only after
identifying the residual coordinates with a spatial velocity $u$ and imposing
\[
m=3,\quad r=u,\quad b=u,\quad D=\nu I,
\quad K=B=\mathcal N=0,
\quad s=-\rho_0^{-1}\nabla P+f,
\quad \nabla\cdot u=0,
\]
with the appropriate dimensional normalization restored. Pressure enforces
the divergence constraint; on the periodic domain it satisfies
\[
\Delta P=-\rho_0\sum_{i,j=1}^3
(\partial_i u_j)(\partial_j u_i)+\rho_0\nabla\cdot f.
\]
For smooth divergence-free periodic fields the pressure contributes zero
to the energy balance, although it can be nonzero pointwise. Unforced
mean-zero Navier--Stokes therefore gives
\[
\frac12\frac{d}{dt}\|u\|_{L^2}^2
+\nu\|\nabla u\|_{L^2}^2=0,
\qquad \alpha=\nu(2\pi/L)^2.
\]
This handles the nonlinear transport in the energy balance without requiring
it to vanish pointwise. The previous shear mode attains an exact exponential
rate $\nu k^2$ because it has one spatial frequency and zero self-advection.
The general energy estimate is conditional on solution regularity or on a
justified weak energy inequality. It does not prove arbitrary global smooth
three-dimensional solutions. Spatial channel labels alone do not justify
the identification $r=u$ or the pressure projection.

\subsection{C4/C6 guard invariance as an operator test}
The energy estimate controls C5 only. To make guard preservation testable,
augment the state by deviations $h$ of the C4 ledger and C6 blocker from
their specified baseline states; a nonzero inherited C4 residual is not
set to zero by this notation. Suppose their actual evolution is
\[
\partial_t h=\mathcal L_g h+\mathcal C r+\mathcal F_g(r,h),
\qquad h(0)=0.
\]
If $\mathcal C r=0$ for every allowed C5 state, $\mathcal F_g(r,0)=0$,
and the augmented evolution has uniqueness, then $h(t)=0$ is invariant.
Indeed it is a solution of the guard equation with the specified initial
data, and uniqueness excludes departure. In projection notation, the
linear no-leakage condition is $P_g\mathcal L P_5=0$ on the declared domain.
These equalities must be proved for the actual native operators; placing
zero in a stored guard coordinate does not verify them. Other strict guard
constraints still require positive margins, and scope is part of the state.

\subsection{QAM flow certificate and the halo boundary}
For a well-posed declared evolution let $\mathcal S_t$ denote its solution
map. The sampled QAM operator is the graph of $\mathcal S_{\Delta t}$ on
the independently verified admissible state domain. A witness payload records
the energy identity, the gap $\alpha$, source and boundary conditions, the
guard-invariance test and the time interval of validity. Closure-class
preservation is still a separate requirement for QAdmTransport.

The pure residual model may first be tested by a fixed, energy-preserving
Fourier--Galerkin projection on a finite number of modes. With prescribed
smooth transport, $\mathcal N=0$, and no source, the resulting finite linear
system inherits the energy estimate by testing with the projected residual.
Its bounded finite-dimensional orbit can support the earlier compact-witness
shadow argument when the actual admissible set is closed and invariant.
Extending that argument to the full PDE needs function-space compactness or
convergence; an $L^2$ energy bound alone is not finite-dimensional compactness.

If the evolution exists for all sampled times and the unforced positive-gap
estimate is uniform, transfer yields infinitesimal $L^2$ residual size at an
infinite index. This is an energy-space statement, not automatic uniform
pointwise halo membership or a derivative estimate. For example a smooth
fixed-height bump of shrinking support has arbitrarily small $L^2$ norm
while retaining height one. A uniform higher-regularity estimate or another
pointwise control is needed to recover the spacetime-film conclusion.

\paragraph{General-model outcome.}
The shear construction is now one exact specialization of a broader
transport--diffusion--reaction candidate. Its general positive result is a
conditional energy-decay certificate with a derived sampled $3/4$ bound.
Its explicit failure modes are undamped zero modes, persistent forcing,
non-invariant guards and insufficient pointwise or regularity control.
The native next step is to extract $b,D,K,B,\mathcal N,s$ and the spatial
score representation from the MM/HC witness data, then verify the spectral
gap, guard block conditions and closure-zero characterization. These
coefficients are proposed mathematical roles, not measured or derived native
parameters.

\section{Structural derivation from the MM/HC witness interface}
\label{sec:structural-native-derivation}
We now derive a candidate operator architecture from witness residuals rather
than postulate a fluid velocity at the outset. The source anchors are the
v4.7.0 r01 sided MM/HC channel split, the r02 residual inventory, the r04
sum-of-squares fit score, and the r27--r31 local constraint and penalty
packaging. Those latter functionals are audit objects, not physical actions.
The construction below introduces an explicit residual-reduction dynamics;
it does not reclassify those source objects as physical laws.

\subsection{Witness dynamics must descend to residual dynamics}
Let $W$ be a finite-dimensional smooth witness chart with $X\in W$, and
let the actual retained residual map be $R:W\to\mathbb R^m$. Denote its
Jacobian by $J_R(X)=DR_X$. A specified witness evolution
\[
\dot X=V_{\rm MM}(X)+V_{\rm HC}(X)+V_{\rm src}(X)
\]
gives by the chain rule
\[
\dot r=J_R(X)V_{\rm MM}(X)+J_R(X)V_{\rm HC}(X)
+J_R(X)V_{\rm src}(X),\qquad r=R(X).
\]
This is already an exact residual equation along a given witness orbit.
For an autonomous equation in $r$ alone, the right-hand side must have
the same value at any two witnesses $X,Y$ with $R(X)=R(Y)$.
That fibre-constancy condition is necessary and sufficient for a well-defined
vector field on the residual image: define its value using any representative;
the condition gives independence of that choice. Regularity of the descended
field needs an additional smooth local quotient or section.
Without this condition the witness variables must remain in the state.
A scalar score is an even coarser readout and cannot determine these
channel dynamics by itself.

\subsection{A guard-tangent compensator from the C5 score}
Suppose the independently evaluated equality guards define a regular local
constraint manifold
\[
W_g=\{X:C(X)=0\},\qquad A_C(X)=DC_X.
\]
With full row rank of $A_C$ in this chart, the Euclidean tangent projector is
\[
\Pi(X)=I-A_C(X)^{\mathsf T}
\bigl(A_C(X)A_C(X)^{\mathsf T}\bigr)^{-1}A_C(X).
\]
Take a symmetric positive semidefinite mobility $Q(X)$ and define the
local C5 fit objective and the proposed compensator step by
\[
\mathcal E_5(X)=\tfrac12\|R(X)\|^2,
\qquad V_{\rm HC}(X)=-\Pi Q\Pi J_R^{\mathsf T}R(X).
\]
This is a new guard-constrained fit-relaxation prescription motivated by
the sum-of-squares score. Identifying it with the actual Heisenberg
compensator requires a native operator calculation. It is not obtained by
varying a physical action or by assigning dynamics to the r31 audit functional.

\begin{proposition}[Derived positive residual damping and guard tangency]
\label{prop:projected-hc-damping}
For the proposed compensator evolution on $W_g$, define
\[
B_{\rm HC}(X)=J_R(X)\Pi(X)Q(X)\Pi(X)J_R(X)^{\mathsf T}.
\]
Then $B_{\rm HC}$ is symmetric positive semidefinite,
\[
\dot r=-B_{\rm HC}(X)r,\qquad
\frac{d\mathcal E_5}{dt}=-r^{\mathsf T}B_{\rm HC}(X)r\le0,
\qquad DC_X V_{\rm HC}(X)=0.
\]
Equality guards remain satisfied on a smooth solution within this regular
chart, provided the initial guards hold. Other witness contributions must
also be tangent for the full evolution to preserve them.
\end{proposition}
\begin{proof}
The chain rule gives the residual equation. For any $v$,
\[
v^{\mathsf T}B_{\rm HC}v
=(\Pi J_R^{\mathsf T}v)^{\mathsf T}Q(\Pi J_R^{\mathsf T}v)\ge0.
\]
Differentiating $\|r\|^2/2$ proves the energy formula. The projector has
$A_C\Pi=0$, so the guard derivative is zero. Integrating it along the
solution proves preservation of the equality constraints.
\end{proof}

Positive definiteness is a further test, not an automatic benefit of the
compensator. If $J_R\Pi Q^{1/2}$ has full row rank and its smallest singular value is
$J_R\Pi Q^{1/2}$ is uniformly at least $\sqrt\gamma>0$, then
$B_{\rm HC}\ge\gamma I$. A rank deficiency leaves residual directions
undamped. In particular, if $J_R\Pi=0$, no allowed tangent correction can
change the residual. Strict inequality guards require a positive-margin
or inward-pointing argument beyond this equality-constraint projector.

\subsection{MM transport splits into rotation and energy-changing parts}
Suppose the MM residual vector field is well defined and smooth near a
reference residual zero. Its expansion is
\[
f_{\rm MM}(r)=f_{\rm MM}(0)+A_{\rm MM}r+\mathcal N_{\rm MM}(r),
\qquad \mathcal N_{\rm MM}(r)=O(\|r\|^2).
\]
In positively normalized residual coordinates write
\[
A_-:=\tfrac12(A_{\rm MM}-A_{\rm MM}^{\mathsf T}),\qquad
A_+:=\tfrac12(A_{\rm MM}+A_{\rm MM}^{\mathsf T}).
\]
Then $A_-$ rotates or exchanges channel content without changing the
quadratic norm, while $A_+$ contributes $r^{\mathsf T}A_+r$ to its rate.
With $K=-A_-$ and the projected compensator,
\[
\dot r=-Kr-\bigl(B_{\rm HC}(X)-A_+\bigr)r
+\mathcal N_{\rm MM}(r)+s_R.
\]
Here $s_R$ includes $f_{\rm MM}(0)$ and the projected witness source.
This derives the skew-coupling, reaction-damping, nonlinear and source
roles of the general transport candidate, conditionally on the specified
witness fields and residual descent. The net damping is
$B_{\rm eff}=B_{\rm HC}-A_+$ and can fail to be positive.
For instance, $B_{\rm HC}\ge\gamma I$ and
$\lambda_{\max}(A_+)\le\eta_+$ give
$B_{\rm eff}\ge(\gamma-\eta_+)I$.
The MM contribution cannot be declared energy-neutral merely by retaining
its skew part and discarding its symmetric part.

\subsection{From certified local routes to spatial diffusion}
Spatial structure requires another explicit input: a finite collection of
local witness cells, a declared adjacency graph and a spatial realization
of its edges. QAM route words alone need not supply distances, volume weights
or a continuum mesh. Let $r_i\in\mathbb R^m$ be residuals at cells, with
positive cell volumes $v_i$ and the weighted inner product
\[
\langle r,z\rangle_V=\sum_i v_i r_i^{\mathsf T}z_i.
\]
Give every allowed edge $e=(i,j)$ a symmetric positive semidefinite
conductance matrix $D_e$ and introduce the exchange energy
\[
\mathcal E_{\rm edge}(r)
=\tfrac12\sum_{e=(i,j)}(r_i-r_j)^{\mathsf T}D_e(r_i-r_j).
\]
Negative weighted gradient flow of this energy is
\[
\dot r_i\big|_{\rm exchange}
=-\frac1{v_i}\sum_{j\sim i}D_{ij}(r_i-r_j)
=-(\mathcal L_Dr)_i.
\]
Direct summation gives
\[
\langle r,\mathcal L_Dr\rangle_V
=\sum_{e=(i,j)}(r_i-r_j)^{\mathsf T}D_e(r_i-r_j)\ge0.
\]
Thus the exchange graph supplies a dissipative Laplacian structurally.
Its conductances are new local coupling data until extracted from actual
MM/HC transitions; positivity and reciprocal pairing must be checked.

For a route generator $\mathcal A$ in this cell representation, let
$\mathcal A^{\dagger_V}=V^{-1}\mathcal A^{\mathsf T}V$, with $V$
the block cell-volume matrix. The split into weighted skew and symmetric
parts gives conservative routing and a separately retained energy-changing
part. The same split used for MM above applies; directed routing is not
automatically a positive diffusion process.

On a periodic cubic mesh of spacing $h$, for prescribed constant velocity
$b$ and diffusivity $D$, the consistent choices are
\[
(\mathcal J_hr)_i
=\sum_{j=1}^3 b_j\frac{r_{i+e_j}-r_{i-e_j}}{2h},
\qquad
(\mathcal L_hr)_i
=\sum_{j=1}^3 D\frac{2r_i-r_{i+e_j}-r_{i-e_j}}{h^2}.
\]
The first is skew-adjoint in the uniform periodic inner product and the
second is positive semidefinite. For a smooth sampled field Taylor
expansion gives
\[
\mathcal J_hr=b\cdot\nabla r+O(h^2),\qquad
-\mathcal L_hr=\sum_j\partial_j(D\partial_jr)+O(h^2).
\]
This derives the continuum transport and diffusion terms as a formal
consistency limit of a declared spatial route model. It is not a convergence
proof for arbitrary native atlases, meshes or variable coefficients.
Boundary routes, volume preservation and the guard block structure must
also be supplied.

\subsection{The derived candidate and the remaining fluid step}
Combining spatial routing, edge exchange and the local witness residual
equation gives the finite model
\[
\dot r=-\mathcal J_hr-\mathcal L_hr-Kr-B_{\rm eff}r
+\mathcal N_{\rm MM}(r)+s_R,
\]
and formally the continuum candidate
\[
\partial_t r+b\cdot\nabla r
=\sum_j\partial_j(D\partial_jr)-Kr-B_{\rm eff}r
+\mathcal N_{\rm MM}(r)+s_R.
\]
The signs of the skew and dissipative terms follow from the stated
construction. Exponential decay needs a positive net gap after symmetric
MM effects and nonlinear growth have been included.

To obtain Navier--Stokes one still needs a native spatial velocity map,
self-transport $b=u$, an incompressibility constraint and its pressure
projection, and a Newtonian momentum-flux law. None follows from a
three-channel fit score alone. The derivation currently supports residual
transport and relaxation; Navier--Stokes is a further constitutive specialization.

\subsection{The scalar recurrence as a consistency check}
In the special one-mode case with no skew contribution, source or nonlinear
growth and constant net damping $\beta>0$,
\[
\dot r=-\beta r,\qquad
\Phi(t)=e^{-2\beta t}\Phi(0).
\]
Matching a measured or independently derived sample interval $\Delta t$
to the C5 ratio requires
\[
e^{-2\beta\Delta t}=\frac34,
\qquad \beta=\frac{\log(4/3)}{2\Delta t}.
\]
The existing recurrence determines this product, not a unique mobility or
viscosity. In contrast, minimizing the r31 penalty
$(\Phi_{n+1}-3\Phi_n/4)^2$ would enforce an already prescribed ratio and
cannot independently derive it. We therefore use the recurrence only as
an endpoint consistency test for the constructed residual dynamics.

\subsection{A verified finite witness and route instance}
Take the structural test witness
\[
X=(r_1,r_2,r_3,c_4,c_6)\in\mathbb R^5,
\quad R(X)=(r_1,r_2,r_3),
\quad C(X)=(c_4-c_4^0,c_6-c_6^0).
\]
These are independently specified model constraints, not extracted native
guard evaluations. They give
\[
\Pi=\operatorname{diag}(1,1,1,0,0),\qquad Q=\gamma I_5,
\qquad B_{\rm HC}=\gamma I_3\quad(\gamma>0).
\]
Choose a tangent MM model with residual generator
\[
A_{\rm MM}=
\begin{pmatrix}\eta&-\omega&0\\\omega&\eta&0\\0&0&\eta\end{pmatrix}.
\]
Its skew part rotates the first two channels and its symmetric part is
$\eta I_3$. The full local model keeps $c_4,c_6$ fixed and has
$B_{\rm eff}=(\gamma-\eta)I_3$. Thus $\gamma>\eta$ gives decay,
$\gamma=\eta$ permits undamped rotation, and $\gamma<\eta$ gives growth
despite a positive compensator mobility.

For a three-cell cyclic exchange graph with scalar conductance $d>0$,
\[
L_3=d\begin{pmatrix}2&-1&-1\\-1&2&-1\\-1&-1&2\end{pmatrix}
\]
has eigenvalues $0,3d,3d$. Its energy form is
$d[(r_1-r_2)^2+(r_2-r_3)^2+(r_3-r_1)^2]$ for scalar cell data.
For vector-valued channels the same identity applies componentwise.
The constant-cell mode is not damped by exchange alone; the positive local
gap $\gamma-\eta$ removes that obstruction in this model.
The guard-projector product, projected damping matrix, MM decomposition,
graph spectrum and graph energy identity were checked symbolically.

\paragraph{Status of this structural derivation.}
The native source supplies the channel split, residual-score form and
guarded witness setting. Proven here are the chain-rule descent criterion,
positive projected-gradient damping, the MM skew/symmetric decomposition,
and the graph-exchange energy identity. Additional model inputs are the
smooth witness vector fields, guard Jacobians, positive mobility, spatial
route realization and conductances. Their actual values and agreement with
the MM/HC objects have not been extracted from the source. The continuum
pass is formal consistency; physical fluid interpretation remains a further
step. This separates an explicit structural derivation from an asserted
native physical derivation.

\section{A curved-pipe route with unchanged fluid properties}
\label{sec:curved-route}
The user's pipe picture gives a different geometric specialization: the
material is fixed while its admissible transport route changes. Keep
$\rho_0$ and $\mu=\rho_0\nu$ constant. Curvature changes the flow domain,
pressure and velocity gradients, not these constitutive constants by definition.
Incompressibility means that pressure redistribution at a bend is not
automatically a change in density.

\subsection{A tube around a planar curved route}
Let $c(s)$ be an arc-length parametrized smooth planar centreline, with
an oriented orthonormal frame $(t,n,b)$ obeying
\[
t'=\kappa n,\qquad n'=-\kappa t,\qquad b'=0.
\]
The signed curvature $\kappa(s)$ may vary. A tube of constant radius $a$
has local coordinates
\[
X(s,\xi,\zeta)=c(s)+\xi n(s)+\zeta b,
\qquad \xi^2+\zeta^2<a^2,
\qquad h=1-\kappa(s)\xi>0.
\]
The local nonsingularity condition $a|\kappa|<1$ is supplemented by no
global self-intersection of the tube. The Euclidean metric and volume are
\[
d\ell^2=h^2ds^2+d\xi^2+d\zeta^2,
\qquad dV=h\,ds\,d\xi\,d\zeta.
\]
For physical orthonormal velocity components $(u_s,u_\xi,u_\zeta)$,
incompressibility is
\[
\nabla\cdot u=\frac1h
\left[\partial_su_s+\partial_\xi(hu_\xi)
+\partial_\zeta(hu_\zeta)\right]=0.
\]
Even with unchanged material properties, the same component profile can
therefore have different derivatives in different bends. The scalar
Laplacian is
\[
\Delta f=\frac1h\partial_s\left(\frac1h\partial_sf\right)
+\frac1h\partial_\xi(h\partial_\xi f)+\partial_\zeta^2f.
\]
The vector viscous operator must additionally account for derivatives of
the moving frame; applying this scalar expression independently to moving
vector components would omit connection terms.

\subsection{Pressure turns the flow; it need not compress the fluid}
For a centreline trajectory with velocity $U(s)t(s)$, differentiation along
the trajectory gives
\[
\frac{Du}{Dt}=UU't+\kappa U^2 n.
\]
In a leading inertial turning balance, neglecting transverse viscous and
body-force terms only for this estimate, the normal momentum equation gives
\[
\partial_nP\simeq-\rho_0\kappa U^2.
\]
Here $n$ points toward the bend centre when $\kappa>0$, so pressure is
higher toward the outside. Across a small-radius cross-section this suggests
a pressure difference of order $\rho_0U^2a|\kappa|$.
This is a turning estimate, not an exact full pipe solution. Secondary
motion and no-slip wall layers require the full velocity and pressure fields.

With the convention
\[
\mathrm{Re}=\frac{2aU}{\nu},\qquad
\mathrm{De}=\mathrm{Re}\sqrt{a|\kappa|},
\]
the Dean number combines flow speed, viscosity and curvature. Definitions
using other reference lengths differ by numerical factors. Different bends
can consequently change the flow profile and secondary circulation while
the fluid remains the same. No universal monotone loss formula or transition
threshold is inferred from these dimensionless numbers alone.

\subsection{Curvature as an energy-neutral route connection}
Let $F(s)$ have columns $(t,n,b)$ and write a transported physical residual
vector as $v(s)=F(s)r(s)$. Then
\[
F'=F\Omega_\kappa,\qquad
\Omega_\kappa=
\begin{pmatrix}0&-\kappa&0\\\kappa&0&0\\0&0&0\end{pmatrix},
\qquad v'=F(r'+\Omega_\kappa r).
\]
Thus a route-based residual equation can be written
\begin{equation}
\label{eq:curved-route-residual}
r'+\Omega_\kappa r=-\Lambda(s)r+q(s),
\end{equation}
where $\Lambda$ is a genuine loss or relaxation operator per unit length
and $q$ a source relative to the declared residual reference. The geometric
connection is skew-symmetric. Therefore
\[
\frac{d}{ds}\frac{|r|^2}{2}
=-r^{\mathsf T}\Lambda_{\rm sym}r+r\cdot q.
\]
Curvature can mix the components but does not itself decrease this norm.
If $\Lambda=0=q$, the components rotate while $|r|$ stays constant.
This also distinguishes mere frame rotation from deformation of an actual
velocity profile: the latter requires the fluid equations and boundary data.

For an isotropic reduced loss $\Lambda=d(s)I$ and $q=0$, the planar
connection matrices commute, so the exact route solution is
\[
r(s)=e^{-\int_0^s d(\sigma)d\sigma}
\exp\left[-\int_0^s\Omega_\kappa(\sigma)d\sigma\right]r(0),
\]
\[
\frac{|r(s)|^2}{|r(0)|^2}
=e^{-2\int_0^s d(\sigma)d\sigma}\quad(r(0)\ne0).
\]
A sampled route segment has C5 ratio $3/4$ only if its integrated loss is
$\int d\,ds=\log(4/3)/2$. Bending angle alone does not fix that integral.
For anisotropic $\Lambda$, bend-induced redistribution can change the loss
rate by rotating the residual into differently damped directions.
That is a mechanism for geometry-dependent behaviour with fixed loss
coefficients, not a derivation of those coefficients from MM/HC.

\subsection{Where the physical energy loss occurs}
On a stationary pipe domain, define
$S(u)=(\nabla u+\nabla u^{\mathsf T})/2$ and
$\sigma=-PI+2\mu S(u)$. For smooth incompressible flow the mechanical
energy balance is
\begin{align*}
\frac{d}{dt}\frac{\rho_0}{2}\int_\Omega|u|^2\,dV
+2\mu\int_\Omega S(u):S(u)\,dV
={}&\int_{\partial\Omega}u\cdot\sigma n\,dA\\
&-\frac{\rho_0}{2}\int_{\partial\Omega}|u|^2u\cdot n\,dA
+\rho_0\int_\Omega f\cdot u\,dV.
\end{align*}
Stationary no-slip walls have zero boundary velocity; the fluid gradients
near them still contribute to the positive volume dissipation.
For a steady unforced segment with equal inlet and outlet kinetic-energy
fluxes, negligible extra viscous work at the chosen cuts, and sectionwise
constant inlet and outlet pressures, this reduces to
\[
Q(P_{\rm in}-P_{\rm out})=2\mu\int_\Omega S:S\,dV.
\]
Hence geometry affects losses through the solved strain field, while $\mu$
can remain unchanged. The local loss in a reduced route equation must be
obtained from this balance or an independent closure law. It cannot simply
be identified with curvature. At fixed pressure drop the flux may change;
at fixed flux the required pressure drop may change.

\subsection{OP5 interpretation and a necessary scale test}
The proposed structural dictionary is: the transported witness or residual
is the fluid-like object, the declared route atlas is the pipe geometry,
and the moving comparison frame supplies the skew connection. Actual
relaxation and source terms remain separate. Native route lengths, curvature,
residual tensor type and compatible frame transformations must be extracted
before this dictionary becomes a native construction.

The reference against which the residual is measured matters. A straight-pipe
profile placed in a curved pipe need not solve its equations; the bend then
generates a mismatch relative to that reference. Relative to the correct
curved-pipe base solution the same steady flow may have zero residual.
One cannot identify ongoing fluid motion with an OP5 closure defect without
specifying that reference.

For infinitesimal speed $U=\varepsilon U_0$ with fixed standard radius,
curvature and viscosity, $\mathrm{Re}$ and $\mathrm{De}$ are infinitesimal
and the inertial turning pressure scale is of order $\varepsilon^2$.
Thus strong inertial bend effects are not automatic for the earlier halo
film. Finite geometric response might instead concern an infinitesimal
perturbation of a standard base flow, an extreme geometric scale, or a
different coupling regime. Those cases require their own analysis.

\paragraph{Curved-route result.}
The same fluid can behave differently along different bends. The model
derives a norm-preserving geometric connection and separates physical
viscous losses from it. It does not assume density compression, invent a
curvature-dependent viscosity, or infer a universal bend-loss coefficient.
The tube formulas and connection are local geometric constructions; the
full curved-pipe velocity field has not been numerically or analytically
solved here.
The tube metric was checked symbolically on a circular centreline, and the
skew-curvature energy cancellation and constant-coefficient route solution
were independently verified by symbolic differentiation.

\paragraph{Curved-pipe reference.}
P. H. M. Bovendeerd, A. A. van Steenhoven, F. N. van de Vosse and G. Vossers,
\emph{Steady entry flow in a curved pipe} (1987),
\href{https://research.tue.nl/en/publications/steady-entry-flow-in-a-curved-pipe}{Eindhoven University of Technology research record}.
This reference supports geometry-dependent axial and secondary flow;
the route-connection and energy identities above are calculated for this note.

\section{Native-data audit: one HC/MM chart and its source-return route}
\label{sec:native-chart-audit}
This pass returns to the external v4.7.0 source rather than selecting another
free fluid ansatz. It fixes one numerical HC/MM chart packet and the
source-return factorization proposed in the adjacent atlas block. Source
line references below apply to the file identified in the companion
Markdown audit; the model assumptions of earlier sections do not change
these source definitions.

\subsection{Three source objects must remain separate}
\begin{enumerate}
\item The r04 C5 witness-fit score is
$\Phi_5(W)=\|\rho_{5,\rm bd}(W)\|^2+
\|\rho_{5,\rm ext}(W)\|^2+\|\rho_{5,\rm mix}(W)\|^2$.
The bounded witness row records $\Phi_5(W_5^{\rm bd})=3/4$ with zero
declared fit guards (source around lines 52323--52365). It does not supply
the component residual values or a formula for that witness update.
\item The asymmetric photon-frame chart records finite radial sizes and
contraction budgets (source around lines 43348--43460). These data can be
evaluated numerically but are not identified with the three C5 fit residuals.
\item The older OP5 Step (B) target has the explicit defect
\[
\Delta_{\Theta R}(T)=\mathcal T(T)-
\left(\frac{3T}{2}-\log T-\frac{T^2}{24}
-\gamma_E-\beta_{\rm sym}^2T^2\right),
\quad T=\gamma_L^{-3/2},
\]
under the source label \texttt{def:delta-theta-r-residual}
(around line 26240). Its historical status is separate from the final
v4.7.0 release scope. Neither the finite radial table nor the C5 score
alone evaluates or proves the vanishing of this defect.
\end{enumerate}
Thus a fluid model for the C5 test is not automatically a model for every
OP5 formulation. A native bridge must state which of these targets it addresses.

\subsection{The concrete chart packet}
The source chart has eight rows
\[
\mathcal S=\{(b_i,q_i,a_i,e_i,w_i,
L_i^{\rm rout},L_i^{\rm unr})\}_{i=1}^8,
\qquad \lambda_i\le c_i:=q_i+a_i+e_i.
\]
Here $b_i$ is the source radial readout, $q_i$ the second-shift radius ratio,
$a_i$ a routed HC/MM absorption contribution, $e_i$ a tolerance allowance
and $w_i$ a positive readout weight. The contribution $a_i$ enters with a
positive sign in a conservative upper bound. It is not, without another
equation, a measured viscous-energy loss or the coefficient of $-\nu\Delta$.

The largest-budget row, row 6, is the specific input packet
\[
(b_6,w_6,q_6,a_6,e_6,L_6^{\rm unr})
=\left(\frac38,\frac{21}{16},\frac{21}{32},
\frac1{64},\frac1{64},0\right).
\]
Its contraction cap is $c_6=11/16$. If the source normalized ratio is
realized as $b_6'/b_6$, it implies
\[
0\le b_6'\le\frac{11}{16}\frac38=\frac{33}{128},
\qquad \frac{(b_6')^2}{b_6^2}\le\frac{121}{256}.
\]
The value $33/128$ is an envelope output, not an observed or uniquely
constructed native output. This chart row has a different contraction
budget from the C5 fit value $3/4$.

\subsection{A checked finite radial-score bound}
For an explicitly defined radial proxy, keeping the source weights fixed,
set $\Psi_b=\sum_i w_i b_i^2$. The table gives
\[
(c_1,\ldots,c_8)=
\left(\frac{17}{32},\frac{37}{64},\frac{39}{64},\frac{21}{32},
\frac{21}{32},\frac{11}{16},\frac{11}{16},\frac{41}{64}\right).
\]
Under the componentwise interpretation $0\le b_i'\le c_i b_i$,
\[
\Psi_b'=\sum_iw_i(b_i')^2
\le\left(\frac{11}{16}\right)^2\Psi_b.
\]
For the envelope choice $b_i'=c_i b_i$, exact rational arithmetic gives
\[
\Psi_b=\frac{107061}{32768},\qquad
\Psi_b'{}_{\rm env}=\frac{155153291}{134217728},
\]
\[
\frac{\Psi_b'{}_{\rm env}}{\Psi_b}
=\frac{155153291}{438521856}\simeq0.35381.
\]
This verifies the declared finite radial budget, not a physical energy
measurement or native C5 score equality.
The source chart margin $9/64$ gives the robust cap
\[
\Lambda_{\rm chart}=\frac{11}{16}+\frac9{64}=\frac{53}{64}<1,
\qquad \Lambda_{\rm chart}^2=\frac{2809}{4096}<\frac34.
\]
If the margin is a uniform operator or component-ratio allowance in the
same unchanged weighted norm, this also bounds the radial proxy score.
The source classifier alone does not prove that interpretation, repeated
admissibility, or equivalence to $\Phi_5$.

\subsection{The native route closest to the pipe picture}
The finite-atlas block (source around lines 43683--43740) proposes
\[
\mathcal R_{\rm OP}^{\rm HC\text{-}cone}
=T_\Gamma^{\vee}\circ J_{\rm HC}^{\rm cone}\circ T_\Gamma.
\]
This offers a definite ordered structural path: forward transport,
cone/collar readout or compensation, then return transport. The symbol
$\vee$ is retained as a native return tag; it is not assumed to mean an
inverse or an adjoint. The source explicitly labels this factorization as
a hypothesis in a conditional closure theorem, not a completed native
operator construction.

In the pipe analogy, $T_\Gamma$ and $T_\Gamma^\vee$ correspond to the
forward and return route, while $J_{\rm HC}^{\rm cone}$ corresponds to
the turning/readout interface. The source periapsis/no-through-collapse
guard is compatible with an outside turn rather than crossing the guarded
core. It does not assign a pipe radius, arclength, bend angle or curvature.
In particular the overlap parameter $\kappa$ in
$H=G[N+\kappa N(N-1)]$ is not identified with the inverse-length
pipe curvature $\kappa(s)$ of the previous section.

For normed native source, collar and return spaces, an actual route estimate
would follow from separately verified bounds
\[
\|T_\Gamma\|\le M_{\rm fwd},\qquad
\|J_{\rm HC}^{\rm cone}\|\le\lambda_{\rm cone},\qquad
\|T_\Gamma^\vee\|\le M_{\rm ret},
\]
\[
\|\mathcal R_{\rm OP}^{\rm HC\text{-}cone}\|
\le M_{\rm ret}\lambda_{\rm cone}M_{\rm fwd}.
\]
This is directly a bound for linear maps; an analogous bound for nonlinear
maps uses suitable zero-referenced bounds or Lipschitz constants and a
verified common reference state. Geometric frame rotations preserve norms,
so they alone cannot explain a cone contraction less than one. Some
additional projection, fit relaxation or dissipative coupling must be
specified. The finite table bounds the combined readout, but does not
uniquely separate these factors or recover their curvature.

\subsection{Concrete atlas and relay checks}
The source five-chart atlas records robust caps
\[
\left(\frac58,\frac{25}{32},\frac{25}{32},\frac{53}{64},
\frac{13}{16}\right)
\]
and declares zero unrouted leakage, unpaid absorption and transition debt.
Its maximum is $53/64$. This is the most concrete local atlas certificate
used in this pass. It lists chart types, not actual overlap transition
matrices or a metric adjacency graph. Its own conditional theorem still
requires coverage, transition stability, cone-factorization, no-through-core
behaviour, pass-count consistency and identification with the OP5 statement.

The later r29 seed anchor $(N,\kappa,G)=(5,3/4,1)$ yields
$H=20$ and relay margin $H-16=4$, as checked exactly. It supplies a
comfortable local relay packet, not a bending radius. Its recurrence
residual $\Phi_{n+1}-3\Phi_n/4=0$ can coexist with $\Phi_n>0$;
zero seed debt therefore does not erase the C5 fit remainder.

\subsection{QAM representation and the next native calculation}
A faithful initial QAM state for this audit stores the chart row or atlas
packet with its native meanings, the source-return route tag, the declared
guard register and an explicit target label: radial proxy, C5 fit or
Theta-balance. It must not silently substitute one target for another.
A QOp graph for actual native transport can be filled only after the maps
$T_\Gamma,J_{\rm HC}^{\rm cone},T_\Gamma^\vee$ are given on their
domains. A diagonal graph $b_i\mapsto c_i b_i$ is a valid envelope-model
operator, but is not reconstructed native transport.

The next source-level calculation is consequently concrete: specify the
state coordinates entering the forward map, calculate the cone/collar map
and return map on the row-6 or atlas packet, then evaluate both the native
radial readout and the actual three C5 residuals. Compatible overlap maps
can then be used to build a metric route only if the source supplies the
necessary distances or frame data. These maps, component residual functions
and distances are not supplied by the inspected chart and fit tables.

\paragraph{Result of the native-data pass.}
We have selected and evaluated an actual source chart, extracted its
proposed native source-return path, and obtained exact finite radial-score
bounds. The path supports the turning-route analogy at the level of order
and guards. It does not yet determine pipe geometry, physical viscosity,
or a native C5 witness update. This is a localized missing construction,
not evidence that the more general OP5 problem is already solved by the
fluid model.

\section{Executing the selected packet: native limits and an explicit completion}
\label{sec:packet-execution}
The source search was extended to the earlier OP transport architecture.
It specifies $T_\Gamma=A_4\circ A_3\circ A_2$ around line 21957, types
forward and reversed transports on a closure sector, and defines
frame-dependent readouts $\rho_\theta$. Its OP closure theorem assumes
HC confinement, native class preservation and compatible invariant readout.
Those assumptions are not numerical evaluation rules for row 6.
The later $T_\Gamma^\vee$ and $J_{\rm HC}^{\rm cone}$ in the atlas
factorization are not thereby identified with concrete matrices or with
the earlier reversed transport. Other occurrences of $A_2,A_3,A_4$ in
coding menus and guard lists have different roles and cannot fill that gap.

\subsection{Native execution trace}
Let $x_0$ be the actual carrier state represented by the row-6 packet.
The intended native trace is
\[
x_2=A_2x_0,\quad x_3=A_3x_2,\quad x_4=A_4x_3,
\quad y=J_{\rm HC}^{\rm cone}x_4,
\quad x_{\rm out}=T_\Gamma^\vee y.
\]
The scalar $b_6=3/8$ does not uniquely specify $x_0$. In particular,
an unshifted radius is not a full carrier coordinate vector, a tensor field
or a three-component fit residual. The inspected native source gives
the route types and scalar budgets, but no row-6 carrier-coordinate
realization, evaluated intermediate states or component formulas
$\rho_{5,\rm bd}(x_{\rm out})$, $\rho_{5,\rm ext}(x_{\rm out})$ and
$\rho_{5,\rm mix}(x_{\rm out})$.
Accordingly these true native outputs cannot be calculated from the
specified packet alone. The following completion is a model execution,
with every introduced rule displayed rather than attributed to the source.

\subsection{A minimal envelope completion of the three route maps}
Represent the row's radial input by a model vector $z_0=(3/8,0,0)$.
Let $O\in SO(3)$ be a freely chosen frame rotation. For row 6 set
\[
T_{\rm env}(z)=\frac{21}{32}Oz,
\qquad J_{\rm env}(y)=\frac{22}{21}y,
\qquad T_{\rm ret,env}(y)=O^{\mathsf T}y.
\]
These factors reproduce the chart cap:
\[
T_{\rm ret,env}\circ J_{\rm env}\circ T_{\rm env}(z)
=\frac{11}{16}z.
\]
For $O=I$, every intermediate value is explicit:
\[
z_0=(3/8,0,0),\quad
z_{\rm fwd}=(63/256,0,0),\quad
z_{\rm cone}=(33/128,0,0),\quad
z_{\rm out}=(33/128,0,0).
\]
The enlargement $22/21$ is the envelope allowance
$(q_6+a_6+e_6)/q_6$, not a derived physical HC response.
The $a_6$ and $e_6$ entries were budget terms, so interpreting their
sum as an actual deterministic gain is an added envelope rule.
Their zero-leakage source registers remain source declarations, not
computed native guard evaluations in this completion.

More generally each source row admits the explicit envelope maps
\[
T_{{\rm env},i}=q_iO_i,\qquad
J_{{\rm env},i}=\frac{c_i}{q_i}I,\qquad
T_{{\rm ret,env},i}=O_i^{\mathsf T},
\]
since all the table's $q_i$ are positive. Their product is $c_iI$.
This executes all eight caps without selecting a bending angle: any
$O_i$ produces the same radial output. The completion therefore proves
that these radial data do not identify frame rotations or pipe curvature.
It is one compatible envelope factorization, not a unique reconstruction
of the native MM/HC maps, nor a decomposition into the native $A_2,A_3,A_4$.

\subsection{What the native C5 scalar does and does not determine}
Even the separate known value $\Phi_5(W_5^{\rm bd})=3/4$ cannot recover
the three channel residuals. For scalar real channel examples,
\[
r^{(a)}=(\sqrt3/2,0,0),\qquad r^{(b)}=(1/2,1/2,1/2)
\]
both have squared norm $3/4$, but entirely different channel distributions.
These are algebraic completions of the score, not proofs that both are
native admissible witnesses. The missing residual functions and guard
evaluations would have to decide their native status.

There is also a useful incompatibility test. If one identifies the row-6
radial vector directly with the retained residual vector, uses an unchanged
homogeneous quadratic score, and identifies this envelope step with the
C5 lowering step, its score ratio would be
\[
\left(\frac{11}{16}\right)^2=\frac{121}{256}
\ne\frac34.
\]
Even the squared robust radial cap $2809/4096$ is less than $3/4$.
Thus that direct identification cannot simultaneously realize the stated
C5 equality under the same cap and norm. This is a failure of the proposed
identification, not a contradiction in the source: the chart measures a
second-shift radial readout, whereas C5 measures a witness-fit score.
A different map, norm, reference or step correspondence would need to be
specified and justified.

\subsection{QAM execution and exact remaining input contract}
The envelope input, intermediate states and output can be stored as QState
payloads with the native packet as provenance. The three displayed linear
maps give functional QOp graphs and their output radii give QReadout graphs.
Their witness certificate is tagged \emph{envelope-model execution}; it
must not carry native C5 residual values or native QAdmTransport status.
No physical viscosity or curvature is inferred from the selected rotations.

To change that tag to a native execution, the required data are now explicit:
the carrier-coordinate representative of the chosen packet; evaluation
rules for $A_2,A_3,A_4,J_{\rm HC}^{\rm cone},T_\Gamma^\vee$ with matching
domains; the three actual C5 residual functions and their norm/reference;
and separately evaluated guard and class-preservation conditions.
The existing radius and score tables do not supply these rules. A new
construction may supply them, but it must be declared as such rather than
reported as already extracted from the Companion.

\paragraph{Execution verdict.}
The full three-stage envelope route has been calculated, and its finite
radial-score bound is verified. The native route has been traced to its
typed factors, but its actual carrier and C5 component outputs remain
underdetermined by the inspected source. The direct radial-to-C5
identification has been tested and rejected. This prevents further fluid
interpretation from silently identifying two different native readouts.

\section{Source-backed explanations for the unequal contraction factors}
\label{sec:source-backed-bridge-explanations}
A focused search through the original source found more explicit OP5-facing
readout formulas than the earlier eight-row table alone supplied. These
explain why direct radial-to-fit identification is unjustified, although
they still do not define the three r04 C5 fit components as functions of
the selected row-6 state.

\subsection{Native radial and cone readouts are different functions}
The halo-channel block around source lines 44160--44280 defines a second
shift $U_\pi$, radial projection $P_{\rm rad}$, and their tilt commutator
\[
\Theta=P_{\rm rad}U_\pi-U_\pi P_{\rm rad}.
\]
In schematic distance notation its radial readout is
\[
\rho_r(x)=d_G(P_{\rm rad}U_\pi x,K).
\]
The HC-residual readout is instead
\[
\rho_h(x)=\frac{|h(x)^2-r(x)^2|}{h(x)^2+r(x)^2+\delta_0}
+\mu_\Theta\|\Theta(x)\|_G,
\quad \delta_0>0,\quad\mu_\Theta\ge0.
\]
Thus kernel distance, normalized cone mismatch and projection tilt are
different quantities already in the native text. The positive comparison
norm and actual projection/shift operators still need their declared
realizations to evaluate a specific carrier.

There is a concrete geometric reason for differing zero sets. With a
nonzero radial collar, $h=r>0$ and zero tilt, the cone mismatch vanishes
while the kernel-distance readout can remain positive. This permits an
exactly cone-normal state outside the excluded kernel, not a through-core
collapse. It does not by itself prove vanishing of the separate C5 residuals.

\subsection{Normalization changes the contraction law}
For the cone-defect part alone, suppose a declared update scales both
$h$ and $r$ by $\lambda\in(0,1)$ and holds $\delta_0$ fixed. Define
$S=h^2+r^2$ and $Q=h^2-r^2$. For $Q\ne0$ the exact ratio is
\[
\frac{|\lambda^2Q|/(\lambda^2S+\delta_0)}
{|Q|/(S+\delta_0)}
=\frac{\lambda^2(S+\delta_0)}{\lambda^2S+\delta_0}.
\]
For $S,\delta_0>0$ it lies strictly between $\lambda^2$ and one.
It is therefore not the squared-radius ratio, despite using the same
underlying scaled coordinates. If the normalization floor scales as well,
or if the tilt term contributes, the law changes again.

As an algebraic illustration, $\lambda=11/16$ and
$\delta_0/S=121/284$ give a cone-defect ratio of $3/4$ exactly.
This is not a source-derived parameter value, and that ratio concerns the
cone-defect readout itself, not its squared norm or the C5 fit score.
Its purpose is to demonstrate an explicit reconciliation mechanism, not
to tune a claimed native proof. For a squared-defect score the ratio would
be the square of the displayed expression; the actual C5 dictionary is
still required.

\subsection{The source has an explicit two-coordinate collar map}
The certified-collar block around source lines 44771--44900 introduces
coordinates $(a,b,z)$ with
\[
\rho_r=r_{{\rm HC},j}^{(1)}(z)+\epsilon_r a,
\qquad \rho_h=\epsilon_h b+\Delta_j(a,b,z),
\]
\[
|\partial_b\Delta_j|\le\theta_j\epsilon_h,
\qquad\theta_j<1.
\]
Here $b$ is a cone-height coordinate, not the earlier chart radius $b_i$.
The source explicitly avoids defining it as $\rho_h/\epsilon_h$, so the
derivative certificate is a nontrivial condition rather than a tautology.
The source symbolic examples are
\[
\frac{\Delta_1}{\epsilon_h}=\frac14ab,\qquad
\frac{\Delta_2}{\epsilon_h}=\frac38b(1-a),\qquad
\frac{\Delta_3}{\epsilon_h}=\frac12ab(1-a).
\]
Direct differentiation gives the normalized determinants
\[
\frac{\det D_{(a,b)}(\rho_r,\rho_h)}{\epsilon_r\epsilon_h}
=1+\frac a4,\quad
\frac{11}{8}-\frac{3a}{8},\quad
1+\frac{a(1-a)}2,
\]
respectively. All are at least one on $0\le a\le1$.
They verify two independent readout directions on these symbolic pieces.
The general source bound $1-\theta_j$ is conservative; the third explicit
piece even has $|\partial_b\Delta_3|\le\epsilon_h/8$.
This supplies a genuine local coordinate-level explanation for why a
single radial contraction does not determine cone-debt evolution.
It does not assign row 6 a specific $(a,b,z)$ or provide the C5 residual map.

\subsection{Collar guards restrict what may approach zero}
The native executed halo requires
\[
0<r_{\rm HC}(x)\le\rho_r(x)\le\epsilon_r,
\qquad\rho_h(x)\le\epsilon_h.
\]
Hence if $r_{\rm HC}$ has a positive standard lower bound, admissible
kernel distance cannot converge to zero. Closure may concern cone/fit
defects instead of the absolute radius, or distance relative to the collar.
The native collar has to be included in any use of the earlier zero-limit
model. To place an infinitesimal radial halo around the kernel, the collar
scale itself must fit inside that halo; a fixed positive standard lower
bound would make such a halo empty.

In the displayed affine coordinates, the active halo inequalities require
\[
a\le1-\frac{r_{{\rm HC},j}^{(1)}(z)}{\epsilon_r},
\qquad b+\frac{\Delta_j}{\epsilon_h}\le1.
\]
The entire unit square is therefore not automatically an admissible halo
piece. The source symbolic formulas have to be used on the appropriate
restricted subregion. If the two epsilon scales are promoted to positive
hyperreal infinitesimals, an unnormalized Jacobian determinant can itself
be infinitesimal but nonzero; taking its standard part then need not
preserve rank. The normalized determinant above is the useful coordinate
certificate. A formal nonstandard model is still needed before transfer
can be invoked for the source's halo terminology.

\subsection{A source-supported chart-calibration mechanism}
The lens-refit block around source lines 47979--48076 defines a bounded
bifocal fit with residual
\[
\Delta_{\rm lens}(x;\widehat a,\widehat c)
=|d(x,\widehat F_-)+d(x,\widehat F_+)-2\widehat a|.
\]
It reports a finite 315-point diagnostic with mean defect reduced from
about $0.0764$ to $0.0033$ by a scale-local refit, and explicitly distinguishes
chart debt from a carrier obstruction. These are source-reported results,
not a rerun of the underlying diagnostic here. They support a mechanism
where fitted residuals change through readout calibration without requiring
the same change in geometric distance or physical energy. Means alone do
not determine the sum of squared residuals or a $3/4$ rate.

\paragraph{Search conclusion.}
The strongest source-backed explanations are independent readout channels,
nonlinear cone normalization, projection/shift tilt, a nonzero collar
reference, and scale-local fit calibration. The radial table concerns a
second-shift readout while the C5 table concerns a later witness-fit step.
These findings support keeping their factors separate and narrow the search
for a bridge to a specific nonlinear residual/readout map. No inspected
definition equates the two scores or derives the C5 factor from row 6.
The scalar-normalization identity and symbolic collar Jacobians were checked
independently with exact symbolic algebra.

\section{A constructed collar-to-fit bridge}
\label{sec:constructed-collar-fit-bridge}
We now supply an explicit candidate for the previously missing local map.
It uses the source's certified collar coordinates and symbolic tilt pieces,
not the absolute chart radius as a surrogate for all fit residuals. The
identification with the native r04 functions remains a hypothesis to test;
all newly assigned fit components are marked with tildes.

\subsection{Recovering collar coordinates from the two readouts}
Fix a source collar piece and its passive coordinate $z$. Set
$r_c=r_{{\rm HC},j}^{(1)}(z)$ and let $\epsilon_r,\epsilon_h>0$.
For the three explicit symbolic pieces write
\[
\Delta_j(a,b)=\epsilon_h b f_j(a),\qquad
f_1(a)=\frac a4,\quad f_2(a)=\frac38(1-a),\quad
f_3(a)=\frac12a(1-a).
\]
Their readout map is
\[
\rho_r=r_c+\epsilon_r a,\qquad
\rho_h=\epsilon_h b(1+f_j(a)).
\]
Since $1+f_j(a)>0$ on $0\le a\le1$, it has the explicit inverse
\[
a=\frac{\rho_r-r_c}{\epsilon_r},\qquad
b=\frac{\rho_h}{\epsilon_h[1+f_j(a)]}.
\]
Thus, given the piece, the collar scale and $z$, the two readouts determine
the collar coordinates. Row 6 alone still lacks this piece/scale/reference
assignment; the inverse is not claimed to fill that missing datum automatically.

\subsection{Three separately recorded fit components}
Define the candidate map on the fixed piece by
\[
\widetilde R_j(a,b,z)=
\begin{pmatrix}
\widetilde\rho_{5,\rm bd}\\
\widetilde\rho_{5,\rm ext}\\
\widetilde\rho_{5,\rm mix}
\end{pmatrix}
:=
\begin{pmatrix}a\\b\\b f_j(a)\end{pmatrix},
\qquad
\widetilde\Phi_j=a^2+b^2+[b f_j(a)]^2.
\]
The proposed roles are bounded collar-excess, the independently defined
cone-height coordinate, and the mixed tilt correction $\Delta_j/\epsilon_h$.
The last is the source's projected local correction, not an identification
with every component of the full commutator $\Theta$.
These role assignments need a native comparison with the r04 residual
functions; the source does not itself declare those equalities.

This is a three-component fit map of a two-coordinate fragment: its third
component is determined by the first two, so it does not pretend to create
three independent physical degrees of freedom. Its first two components
also show that it is injective on the fixed $(a,b)$ piece. The scalar
HC readout sums $b+bf_j(a)$, whereas the fit score retains the two terms
separately before squaring.

\begin{proposition}[Local zero set and collar compatibility]
On a fixed symbolic piece with a fixed passive reference $z$,
\[
\widetilde\Phi_j=0
\Longleftrightarrow a=b=0.
\]
The target has $\rho_r=r_c>0$ and $\rho_h=0$. If
$0<r_c\le\epsilon_r$, it satisfies the radial and HC halo inequalities.
It is a collar-normal target outside the kernel, not absolute radial zero.
\end{proposition}
\begin{proof}
The score is a sum of nonnegative squares containing $a^2$ and $b^2$.
At their common zero the third term also vanishes. Substitution in the
source readouts gives the target values and its halo membership.
\end{proof}
Any obstruction retained in the passive data $z$, the extra r02 slots,
or another native chart is outside this local zero-set certificate.

\subsection{A genuine coordinate update and its nonlinear score law}
For $0<t\le1$, define the witness-level candidate
\[
F_t(a,b,z)=(ta,tb,z).
\]
Its induced residual update is exactly
\[
\widetilde R_j(F_t(a,b,z))
=(ta,tb,tb f_j(ta)).
\]
It is not generally $t\widetilde R_j(a,b,z)$. For example, on piece 1
the mixed component contracts by $t^2$, not $t$.
This is a nonlinear witness-to-residual lift, rather than a scalar-score
update with unspecified component action.

\begin{proposition}[Readout-guard invariance of the symbolic pieces]
Suppose $0<r_c<\epsilon_r$ and initially
\[
0\le a\le1-r_c/\epsilon_r,\qquad
0\le b,\qquad b(1+f_j(a))\le1.
\]
Then $F_t$ preserves these halo inequalities for all $0<t\le1$ on all
three symbolic pieces.
\end{proposition}
\begin{proof}
The radial excess decreases, so the lower collar and upper radial bounds
remain valid. Put $H_j(t)=tb[1+f_j(ta)]$. Direct differentiation gives
\[
H_1'(t)=b(1+at/2),\quad
H_2'(t)=b(11/8-3at/4),\quad
H_3'(t)=b(1+at-3a^2t^2/2).
\]
For $0\le at\le1$ these are nonnegative, with respective lower bounds
$b$, $5b/8$ and $b/2$. Thus $H_j(t)\le H_j(1)\le1$.
\end{proof}
This proves the displayed radial/cone readout guards. Preservation of native
C4/C6, leakage and closure classes still requires their separate witness
definitions; fixed passive data are not by themselves that proof.

\subsection{Strict fit lowering without prescribing a ratio}
\begin{proposition}[Monotone local fit contraction]
For a nonzero admissible $(a,b)$ in the unit square, define
\[
P_j(t)=\widetilde\Phi_j(ta,tb)
=t^2(a^2+b^2)+[tb f_j(ta)]^2.
\]
Then $P_j$ is strictly increasing on $0<t\le1$ and $P_j(0)=0$.
Consequently every $F_t$ with $0<t<1$ gives strict score lowering.
\end{proposition}
\begin{proof}
For piece 1 the mixed square is $a^2b^2t^4/16$, so every nonzero term
has positive derivative. For piece 2, put $x=ta\in[0,1]$. Its derivative is
\[
P_2'(t)=2t\left[a^2+b^2+\frac{9}{64}b^2(1-x)(1-2x)\right].
\]
Since $(1-x)(1-2x)\ge-1/8$ on this interval,
$P_2'(t)\ge2t[a^2+(503/512)b^2]>0$.
For piece 3,
\[
P_3'(t)=2t(a^2+b^2)
+\frac{a^2b^2}{2}t^3(1-x)(2-3x).
\]
The product is at least $-1/12$, so
$P_3'(t)\ge2t(a^2+b^2)-a^2b^2t^3/24>0$ for nonzero $(a,b)$
in the unit square. Continuity gives the endpoint value at zero.
\end{proof}
The mixed component on piece 2 can increase as $a$ decreases near the
outer part of the chart; the proposition controls the total score rather
than silently claiming that each component must decrease.

\subsection{Comparison with the observed C5 ratio}
As a valid model probe take $a=b=1/2$ and $r_c/\epsilon_r=1/4$.
All three pieces satisfy the displayed halo guards. Their initial scores are
$129/256$, $521/1024$ and $129/256$ respectively. The uncalibrated half-step
$t=1/2$ gives the exact score ratios
\[
\frac{171}{688},\qquad\frac{2129}{8336},\qquad\frac{2057}{8256}.
\]
They are not equal to $3/4$. This is a direct test of the constructed map,
not a claim that the source's bounded witness uses this probe or step.

Monotonicity also proves that, for each nonzero input and each desired
score ratio $q\in(0,1)$, there is one and only one $t\in(0,1)$ with
\[
P_j(t)=qP_j(1).
\]
For $q=3/4$ and the probe above the calibrated steps are approximately
\[
t_1=0.8668609233,\qquad t_2=0.8638582271,\qquad t_3=0.8661446361.
\]
Normalization of all components by the fixed square root of the initial
probe score gives reference score one and output score $3/4$.
This demonstrates a possible nonlinear collar-to-fit realization of the
recorded ratio. It does not independently derive that ratio: these $t_j$
were solved from it. If the same calibration is applied repeatedly, the
step must generally be recomputed from the current coordinates.
Near the common zero its leading amplitude factor tends to $\sqrt3/2$,
consistent with a quadratic fit norm, not with absolute radius contraction.

The actual radial readout transforms as
\[
\frac{\rho_r(F_t x)}{\rho_r(x)}
=\frac{r_c+\epsilon_r ta}{r_c+\epsilon_r a},
\]
which tends to one as the collar target is approached. This is why a
nonzero collar permits decreasing fit defects without forcing the absolute
radius to zero. The row-6 cap $11/16$ cannot be imposed on that absolute
radius indefinitely with a fixed positive collar.

\subsection{QAM bridge contract and present result}
The candidate QState payload is $(j,a,b,z,r_c,\epsilon_r,\epsilon_h)$,
with its actual readout-guard certificate. The formulas above supply a
well-defined QReadout for the three-component fit vector and a QOp graph
for $F_t$ or the calibrated $F_{t(x)}$. The proved witness content is
local score lowering, the collar/cone halo inequalities and the explicit
zero set. Native QAdmTransport and equality with the source C5 functions
are not part of that certificate until separately verified.

We now have an explicit local map rather than only a request for one.
It is source-shaped, nonlinear, guard-compatible on the tested symbolic
pieces, and can realize the recorded score decrease by stated calibration.
The remaining native task is to test whether the proposed bounded, exterior
and mixed roles match the actual r04 residual functions and whether the
native witness update induces the constructed $F_t$. No viscosity, spatial
curvature or physical time is recovered from this fit map alone.
The inverse readout, Jacobians, score polynomials, exact probe ratios and
calibration roots were independently checked with symbolic algebra.

\section{A new local C5 normal form with a proved return operator}
\label{sec:new-local-c5-normal-form}
This section makes the preceding candidate into a defined, self-contained
local extension. Its component functions and update rule are supplied here,
not claimed to have been recovered from the older unspecified native maps.
The result is a theorem for a fixed collar piece, with a separate compatibility
contract for the inherited C5 and OP5 statements.

\paragraph{Chartwise scope.}
The score and flow below are local to the fixed piece $j$. They are not
automatically one common operator or score on an atlas. The explicit
overlap test in Section~\ref{sec:explicit-native-readout-bridge} detects
a failure of naive gluing and supplies a different, readout-compatible
construction. This does not invalidate the fixed-piece estimates below.

\subsection{Defined witness domain and local residual functions}
Fix $j\in\{1,2,3\}$, a passive reference $z$, and standard scales
$0<r_c<\epsilon_r$, $\epsilon_h>0$. Define the closed domain
\[
A_j=\left\{(a,b):0\le a\le1-r_c/\epsilon_r,\quad
0\le b,\quad b[1+f_j(a)]\le1\right\}.
\]
The functions $f_j$ are the three source symbolic collar corrections already
listed. Since $0\le f_j\le3/8$, this domain is compact and contains $(0,0)$.
Within this \emph{new local normal form}, define the three retained residuals
to be
\[
R^{\rm NF}_j(a,b)=(a,b,bf_j(a)),\qquad
\Psi_j(a,b)=\|R^{\rm NF}_j(a,b)\|^2.
\]
For a fixed nonzero reference $x_0=(a_0,b_0)$, the normalized score is
$\Phi^{\rm NF}_j(x)=\Psi_j(x)/\Psi_j(x_0)$.
These are now executable definitions. Their bounded/exterior/mixed role
labels refer to the restricted fit fragment, not to all native residuals.
Additional fit or leakage slots remain outside this definition.

The exact local normal-form closure set is defined by the geometric target
\[
Z_j^{\rm NF}=\{(a,b):\rho_r=r_c,\ \rho_h=0\}=\{(0,0)\}.
\]
The last equality follows from the invertible collar readout, not from a
chosen score threshold. Consequently $\Psi_j=0$ characterizes this local
target exactly.

\subsection{Autonomous collar relaxation and quantitative decay}
Introduce a dimensionless internal relaxation parameter $\tau$ and the
autonomous witness vector field
\[
\dot a=-a,\qquad\dot b=-b,\qquad\dot z=0.
\]
Its flow is $S_\tau(a,b,z)=(e^{-\tau}a,e^{-\tau}b,z)$.
This rule is a proposed local relaxation dynamics. It is not asserted to
be the inherited MM/HC operator, physical time or a viscosity law.

\begin{theorem}[Local normal-form witness and score theorem]
\label{thm:local-normal-form-flow}
For every initial state in $A_j$, the flow exists for all $\tau\ge0$, stays
in $A_j$, and retains a positive absolute collar radius.
For every nonzero initial state its score is strictly decreasing, positive
at finite $\tau$, and satisfies
\[
e^{-3\tau}\Psi_j(x_0)
\le\Psi_j(S_\tau x_0)
\le e^{-(376/219)\tau}\Psi_j(x_0).
\]
It converges to the unique local target $(0,0)$, with
$\rho_r\to r_c>0$ and $\rho_h\to0$.
\end{theorem}
\begin{proof}
The explicit flow and the previously proved readout-guard invariance give
existence and invariance. Put $Q=a^2+b^2$ and
$\mathcal D_j=a\partial_a\Psi_j+b\partial_b\Psi_j$. Directly
\[
Q\le\Psi_j\le\frac{73}{64}Q,\qquad
\frac{47}{24}Q\le\mathcal D_j\le3Q.
\]
For piece 1, $\mathcal D_1=2Q+a^2b^2/4$.
For piece 2 it is
$2[Q+(9/64)b^2(1-a)(1-2a)]$; the product is between $-1/8$ and one.
For piece 3 it is
$2Q+(a^2b^2/2)(1-a)(2-3a)$, where the product is between $-1/12$ and two.
Using $a,b\in[0,1]$ gives the stated conservative common bounds.
Since the flow has
\[
\frac{d\Psi_j}{d\tau}=-\mathcal D_j,
\qquad \frac{376}{219}\Psi_j\le\mathcal D_j\le3\Psi_j,
\]
integration proves the score estimates. The explicit flow gives positivity
for finite time at nonzero input and convergence of both coordinates to
zero. Substitution in the collar readouts gives the endpoint statements.
\end{proof}

\subsection{A unique score-level return map}
\begin{theorem}[Event return and discrete C5 realization]
\label{thm:local-c5-event-return}
For $q\in(0,1)$ and $x\in A_j\setminus\{0\}$ there is a unique finite
$\tau_q(x)>0$ such that
\[
\Psi_j(S_{\tau_q(x)}x)=q\Psi_j(x).
\]
Define $F_q(x)=S_{\tau_q(x)}x$ and $F_q(0)=0$. Then $F_q:A_j\to A_j$
is continuous, preserves the stated collar/cone guards, and obeys
\[
\Psi_j(F_qx)=q\Psi_j(x),\qquad
F_p\circ F_q=F_{pq}.
\]
In particular, for $q=3/4$ and a nonzero reference $x_0$,
\[
x_n=F_{3/4}^n x_0\in A_j,\qquad
\Phi^{\rm NF}_j(x_n)=(3/4)^n>0\quad(n<\infty),
\qquad x_n\longrightarrow(0,0).
\]
The threshold $1/256$ is first reached at step $20$ for this defined
normalized score. Exact local closure occurs at the limit, not at a finite
step for a nonzero initial witness.
\end{theorem}
\begin{proof}
The flow theorem gives continuous strict decay from $\Psi_j(x)$ to zero,
so the intermediate value theorem gives existence and uniqueness of the
event. Its time is bounded by
\[
\frac{\log(1/q)}3\le\tau_q(x)
\le\frac{219}{376}\log(1/q).
\]
The derivative of the event equation with respect to time is strictly
negative. The implicit function theorem therefore gives continuity of the
event time on nonzero inputs. At zero, $\|F_qx\|\le\|x\|$ proves
continuity of the map itself. The flow and domain invariance give the guards.
The flow semigroup and uniqueness of the score-level event prove the
composition rule, including at zero. Induction gives the score recurrence.
The lower bound on every event duration makes cumulative relaxation time
tend to infinity, so the coordinates tend to zero. Strict monotonicity of
the geometric score and the exact step-19/20 comparisons prove the first
threshold statement.
\end{proof}

The factor $q=3/4$ in this theorem is an explicitly selected \emph{return
level}, matching the inherited discrete specification. It is not a ratio
predicted independently by the autonomous flow. In particular, successive
event times need not be equal; a fixed physical-time recurrence would be a
stronger claim requiring a native clock and operator calculation.

\subsection{Guard scope and QAM encoding}
The theorem proves actual inequalities of the defined collar readouts.
The extended local state may carry additional C4 and C6 baseline coordinates
$g_4,g_6$ with the prescribed update
\[
\dot g_4=0,\qquad\dot g_6=0.
\]
Then $g_4=g_4^0$ and $g_6=g_6^0$ are proved invariant in this product model.
This does not prove invariance of an arbitrary native leakage functional
depending on $a,b,z$. For such a function $g$ the necessary local check is
\[
Dg(a,b,z)[-a,-b,0]=0
\]
on its equality-guard set, or an appropriate inward-pointing condition for
inequality guards. The normal form does not absorb a nonzero C4 residual
or unlock C6; it must be withheld from native admissibility if that check fails.

On the fixed piece, QAM state codes $s_x=\langle0,(j,x,z)\rangle$ and the
graph of $F_q$ provide a functional QOp. The graph of $R_j^{\rm NF}$ gives
a three-component QReadout. The local QWitness payload can cite
Theorems~\ref{thm:local-normal-form-flow} and
\ref{thm:local-c5-event-return}, the scale and domain data, and the declared
guard scope. A newly defined local closure label $(j,z)$ is preserved, but
it is not thereby the inherited native carrier class $K_{\rm ncl}$.
Native QAdmTransport requires the actual class-preservation bridge.

For an infinite hypernatural index $H$, transfer of this explicitly defined
discrete orbit and its convergence gives
\[
\Phi^{\rm NF}_j({}^*x_H)=(3/4)^H>0,
\qquad {}^*x_H\approx(0,0),\qquad
\operatorname{st}({}^*x_H)=(0,0)\in A_j.
\]
Its absolute radial readout is near $r_c$, not near zero when $r_c$ is
positive standard. The standard witness has exact local normal-form closure.

\subsection{Compatibility theorem for the inherited witness system}
\begin{theorem}[Conditional transfer to native C5]
\label{thm:native-normal-form-compatibility}
Suppose a native witness subset $W_j$ has an invertible collar chart
$\chi:W_j\to A_j$, a specified native update $F_{\rm nat}$, and actual
native residual functions. Suppose it is independently proved that
\[
R_{\rm nat}(X)=c_*R_j^{\rm NF}(\chi(X)),\qquad
\chi(F_{\rm nat}X)=F_{3/4}(\chi(X)),
\]
with fixed $c_*>0$ and matching native fit norms. Suppose the native guards
and closure class are preserved and the native local closure target is
exactly $\chi^{-1}(0)$. Then the local normal-form score, threshold and
finite-positive statements transfer to that native subset after fixed
reference normalization. If $\chi^{-1}$ is continuous at zero, its iterates
converge to the admissible native local closure witness $\chi^{-1}(0)$.
\end{theorem}
\begin{proof}
The update intertwining gives $\chi(X_n)=F_{3/4}^n\chi(X_0)$ by induction.
The component and norm identification gives the normalized score equality;
the fixed factor cancels in reference normalization. Apply the event-return
theorem. Continuity of the inverse chart gives the native endpoint; the
assumed independent target characterization supplies its local meaning.
\end{proof}
The first equality may be replaced by a fixed componentwise isometry in
the native residual spaces. Equality of total scores alone is insufficient
to identify the bounded, exterior and mixed roles. Equivalence to a broader
OP5 statement additionally requires its own target identification, atlas
overlap/coverage and the retained obstruction checks. These are not
conclusions of a one-piece normal-form theorem.

\paragraph{What has now been established.}
The local component definitions, invertible readout chart, witness dynamics,
global-in-relaxation-time collar invariance, quantitative decay, continuous
score-level return operator, exact discrete $3/4$ realization and admissible
local endpoint have been proved within the new fixed-piece construction.
The remaining task is no longer to invent an unspecified local operator:
it is to verify or reject the displayed compatibility equalities against
actual native definitions. Selecting the return level and carrying passive
guard coordinates have been kept explicit so they cannot masquerade as
native physical derivations.

\paragraph{Independent iteration check.}
For $a_0=b_0=1/2$ and $r_c/\epsilon_r=1/4$, twenty consecutive
score-level returns were evaluated on all three pieces with 60-digit
arithmetic and bracketed root finding. The normalized step-20 score was
$0.00317121193893$, below $1/256$, and all tested readout guards held.
The maximum score discrepancy from $(3/4)^n$ was below $8\cdot10^{-49}$.
The event durations varied across steps, as the theorem permits.
This numerical check supports the implementation of the constructed
return rule; the proofs above, rather than the numerical tolerance,
establish its exact properties.

\section{Re-audit of the origin and necessity of the factor three quarters}
\label{sec:three-quarter-origin-audit}
The number $3/4$ is arithmetically consistent throughout the proposed C5
chain. This section distinguishes its recorded origin from independent
operator mechanisms that could explain it.

\subsection{Where the native C5 chain first introduces it}
In the source C5 chain, the first explicit $3/4$ value is the r04 finite
witness ledger, around line 52359:
\[
\Phi_5(W_5^{(0)})=1,\qquad
\Phi_5(W_5^{\rm bd})=3/4,\qquad \Delta_5=-1/4.
\]
The r04 lowering lemma uses these table values and declared zero guards.
It does not calculate $3/4$ from explicit component residuals or a
specified witness operator. The later r05--r07 ledgers record matching
finite ratios; r08 defines $\mathcal U_5^{\rm bd}(\Phi)=3\Phi/4$;
r09 proves the geometric formula conditional on applying that update.
The induction is correct for the defined recurrence, but it is not a
separate derivation of the coefficient.

The local r27--r31 penalty for
$\Phi_{n+1}-3\Phi_n/4$ also prescribes the same coefficient rather than
explaining it. The later relay parameter $\kappa=3/4$ has a different
role in $H=G[N+\kappa N(N-1)]$; no inspected source relation identifies
it with the C5 contraction factor. No raw fit data or component evaluation
for the r04 bounded witness were found in the inspected source and
available local companion files. This is a documentation limit, not a
proof that no such construction can exist.

\subsection{An exact operator criterion for a quadratic score}
Suppose actual retained residuals lie in an invariant finite-dimensional
linear space with fixed positive definite metric $G$, and one native step
is a linear map $r'=Ur$. A uniform $3/4$ score law is equivalent to
\[
U^{\mathsf T}GU=\frac34G.
\]
Indeed the score ratio is
\[
q(r)=\frac{r^{\mathsf T}U^{\mathsf T}GU r}{r^{\mathsf T}Gr}.
\]
If $q(r)=3/4$ for all nonzero vectors in that space, polarization gives
the matrix identity; conversely it immediately gives the score law.
In whitened coordinates this requires every singular value of the
induced step to be $\sqrt3/2$; equivalently
\[
G^{1/2}UG^{-1/2}=\frac{\sqrt3}{2}O
\]
for an orthogonal matrix $O$. For one witness alone, $q(r)=3/4$ is only
one direction-dependent quadratic identity and does not prove this
uniform operator condition. Thus an actual MM/HC residual matrix and
declared fit metric would make the coefficient directly testable.

\subsection{A quarter-removal mechanism needs its own witness}
For an orthogonal least-squares correction $r'=r-d$ with
$\langle r-d,d\rangle=0$, Pythagoras gives
\[
\|r'\|^2=\|r\|^2-\|d\|^2.
\]
It yields $3/4$ exactly if the removable component has one quarter of
the initial residual energy. For example, projecting
$r=(\sqrt3/2,1/2)$ onto its first coordinate does so.
The source's three retained channel labels do not impose that energy
partition. Nor does the same orthogonal projection generate a geometric
$3/4$ recurrence: applying it again makes no further change because
$P^2=P$. Iterated quarter removal would require refreshed fit directions
and a new one-quarter identity at every step.
The value $\Delta_5=-1/4$ in the original table is the subtraction of
its listed scores, not an independent demonstration of this mechanism.

\subsection{The local normal form predicts a variable ratio at fixed time}
In the uncalibrated relaxation $S_\tau(a,b)=(e^{-\tau}a,e^{-\tau}b)$,
choose $t=e^{-\tau}=\sqrt3/2$. At the reference probe $a=b=1/2$,
the three exact ratios are
\[
\frac{515}{688},\qquad
\frac{6657}{8336}-\frac{27\sqrt3}{1042},\qquad
\frac{2105}{2752}-\frac{3\sqrt3}{344},
\]
approximately $0.7485465$, $0.7537041$ and $0.7497932$.
They are not identically $3/4$. The mixed collar correction changes the
score law even when the two base coordinates have a common amplitude factor.
Near the common zero, however, for fixed $t\in(0,1)$,
\[
\frac{\Psi_j(ta,tb)}{\Psi_j(a,b)}\longrightarrow t^2
\qquad\text{as }(a,b)\to(0,0),\ (a,b)\ne0.
\]
For pieces 1 and 3 the mixed term is higher order; on piece 2 its leading
linear component is $(3/8)b$, which still scales by $t$.
Thus $3/4$ is a possible leading quadratic-score factor if an independent
clock or native step gives $t=\sqrt3/2$. The collar geometry alone does
not select that amplitude or clock interval.

For a pure mode $\dot r=-\beta r$, an independent physical or structural
sample interval gives $q=e^{-2\beta\Delta t}$. The desired value requires
$\beta\Delta t=\log(4/3)/2$. Likewise, for a gradient step on an
eigenmode with $r'=(1-\eta\lambda)r$, one has
$q=(1-\eta\lambda)^2$. These are checkable derivation routes only when
the rate, interval, mobility and eigenvalue come from independent data.
Solving for them after requiring $q=3/4$ is calibration.

\subsection{Step twenty does not identify a unique coefficient}
For a normalized geometric score $q^n$ and detector threshold $1/256$,
the first crossing at step 20 means
\[
q^{20}\le1/256<q^{19},
\qquad (1/256)^{1/19}<q\le(1/256)^{1/20}.
\]
This is the whole interval approximately
$(0.74687948,\,0.75785828]$, not the single number $3/4$.
The arithmetic threshold check therefore verifies the consequences of
the selected factor but cannot explain its origin.

\paragraph{Coefficient verdict.}
The coefficient is currently a recorded native fit-table value, a defined
recurrence coefficient and, in the new local normal form, an explicitly
selected return level. Its finite-positive sequence and detector crossing
are verified. An independent origin remains unproved. The most direct next
test is an actual residual update and fit metric, or a resolved fit correction
with its energy partition; a dynamics-based alternative needs an independent
rate and sample clock. Until one is supplied, exact $3/4$ should be described
as a specified or calibrated contraction value, not a geometrically or
physically derived constant. The local existence and closure-limit results
do not require that particular fraction: they hold for the defined event
return at any $q\in(0,1)$.

\section{A proved common upper contraction bound at fixed sampling time}
\label{sec:common-upper-contraction}
The exact event-return construction and a fixed-time upper bound are
different statements. We now prove the latter for the uncalibrated local
normal-form flow. No input-dependent stopping rule is used in this section.

\begin{theorem}[Uniform decay rate for the three collar pieces]
\label{thm:sharper-common-decay}
For every nonzero state in each defined collar domain $A_j$, the flow
$S_\tau(a,b)=(e^{-\tau}a,e^{-\tau}b)$ satisfies
\[
\frac{d\Psi_j(S_\tau x)}{d\tau}
\le-\frac{247}{128}\Psi_j(S_\tau x).
\]
Thus, for all $\tau\ge0$,
\[
\Psi_j(S_\tau x)\le e^{-(247/128)\tau}\Psi_j(x).
\]
The same inequalities hold for a score divided by any fixed positive
reference score. At the zero state the inequalities hold trivially.
\end{theorem}
\begin{proof}
Write $f=f_j(a)$ and $\Psi=a^2+b^2+b^2f(a)^2$. Its Euler derivative is
\[
\mathcal D=a\partial_a\Psi+b\partial_b\Psi
=2\Psi+2ab^2f(a)f'(a),\qquad\dot\Psi=-\mathcal D.
\]
On piece 1, the correction is $a^2b^2/8\ge0$, so
$\mathcal D_1\ge2\Psi_1$.
On piece 2 it is $-(9/32)a(1-a)b^2$. Since
$a(1-a)\le1/4$ and $b^2\le\Psi_2$,
\[
\mathcal D_2\ge2\Psi_2-\frac9{128}b^2
\ge\frac{247}{128}\Psi_2.
\]
On piece 3 it is $a^2b^2(1-a)(1-2a)/2$.
For $a\le1/2$ it is nonnegative. For $a\ge1/2$,
$(1-a)(2a-1)\le1/8$ and $a^2b^2\le\Psi_3$, so
\[
\mathcal D_3\ge2\Psi_3-\frac1{16}a^2b^2
\ge\frac{31}{16}\Psi_3
\ge\frac{247}{128}\Psi_3.
\]
The uniform differential inequality integrates to the exponential estimate.
Domain invariance proved earlier permits its repeated use at every step.
\end{proof}
This improves the earlier conservative common rate $376/219$.
Neither rate is asserted to be the sharp optimal decay constant.

\begin{theorem}[Fixed-step three-quarter upper certificate]
\label{thm:fixed-step-upper-certificate}
Choose one fixed internal sampling interval
\[
\Delta\tau_*:=\frac{128}{247}\log\frac43
\simeq0.1490822076,
\]
and define $U_*:=S_{\Delta\tau_*}$ on every collar piece. Its common
coordinate multiplier is
\[
t_*:=e^{-\Delta\tau_*}=(3/4)^{128/247}
\simeq0.8614982903.
\]
Then for every state in each $A_j$,
\[
\Psi_j(U_*x)\le\frac34\Psi_j(x).
\]
For a nonzero reference $x_0$ and $x_n=U_*^n x_0$,
\[
0<\Phi_j^{\rm NF}(x_n)\le(3/4)^n\quad(n<\infty),
\qquad x_n\longrightarrow(0,0).
\]
The collar/cone inequalities remain satisfied, and detector threshold
$1/256$ is reached no later than step $20$.
\end{theorem}
\begin{proof}
Apply Theorem~\ref{thm:sharper-common-decay} for $\Delta\tau_*$:
$e^{-(247/128)\Delta\tau_*}=3/4$. Iteration gives the normalized
upper bound. The explicitly positive coordinate multiplier gives finite
score positivity and coordinate convergence for a nonzero initial state.
The existing flow invariance gives the readout guards. Finally,
$(3/4)^{20}<1/256$ guarantees the stated upper hitting-time bound.
\end{proof}

The inequality does not imply equality, a first crossing exactly at step
$20$, or a universal measured factor. More generally, any fixed
$\Delta\tau\ge\Delta\tau_*$ gives the same or a stronger upper bound.
For an independently selected interval the proved bound is
$q_{\rm bound}(\Delta\tau)=e^{-(247/128)\Delta\tau}$; $3/4$ is a
convenient selected upper target achieved by the displayed common interval.
The interval is dimensionless model time, not a derived MM/HC physical clock.

\subsection{Check against the three explicit reference states}
For $a_0=b_0=1/2$ and $r_c/\epsilon_r=1/4$, twenty fixed-time steps
were evaluated independently. The largest observed one-step ratios over
those runs were approximately
\[
0.74217412,\qquad0.74597663,\qquad0.74216026,
\]
all below $3/4$. The three normalized scores at step $20$ were
approximately $0.00255156$, $0.00269581$ and $0.00255170$.
Every probe first crossed $1/256$ at step $19$, while retaining the
radial/cone halo inequalities. These are numerical checks for the stated
probes; the uniform theorem covers the complete defined domains.
The Euler-derivative correction identities were also checked symbolically.

\subsection{What this bound can certify natively}
The QAM graph of $U_*$ and its fit readout can now carry a local witness
certificate for an upper contraction bound with fixed sampling. To transfer
it to the inherited C5 system one still needs actual component/norm
identification, native update intertwining with $U_*$ (or a separately
proved stronger native decrease), and native guard/class preservation.
The earlier compatibility theorem for the exact event map does not silently
become a fixed-time theorem: these are different operators.

For the local normal form the result justifies using
$\Phi_{n+1}\le3\Phi_n/4$ as a conservative common bound. It does not
retroactively prove that the original r04 fit-table equality came from
this dynamics. The native scalar equality, the exact event-level
realization and this fixed-time upper certificate retain distinct statuses.

\section{Explicit transfer to collar readout witnesses}
\label{sec:explicit-readout-transfer}
We now complete the local operator and component definitions in the source's
radial/HC readout variables. This supplies a conservative local realization
that can be checked against a native witness, not an assertion that the
older unspecified witness functions have already been identified.
The source symbolic collar formulas are used on their admissible interior
extension as well as their boundary; existence of corresponding native
carriers throughout that extension is a separate realization condition.

\subsection{Readout witness and invertible state bridge}
Fix a piece $j$, passive data $z$, scales $0<r_c<\epsilon_r$,
$\epsilon_h>0$, and declared baseline guard coordinates $g=(g_4^0,g_6^0)$.
Define the new local readout witness
\[
X=(j,\rho_r,\rho_h,z,g),\qquad
r_c\le\rho_r\le\epsilon_r,\quad0\le\rho_h\le\epsilon_h.
\]
These are data tuples, not yet coordinate representatives of every abstract
carrier in the Companion. For each fixed $j,z,g$, the map
\[
\chi_j(X)=(a,b),\qquad
a=\frac{\rho_r-r_c}{\epsilon_r},\qquad
b=\frac{\rho_h}{\epsilon_h[1+f_j(a)]}
\]
is a continuous bijection onto $A_j$, with continuous inverse
\[
\rho_r=r_c+\epsilon_r a,\qquad
\rho_h=\epsilon_h b[1+f_j(a)].
\]
Its Jacobian is nonsingular because $1+f_j(a)>0$.
The rectangle inequalities and the previously defined $A_j$ constraints
are exactly equivalent under this bridge.

\subsection{Executable components and fixed-time operator}
Define the local readout fit components directly from $X$:
\[
R_j^{\rm read}(X)=
\begin{pmatrix}
(\rho_r-r_c)/\epsilon_r\\
\rho_h/\{\epsilon_h[1+f_j(a)]\}\\
\rho_h/\epsilon_h-\rho_h/\{\epsilon_h[1+f_j(a)]\}
\end{pmatrix}.
\]
Their proposed bounded, exterior and mixed roles are those of the new
local normal form. In particular the third component is the difference
between the total normalized HC readout and its independently specified
cone-height component. It is not silently the native C4 leakage term.
Let $\Psi_j^{\rm read}(X)=\|R_j^{\rm read}(X)\|^2$ in the same
dimensionless positive fit norm.

For any fixed $0<t\le1$, define the readout operator by
\[
\widehat F_t(X)=
\left(j,\ r_c+t(\rho_r-r_c),\
t\rho_h\frac{1+f_j(ta)}{1+f_j(a)},\ z,\ g\right).
\]
Every entry is now evaluated from the input witness and the fixed piece.
This is an explicit nonlinear local operator, not the scalar rule
$\Phi\mapsto3\Phi/4$ without a witness action.

\begin{theorem}[Exact readout-to-normal-form transfer]
\label{thm:explicit-readout-transfer}
The new local readout functions and operator satisfy
\[
R_j^{\rm read}(X)=R_j^{\rm NF}(\chi_jX),\qquad
\chi_j(\widehat F_tX)=F_t(\chi_jX).
\]
Consequently, for the fixed $t_*=(3/4)^{128/247}$,
\[
\Psi_j^{\rm read}(\widehat F_{t_*}X)
\le\frac34\Psi_j^{\rm read}(X).
\]
The radial and HC readout inequalities are invariant. The local zero set
is exactly $(\rho_r,\rho_h)=(r_c,0)$. Fixed product guard coordinates
remain unchanged by this defined operator.
\end{theorem}
\begin{proof}
Substituting the inverse readout gives
$R_j^{\rm read}=(a,b,bf_j(a))$.
The new radial readout gives $a'=ta$. The new HC readout gives
\[
b'=\frac{t\rho_h[1+f_j(ta)]/[1+f_j(a)]}
{\epsilon_h[1+f_j(a')]}
=tb.
\]
This proves the two bridge identities. Apply the fixed-step upper
certificate to the identified score. Domain invariance transfers through
the bijection, and the identified zero set follows from $a=b=0$.
The last two coordinates are unchanged by the defined product update.
\end{proof}
For a fixed nonzero reference, normalizing this score by its initial value
gives the same geometric upper bound and step-at-most-20 detector guarantee.
This is a proved transfer to an explicitly defined local witness realization.
It is not yet the inherited native C5 score or inherited MM/HC step.

\subsection{A differential guard check in readout variables}
The continuous witness flow behind $\widehat F_t$ is
\[
\dot\rho_r=-(\rho_r-r_c),\qquad
\dot\rho_h=-\rho_h\left[1+
\frac{a f_j'(a)}{1+f_j(a)}\right],\qquad\dot z=0.
\]
For $0\le a\le1$, the bracket is positive on all three pieces:
\[
g_1(a)=\frac{4+2a}{4+a},\qquad
g_2(a)=\frac{11-6a}{11-3a},\qquad
g_3(a)=\frac{2+2a-3a^2}{2+a-a^2}.
\]
Their lower bounds are respectively $1$, $5/8$ and $1/2$.
Hence the vector field is tangent to the lower radial and HC-zero faces
and points inward at the upper radial and HC faces. This independently
checks the actual readout inequalities rather than only stored flags.
For another native equality guard $G(X)=0$, the required tangency test is
\[
\partial_{\rho_r}G\,[-(\rho_r-r_c)]
-\partial_{\rho_h}G\,\rho_h g_j(a)=0
\quad\text{on its guard set},
\]
with the other passive derivatives included if a native update moves them.
No claim is made that an unspecified native leakage functional passes this test.

\subsection{What must match an actual MM/HC carrier}
Let $W_j^{\rm nat}$ be an actual admitted carrier subset with its native
radial/HC readouts, fixed piece and passive reference. To use the local
realization natively, the following are independent obligations:
\begin{enumerate}
\item The readouts and passive data must parameterize the relevant native
witness states, or all retained residuals and updates must factor through
them. Otherwise hidden carrier data have been discarded.
\item The actual three native C5 residuals must agree with
$R_j^{\rm read}$ through a fixed component-preserving fit isometry and
constant normalization, or verified comparison bounds must be supplied.
\item The actual MM/HC update must induce $\widehat F_{t_*}$ on these
readouts, or a separate stronger decrease must be proved in the same score.
\item The remaining native guards and carrier class must be invariant,
and overlaps between pieces must preserve the stated fit meaning.
\end{enumerate}
These statements are not made true by renaming the newly defined functions.
They require calculations from actual carrier and operator definitions.

\subsection{A transfer theorem with imperfect norm agreement}
\begin{theorem}[Comparison-metric upper-bound transfer]
\label{thm:metric-comparison-transfer}
Suppose a native subset and its chart obey
\[
 c_-\Psi_j(\chi X)\le\Phi_{\rm nat}(X)
\le c_+\Psi_j(\chi X),\qquad 0<c_-\le c_+,
\]
with fixed comparison constants. Positivity at a nonzero chart state
follows from the lower bound. Suppose the native step preserves
the domain and induces $S_{\Delta\tau}$, or has at least its chart-score
decrease. Set $\kappa_G=c_+/c_-$. Then
\[
\Phi_{\rm nat}(F_{\rm nat}X)
\le\kappa_G e^{-(247/128)\Delta\tau}\Phi_{\rm nat}(X).
\]
In particular a native three-quarter upper bound follows if
\[
\Delta\tau\ge\frac{128}{247}\log\left(\frac{4\kappa_G}{3}\right).
\]
\end{theorem}
\begin{proof}
Use the upper score comparison at the output, the chart decay estimate,
and the lower score comparison at the input, in that order. This gives
the product factor $c_+/c_-$; the stated time condition makes it at most
$3/4$.
\end{proof}
Thus comparable native norms are sufficient, but their distortion costs
extra decay. The original $\Delta\tau_*$ alone does not certify $3/4$
when $\kappa_G>1$. Variable reference normalization, state-dependent
weights or changing charts need further uniform control.

\paragraph{Transfer status.}
The local missing definitions have been supplied, and their component
agreement, update intertwining, readout-guard invariance and score-bound
transfer have been proved. The comparison-metric theorem specifies how
to handle a nonidentical native norm. All these results apply to the
declared local realization. Identity with the old MM/HC carrier update,
the full C4/C6 guards or the original OP5 statement remains subject to the
explicit independent obligations above.

\section{Actual native-contract comparison and transfer verdict}
\label{sec:native-contract-verdict}
The source was revisited specifically for native update rules, residual
definitions and admissible domains. This audit checks the conditions of
the transfer theorems instead of assuming them.

\subsection{What matches, and what does not yet match}
\begin{enumerate}
\item \textbf{Symbolic readout chart: matched on the declared pieces.}
The source's normal collar formulas and three symbolic corrections have
the same algebraic form as the local coordinate bridge. Their inverse
readouts and nonzero normalized Jacobians were checked. This establishes
formula agreement on that symbolic fragment, not reconstruction of a
complete native carrier from two readouts.
\item \textbf{Native shift and projection: not executable from the inspected
rules.} The source expands $U_\pi^{(2)}$ and
$\Pi_{{\rm cone}\to{\rm rad}}^{{\rm HC},\gamma}$ as named symbols and
uses them in distance and commutator readouts. The inspected local blocks
do not supply row-specific action matrices or evaluated carrier inputs.
The macro definitions near source line 44126 name the operators; they are
not mathematical evaluation algorithms.
\item \textbf{Native C5 components: not identified.} The r04 three-norm
definition is available, but actual component functions of the selected
carrier are not given by the readout formulas. Assigning them the new
local functions remains an extension, not a proved native equality.
\item \textbf{Admissible geometric domain: needs an interior extension.}
The source certified pieces are stated as subsets of the halo boundary.
The new $A_j$ is a full collar region, and its fixed-time flow can leave a
boundary piece inward. Such a move is harmless for halo containment but
does not define an endomorphism of the original boundary-only chart.
\item \textbf{Native equality recurrence: fails for the fixed-time upper
model under direct score identification.} The original local monitor
tests equality, not merely an upper bound. This distinction is evaluated
explicitly below.
\item \textbf{Native guard coupling: remains conditional.} The actual C4
mixed-kernel projection, clock window, selector switches and overlap maps
must still be evaluated on the carrier update. Preserving passive product
flags is not a substitute for those evaluations.
\end{enumerate}

\subsection{A definite equality-monitor obstruction}
The r36 interface table near source line 59323 uses the signed monitor
\[
E_{\rm col}=\Phi_5^{(n+1)}-\frac34\Phi_5^{(n)},
\]
with zero required for the stated interface certificate. The r50 domain
\texttt{subsec:r50-admissibility-domains} near source line 61805 explicitly
requires $\Phi_5^{(n+1)}=3\Phi_5^{(n)}/4$. The r31 audit functional also
contains $E_{\rm col}^2$.

For the fixed-time local witness at $a_0=b_0=1/2$, the normalized first-step
monitor values on the three pieces are approximately
\[
-0.00930402283529,\qquad -0.00402336563129,\qquad -0.00803930773017.
\]
They are negative rather than zero. Thus even after a hypothetical native
score identification, the new fixed-time step would not meet that literal
old recurrence-domain contract. The squared equality penalty is positive
despite improved contraction.

There is an analytic version of this obstruction. Piece 1 has
$\dot\Psi_1\le-2\Psi_1$, while $247/128<2$. At the selected fixed step,
for any nonzero input,
\[
\frac{\Psi_1(U_*x)}{\Psi_1(x)}
\le e^{-2\Delta\tau_*}
=(3/4)^{256/247}<\frac34.
\]
Therefore its old equality monitor cannot vanish. This is a genuine failure
of literal compatibility, not merely missing numerical precision.

\subsection{The required bounded-contraction monitor extension}
For a proposed revised local contract, define a violation-only monitor
\[
E_{{\rm col},+}:=\max\left(0,
\Phi_{n+1}-\frac34\Phi_n\right),
\]
and the finite-positive domain
\[
\mathbb D_{{\rm col},\le}^{\rm new}
=\{(\Phi_n,\Phi_{n+1}):0<\Phi_{n+1}\le(3/4)\Phi_n\}.
\]
Then the fixed-time normal form has $E_{{\rm col},+}=0$ and preserves
finite positivity. The exact positive recurrence domain is a subset of
this new domain, but the reverse inclusion fails. Consequently this
change is a deliberate relaxation of an equality specification. It is
not a logical consequence of the old domain definition.

In a revised audit functional $E_{{\rm col},+}^2$ can replace the signed
square if the intended task is to penalize insufficient contraction only.
The signed $E_{\rm col}$ should remain available as a diagnostic to record
stronger-than-target lowering. Statements that used exact geometric values,
constant-ratio identities or exact first hitting at step 20 must then be
changed to upper bounds and reaching the threshold by step 20.
This note proposes that extension only for its local model; it does not
silently edit the external v4.7.0 source or its release claims.

\subsection{Boundary charts need a collar, not just a renamed domain}
Consider a source radial outer face $\rho_r=\epsilon_r$ with
$\rho_h<\epsilon_h$. For $0<t<1$, the constructed step gives
\[
\rho_r'=r_c+t(\epsilon_r-r_c)<\epsilon_r.
\]
For a point with positive radial excess and strictly interior HC coordinate,
the HC readout also stays interior. The output is therefore inside the halo,
not on this outer boundary. The required extension is a realized tubular
collar around the certified boundary piece, with the same readout formulas
and guard definitions extended consistently into that collar.
The symbolic algebra supplies a possible such extension; its realization
by actual native carriers and agreement on overlaps is not proved by the
nonzero boundary Jacobian alone.

\subsection{Native monitor values that can be checked at the source anchor}
At the recorded comfortable source packet
\[
(N,\kappa,G)=(5,3/4,1),\quad H=H_\partial=20,
\quad\mathcal C_\gamma(\tau_v)=1,
\]
the native relay and clock monitor formulas give
\[
E_{\rm rel}^{-}=E_{\rm src}^{-}=0,
\qquad E_{\rm clk}=\max(0,4/5-1)=0.
\]
If the actual transition seam is zero and the recorded mixed guard data
have $P_\partial S_{4,\mu}^{\rm mixker}=0$ and $\beta_4=1/256$, their
monitor formulas also evaluate to zero. This does not make the full C4
mixed-kernel remainder vanish; only its specified projection is zero and
the positive ledger bound is retained.
Keeping those \emph{actual source data} fixed in an independent product
extension preserves these evaluated monitor values. If the native update
couples them to the collar coordinates, the proposed product construction
does not establish preservation. The recurrence slot still needs the
contract choice just discussed.

\subsection{Two distinct continuation paths}
\begin{enumerate}
\item \textbf{Keep the original exact recurrence contract.} Use the
score-level event map locally. It gives $E_{\rm col}=0$ by construction,
with state-dependent event times. Its identity with the actual MM/HC step,
the component map and full native guards remains unproved. It does not
derive an exact physical constant-ratio law.
\item \textbf{Use a fixed-time contraction bound.} Adopt the explicitly
new violation-only monitor, extend the boundary chart to a realizable
interior collar, and check native residual/norm and guard agreement.
Then the fixed-time upper-bound theorem is applicable on that checked
local witness subset. This path accepts stronger lowering instead of
rejecting it through an equality penalty.
\end{enumerate}

\paragraph{Native transfer verdict.}
The symbolic chart comparison succeeds. An unconditional transfer to the
old native operator does not: its action and component residuals remain
unidentified, the original chart needs an interior realization, and the
fixed-time candidate fails the old equality-recurrence monitor outright.
The revised inequality monitor removes that particular contract obstruction
as a declared new extension, not as a proof of the old equality.
The anchor monitor evaluations, inward boundary move and recurrence-monitor
counterexample were checked explicitly. These checks narrow the unfinished
native construction instead of treating another fitted model as its solution.

\section{Identifying the native MM/HC objects before transferring dynamics}
\label{sec:native-mmhc-identification}
The source contains more explicit MM/HC algebra than a named transport
symbol alone. The following identification uses its shared-leakage probe,
tensor admissibility and closure-functional definitions. It separates what
can be calculated from those formulas from the still-unspecified C5 fit
update. Source line numbers refer to the inspected external v4.7.0 file;
the source labels are the more stable anchors.

\subsection{A concrete splitting, but not yet a damping generator}
In the source's active linear realization, set
\[
 J^2=I,\qquad C=\tfrac12(I+J),\qquad P_-=I-C.
\]
The native HC object here is the projector $C$, and the native MM object
is the layer transport matrix $\mathcal M$. Direct algebra gives
\[
 C^2=C,\quad P_-^2=P_-,\quad CP_-=P_-C=0,
 \quad JC=C,\quad P_-J=-P_-.
\]
Thus $V=\operatorname{im}C\oplus\operatorname{im}P_-$ is an algebraic
even/odd splitting. Orthogonality additionally requires a compatible
inner product, for example $J^\dagger=J$; involution alone is insufficient.
Indeed
\[
 J=\begin{pmatrix}1&2\\0&-1\end{pmatrix},\qquad
 C=\begin{pmatrix}1&1\\0&0\end{pmatrix}
\]
satisfies the involution and projector identities, but sends $(1,1)$ to
$(2,0)$, increasing the Euclidean squared norm from $2$ to $4$.

Even an orthogonal $C$ is a selector rather than a repeated constant-rate
damper: $C^n=C$ for every $n\ge1$. It preserves an even state exactly and
removes an odd state in one step. Therefore this native HC definition does
not identify the earlier strictly positive three-quarter fit iteration.

The tensor kernel's condition (T4), at
\texttt{def:tensor-structural-kernel} near source line 8797, imposes
$[\mathcal M,J]=[\mathcal M,C]=0$. These equations constrain transport;
they do not specify matrix entries, a physical clock or a dissipative
generator. For instance, in an orthogonal splitting the family
\[
 \mathcal M_s=sC+2P_-,\qquad 0<s\le1,
\]
obeys both equations and has scalar action $sI$ on $\operatorname{im}C$.
Its squared state-norm multiplier there is $s^2$, which is free. This is
an underdetermination example for the displayed algebraic conditions,
not a verification of every other native carrier requirement. Neither
those conditions nor the displayed scalar-action form selects
$s=\sqrt3/2$, and a state norm is not automatically the native C5 score.

\subsection{The leakage residuals are identifiable, with a scope condition}
The source at \texttt{def:working-shared-matrix-leakage}, near line 4773,
defines actual matrix-valued residuals
\[
 E_{\rm comm}=[\mathcal M,J],\qquad E_{\rm leak}=P_-\mathcal M C.
\]
Their relation is explicit:
\[
 [\mathcal M,J]=2(P_-\mathcal M C-C\mathcal M P_-).
\]
On the strict (T4) domain both residuals vanish. Conversely,
$E_{\rm leak}=0$ alone need not make the opposite block
$C\mathcal M P_-$ vanish. The two-channel diagnostic is consequently
meaningful on relaxed or provisional transport data, but cannot be a
nonzero fit variable on the same strict commuting realization without
changing its scope.

This is also why a persistent C4 mixed-kernel ledger value must not be
silently identified with $\|E_{\rm leak}\|$: the source retains C4
separately while imposing (T4) on its admissible tensor converter.
A different projector, carrier layer or readout could distinguish them,
but that connecting map has to be exhibited.

\subsection{The native probe gives an exactly computable response}
At \texttt{def:working-shared-leakage-probe-family}, near source line 4902,
the source supplies a bounded probe matrix $K$, an even input $Cv=v$,
fixed linear masks $L_1,L_5$, and
\[
 \Delta_K=P_-KC+CKP_-,\quad
 \mathcal M_\lambda=\mathcal M+\lambda\Delta_K,\quad
 y_\lambda=\mathcal M_\lambda Cv.
\]
For these supplied definitions, not just for a formal residual analogy,
\[
 P_-y_\lambda=P_-\mathcal M Cv+\lambda P_-KCv.
\]
Writing
\[
 d(\lambda)=
 \begin{pmatrix}L_1(P_-y_\lambda)\\L_5(P_-y_\lambda)\end{pmatrix},
 \qquad
 a=\begin{pmatrix}L_1(P_-KCv)\\L_5(P_-KCv)\end{pmatrix},
\]
gives the exact affine relation $d(\lambda)=d(0)+\lambda a$.
Its image has rank at most one. The masks' coefficients, $K$, $v$ and
the entries of $\mathcal M$ remain unevaluated, so this is an executable
algebraic formula once those data are supplied, not a measured response.

\begin{proposition}[Reversal cannot supply a second probe direction]
\label{prop:native-reversal-probe-rank}
For every $K$ in this fixed linear probe setup, the reversal-paired matrix
$\widetilde K=JK$ gives $\widetilde a=-a$. If $a\ne0$ and the masks have
real outputs, the source's normalized determinant and absolute dot-product
defects therefore satisfy
\[
 \mathfrak d_{\parallel}=0,\qquad \mathfrak d_{\perp}=1.
\]
\end{proposition}
\begin{proof}
Since $P_-J=-P_-$,
$P_-\widetilde KCv=P_-JKCv=-P_-KCv$. Apply the same linear masks.
The response matrix has rows $a^T$ and $-a^T$, hence zero determinant
and absolute dot product $\|a\|^2$; both row norms are $\|a\|$.
\end{proof}
The source already calculates this result for its canonical low-rank
class at \texttt{prop:working-response-geometry-canonical-low-rank-class},
near line 5114. The calculation above extends it to \emph{all} $K$ under
the stated fixed projector and linear-mask definitions. Increasing the
rank of $K$ alone does not rescue the earlier $JK$ quadrature diagnostic.
An unnormalized small dot product can merely reflect a small response.

The source's later two-directional construction instead introduces a
different operator $R_W$ on an odd plane, with
$R_Ww_1=w_2$ and $R_Ww_2=-w_1$. That is a genuine quarter-turn, whereas
$J|_{\operatorname{im}P_-}=-I$ is not. The resulting outputs can be
orthogonal on the carrier plane. Their masked visible outputs are
orthogonal only if the masks preserve that geometry on the plane, for
example through a scaled isometry. These are distinct constructions;
neither constitutes an identified C5 lowering step.

\subsection{What damping the identified probe would actually require}
For real masks and a declared fixed positive fit metric $G$, let
\[
 E(\lambda)=d(\lambda)^TGd(\lambda),\qquad
 \beta=a^TGa>0,\qquad
 \lambda_*=-\frac{a^TGd(0)}{\beta}.
\]
Completing the square gives
\[
 E(\lambda)=E_{\rm floor}+\beta(\lambda-\lambda_*)^2,
 \quad E_{\rm floor}=d(\lambda_*)^TGd(\lambda_*)\ge0.
\]
Thus one scalar probe can remove only the component along its response
direction. It cannot remove a nonzero metric-orthogonal floor.

A \emph{new}, explicitly declared relaxation
\[
 \lambda'=\lambda_*+(1-\eta)(\lambda-\lambda_*),\quad 0<\eta<1,
\]
would give
\[
 E(\lambda')-E_{\rm floor}
 =(1-\eta)^2\bigl(E(\lambda)-E_{\rm floor}\bigr).
\]
It stays in the source's bounded probe interval if both the input and
$\lambda_*$ lie there. Taking $\eta=1-\sqrt3/2$ makes the
\emph{excess above the floor} contract by $3/4$. If the floor is positive,
\[
 E(\lambda')=\tfrac34 E(\lambda)+\tfrac14 E_{\rm floor}
 >\tfrac34 E(\lambda).
\]
This identifies a possible relaxation of the actual affine probe, but
not an inherited relaxation rule or native C5 contraction. The chosen
metric and relaxation parameter are additional data. In particular a
positive floor obstructs decay of the full probe score to zero.

\subsection{A concrete guard-coupling test, not passive guard flags}
Changing the probe parameter changes the matrix leakage by
\[
 E_{\rm leak}(\lambda')-E_{\rm leak}(\lambda)
 =(\lambda'-\lambda)P_-KC.
\]
Starting from a strictly commuting $\mathcal M$, preservation of (T4)
at a nonzero $\lambda$ requires both cross blocks of $\Delta_K$ to
vanish, hence $\Delta_K=0$. A nontrivial cross-sector probe therefore
leaves that strict domain. This does not invalidate its diagnostic use
on provisional data, but precludes treating it as an admissible (T4)
fit step without a separately justified domain extension.

If instead a relaxed native guard is realized by a fixed linear mask
$Q_4$ of the complementary output, exact preservation on this input
requires
\[
 (\lambda'-\lambda)Q_4P_-KCv=0.
\]
Preserving the \emph{entire} leakage output during a nontrivial parameter
change would force $P_-KCv=0$, and consequently $a=0$: this particular
probe could no longer lower its masked residual. Preserving only selected
guard components can leave a fit direction available, but requires
$Q_4P_-KCv=0$ together with a nonzero fit response.
The source's actual $P_\partial S_{4,\mu}^{\rm mixker}$ and leakage
budgets have not been identified with $Q_4P_-y_\lambda$. The displayed
kernel test is conditional on that additional identification, not a
claim that every C4-preserving C5 step is impossible.

\subsection{The positive HC closure functional is a different object}
Near source line 8077 the source defines
\[
 \mathcal H_C=\rho[\mathcal M_{16},Q]^\dagger[\mathcal M_{16},Q]
 +\sigma\mathfrak R_Q^\dagger\mathfrak R_Q,\quad \rho,\sigma\ge0,
\]
where $Q=\hat P_{C_7}$ and
$\mathfrak R_Q=(I-Q)\mathcal M_{16}Q+Q\mathcal M_{16}(I-Q)$.
Neither $Q=C$ nor $\mathcal H_C=C$ follows from these definitions.

For an orthogonal $Q$, in its adapted splitting write
\[
 \mathcal M_{16}=\begin{pmatrix}A&B\\D&F\end{pmatrix}.
\]
Then direct block multiplication yields
\[
 [\mathcal M_{16},Q]=\begin{pmatrix}0&-B\\D&0\end{pmatrix},
 \quad\mathfrak R_Q=\begin{pmatrix}0&B\\D&0\end{pmatrix},
\]
and therefore
\[
 \mathcal H_C=(\rho+\sigma)
 \begin{pmatrix}D^\dagger D&0\\0&B^\dagger B\end{pmatrix}.
\]
This is an explicit positive mismatch operator. It does not specify that
a state or matrix evolves by its negative gradient. Even the source's
displayed variational penalty does not by itself supply such a
dissipative relaxation or an effective viscosity. Those require a
generator, metric, clock and a checked readout/coarse-graining map.

\subsection{Why the three C5 fit components still cannot be identified}
The r04 score's three functions occur in the declaration at
\texttt{def:v470-r04-c5-lowering-score}; the searched literal component
names have no separate evaluated definitions elsewhere in the inspected
file. The corresponding r05 functions are declared again with a
second-step superscript. Their native fit-package arguments are not
specified as $\mathcal M_\lambda$, $y_\lambda$ or the normal-collar readouts.
Consequently neither the three-component map nor its fit metric is
recoverable from the probe formulas by a stated identity.

Even at a single ledger row, a positive scalar value $\Phi$ cannot supply
the three components: for arbitrary $p_i\ge0$ with $\sum_i p_i=1$, scalar
components $\sqrt{p_i\Phi}$ have that same sum of squared norms. This is
only an algebraic ambiguity demonstration, not a proposed assignment of
the native components. The ledger supplies aggregate values and guard
classifications, not the missing component evaluations or their derivative
along an actual update.

\paragraph{Identification result and next irreducible calculation.}
The native HC projector, MM admissibility identities, two matrix leakage
residuals, affine probe response and positive closure functional are now
identified explicitly. The reversal response rank, the probe energy floor
and its leakage change have been derived, rather than assumed. What remains
unidentified is the link from an actual C5 fit package to those objects.
The next useful input is one realized bounded witness $W$ with
$R_5(W)=(\rho_{5,\rm bd},\rho_{5,\rm ext},\rho_{5,\rm mix})(W)$ and
an evaluated update $W\mapsto W'$. On those same data one must compute
$R_5(W')$ and the actual guard tuple $\Gamma_5(W')$, not borrow a
cross-sector probe that already changes the strict MM domain.
Only then can the earlier normal-collar or fluid model be compared to
the native action. The source's own \texttt{hyp:fluid-readout}, near
line 10874, likewise leaves the concrete coarse-graining map open.

\paragraph{Verification.}
The projector, cross-block, reversal, commutator and closure-functional
identities were checked by symbolic matrix multiplication on general
two-by-two block realizations. The affine-score square completion and
relaxation identity were checked independently as polynomial identities.
These checks validate the displayed calculations, not the missing native
fit-package identification.

\section{Following the C5 fit provenance to an explicit guarded update}
\label{sec:c5-bifocal-update-investigation}
The preceding negative identification result concerns a direct substitution
of the cross-sector leakage probe for the C5 fit. The source points to a
different, more concrete route: its finite witness-fit register inherits
the inverse lens-fit calibration method. This section investigates that
route and constructs a conditional readout update from its actual defect
formula. It is not a claim that the r04 ledger's missing witness data have
been recovered.

\subsection{The provenance chain selects a fit, not matrix leakage}
The relevant source anchors are
\texttt{def:v470-r01-promotion-ledger} near line 51704,
\texttt{def:v470-r01-extmix-field-equation-candidate} near line 51804,
and \texttt{def:v470-r02-extracted-candidate-family} near line 51960.
They promote inverse lens-fit calibration as a finite test method, assign
that provenance to $\mathfrak R_{\rm fit}$, and define
\[
 C_5={}^{\rm fit}\mathcal T_{\mathfrak R_{\rm fit}}^{\rm bd}.
\]
In contrast $C_4$ is the genuinely mixed MM/HC term. Thus the source's
own typing supports testing C5 as a bounded calibration of a readout;
it does not require deforming the strict commuting MM matrix.

This is a method-level link, not yet a component-level link. The r02 C5
tuple is $(0,0,\epsilon_{\rm mix}^{\rm fit},0,\rho_{\rm fit})$, whereas
r04 introduces three squared component norms without an evaluation map
between those two presentations. The r02 scope remark explicitly keeps
older lens and route fragments as provenance until mapped into the
current candidates. Therefore the following component construction is
a proposed implementation of that link, not an inherited identity.

\subsection{A residual map derived from the displayed lens formula}
The source's \texttt{sec:lens-extreme-local-refitting}, near line 47978,
uses symmetric foci $F_-=-ce$, $F_+=ce$ and adjusts an effective semi-major
value. On a declared Euclidean chart with fixed samples $x_i$, unit axes
$e_i$, and positive constant weights $w_i$, write that value as $A$ and set
\[
 S_i(c)=\|x_i+ce_i\|+\|x_i-ce_i\|,\qquad
 \ell_i(A,c)=S_i(c)-2A.
\]
The source defect is $|\ell_i|$. Squaring it equals $\ell_i^2$, so a
signed residual is legitimate for the squared fit objective, including
at zero defect. Fix a disjoint, exhaustive assignment of samples to
$I_{\rm bd},I_{\rm ext},I_{\rm mix}$ and define
\[
 R_{\alpha}^{\rm lens}(A,c)
 =\bigl(\sqrt{w_i}\ell_i(A,c)\bigr)_{i\in I_{\alpha}},
 \qquad \alpha\in\{{\rm bd},{\rm ext},{\rm mix}\}.
\]
Then the explicitly realized three-block score is
\[
 E(A,c)=\sum_\alpha\|R_\alpha^{\rm lens}\|^2
       =\sum_iw_i\ell_i(A,c)^2.
\]
The source supplies neither this channel assignment nor the weights for
its r04 witness. They must be attached to actual fit data and kept fixed
throughout the update. Relabelling arbitrary samples is not proof of their
bounded, exterior or mixed meaning. For general metric charts the distance
formula still makes sense, but the Euclidean derivatives below require
this additional chart realization.

\subsection{The actual parameter response and its identifiability test}
Away from samples coinciding with a focus, put
\[
 s_i=x_i\cdot e_i,\quad h_i^2=\|x_i\|^2-s_i^2,\quad
 d_{i,-}=\sqrt{(s_i+c)^2+h_i^2},\quad
 d_{i,+}=\sqrt{(s_i-c)^2+h_i^2}.
\]
The weighted residual Jacobian has rows
\[
 \bigl(DR^{\rm lens}\bigr)_i
 =\sqrt{w_i}\,(-2,\tau_i),\qquad
 \tau_i=\frac{s_i+c}{d_{i,-}}+\frac{c-s_i}{d_{i,+}},
 \qquad |\tau_i|\le2.
\]
Moreover
\[
 \partial_c^2 S_i
 =h_i^2(d_{i,-}^{-3}+d_{i,+}^{-3})\ge0.
\]
With $W=\sum_iw_i$, $T=\sum_iw_i\tau_i$ and
$U=\sum_iw_i\tau_i^2$, the Gram matrix and its determinant are
\[
 (DR^{\rm lens})^TDR^{\rm lens}
 =\begin{pmatrix}4W&-2T\\-2T&U\end{pmatrix},\qquad
 \det=4(WU-T^2).
\]
Thus both parameters are locally identifiable exactly when the
$\tau_i$ are not all equal on the positive-weight samples. At $c=0$
all $\tau_i=0$, so focal separation is not first-order identifiable
there. Having three residual blocks does not imply three independent
adjustable degrees of freedom; here there are at most two.

\subsection{An exact bounded update without changing the scan carrier}
Hold $c$ fixed. Define
\[
 A_*(c)=\frac{\sum_iw_iS_i(c)}{2W},\qquad
 E_{\rm floor}(c)=\sum_iw_i\bigl(S_i(c)-2A_*(c)\bigr)^2.
\]
Completing the square now uses an actual source-supported fit parameter:
\[
 E(A,c)=E_{\rm floor}(c)+4W(A-A_*(c))^2.
\]
For $0\le t<1$ define
\[
 A'=A_*+t(A-A_*),\qquad c'=c.
\]
If $A$ and $A_*$ lie in the same admissible interval, the whole segment
stays in that interval. It gives the exact identity
\[
 E(A',c)=E_{\rm floor}+t^2(E(A,c)-E_{\rm floor}).
\]
There is no need to change $\mathcal M$, $J$, $C$, the sampled carrier
or the protected core to implement this readout calibration.

In particular the source's actual geometric guard is
\[
 g_i=d(x_i,K)-\varepsilon,
\]
from \texttt{def:finite-bifocal-guarded-distance-registers} near line 930.
Fixing $x_i,K,\varepsilon$ preserves each $g_i$ exactly, not merely to
first order. The index-level parity identities also remain unchanged
when the cyclic sample indices and their transport are fixed. If
$c=\kappa a_s$ with fixed inherited scale $a_s$, fixing $c$ also fixes
the bounded focal-ratio parameter $\kappa$.

This proves geometric guard preservation for the actual displayed
guard formula. Full r04/r05 preservation additionally requires verifying
that the MM/C4 data, selector assignments and all-readout block do not
depend on $A$, or evaluating their dependence explicitly. The source's
separation requirement permits a fit-only implementation but does not
by itself provide those dependency maps.

\begin{proposition}[Exact three-quarter attainability for the lens fit]
\label{prop:lens-quarter-attainability}
Suppose $E(A,c)>E_{\rm floor}$ and the above interval condition holds.
For $0<q<1$, an update of this form has $E(A',c)=qE(A,c)$ if and only if
$E_{\rm floor}\le qE(A,c)$. Its nonnegative relaxation factor is then
\[
 t=\sqrt{\frac{qE(A,c)-E_{\rm floor}}{E(A,c)-E_{\rm floor}}}.
\]
For $q=3/4$, a positive floor therefore limits how far the exact recurrence
can continue at fixed $c$.
\end{proposition}
\begin{proof}
The segment's attainable scores range from $E_{\rm floor}$ to the
input score. Solving the exact identity for $t^2$ gives the displayed
condition and factor. If an infinite recurrence at fixed $c$ had score
$q^nE(A_0,c)$, it would eventually fall below any positive fixed floor,
which is impossible.
\end{proof}
If the floor is positive, $t=\sqrt3/2$ contracts only the excess above
the floor and gives $E'=3E/4+E_{\rm floor}/4$. If the floor is zero, the
same factor contracts each signed residual, each fixed channel block
and the entire score exactly. The zero-floor condition is equivalent
to all sampled focal sums $S_i(c)$ being equal. It is satisfied by an
exact bifocal ellipse at that focal separation, but is not guaranteed by
the source's approximate finite fit summaries.

\subsection{An exact analytic test of the proposed implementation}
This test is constructed here; it is not the source's 315-point refit
or its unevaluated r04 bounded witness. Set $c=1$, $e_i=(1,0)$ and take
the six points
\[
 (2,0),\ (-2,0),\ (0,\sqrt3),\ (0,-\sqrt3),\
 (1,3/2),\ (-1,-3/2).
\]
All lie on the ellipse $x^2/4+y^2/3=1$ and have $S_i(1)=4$.
Choose unit weights and, solely to test the three-block interface, assign
the consecutive pairs to its three blocks. With $A_0=9/4$,
\[
 A_*=2,\quad E_{\rm floor}=0,\quad E(A_0,1)=3/2.
\]
The explicit first update
\[
 A_1=2+\frac{\sqrt3}{8}
\]
gives $E(A_1,1)=9/8$, hence the normalized score
$\Psi=E/(3/2)$ changes from $1$ to $3/4$. More generally
\[
 A_n=2+\frac14\left(\frac{\sqrt3}{2}\right)^n,\qquad
 \Psi_n=\left(\frac34\right)^n>0\quad(n<\infty).
\]
For $K=0$ and $\varepsilon=1/4$, the fixed minimum geometric guard is
$\sqrt3-1/4>0$ throughout. The readout parameter changes; the scan
points do not move toward the core. On this sample the Jacobian's focal
responses are $(0,0,1,1,4/5,4/5)$ and its Gram determinant is $672/25>0$.

This realizes the earlier squared-norm contraction through the source's
bifocal residual formula rather than through an arbitrary abstract
three-vector. However the sample, channel assignment and normalization
are declared test data. Their correspondence with the original C5
components and all native guards still needs proof. The factor
$\sqrt3/2$ was chosen to test the ledger value: any $\sqrt q$ with
$0<q<1$ would produce the corresponding exact normalized recurrence.
Thus the test establishes compatibility, not selection of $3/4$.

\subsection{The constrained test when focal adjustment is also needed}
If the fixed-$c$ floor obstructs lowering, the computed second Jacobian
column supplies a concrete additional fit direction. Whether that
direction is allowed is determined by the actual guard dependencies.
On a regular smooth stratum of equality guards $g(\theta)=g_0$, with
strict margins for the remaining inequalities, let $B$ have orthonormal
columns spanning $\ker Dg(\theta)$, and let
\[
 H=DR(\theta)B,\qquad P_H=HH^+.
\]
Here $R$ is the fixed weighted residual vector, not a changing score
normalization, and $H^+$ is its Moore--Penrose inverse. The tangent
direction $v=-BH^+R$ satisfies
\[
 Dg\,v=0,\qquad DR\,v=-P_HR,\qquad
 D\|R\|^2[v]=-2\|P_HR\|^2.
\]
Therefore a strict first-order guard-compatible lowering direction exists
exactly when $P_HR\ne0$. If it is zero, no allowed tangent direction has
strictly negative first derivative; higher-order decrease is not excluded.
Regularity supplies a local curve in the guard level set with this tangent,
and continuity preserves strict inequality margins for a sufficiently
small step. A raw finite Euler step tangent to a nonlinear guard is not
automatically guard-preserving. Active inequality faces require a tangent
cone analysis rather than this interior-stratum test.

For the lens implementation, $DR$ is the explicit two-column matrix
above. This turns the missing guard question into an evaluable rank and
projection test once the source's actual fit package is supplied. It does
not replace the missing data by a newly assumed invariant flag.

\paragraph{Investigation result.}
There is now a source-supported conditional realization of C5 lowering:
bounded bifocal readout calibration, with an exact update, an explicit
Jacobian, a geometric guard-preservation proof and a sharp fixed-focal
floor test. This route avoids the strict-MM obstruction of the leakage
probe. It can realize the three-quarter recurrence on a zero-floor sample,
but neither forces that numerical factor nor identifies the old ledger
witness. Required native data have narrowed to the sampled fit package,
its fixed channel/weight assignment and the non-geometric guard dependency
maps. Published point counts and mean absolute defects alone do not supply
those quantities or a squared-norm contraction certificate.

\paragraph{Verification.}
The bifocal first and second derivatives, Gram determinant, weighted square
completion and relaxation identity were checked symbolically. The six-point
focal sums, two score values, positive core reserve and Jacobian determinant
were also checked exactly. These are checks of a declared analytic fit
realization, not empirical evidence of fluid motion or a recovered r04 run.

\section{Coherent C5 readouts and coupled guard preservation}
\label{sec:coherent-c5-coupled-guards}
The preceding fit construction leaves two separate questions: how its
residual is read into C5 components, and whether changing a readout can
move a guard projector or relay parameter. This section derives conditions
that address both questions and evaluates a coupled finite test packet.
It does not identify that packet with the missing historical r04 witness.

\subsection{A common mode removes dependence on arbitrary partitions}
For a zero-floor sample at fixed $c$, write $\delta=A-A_*>0$ and let
$z(A)=(\sqrt{w_i}\ell_i(A,c))_i$ be the full signed lens residual.
Then $z(A)=-2\delta(\sqrt{w_i})_i$. Under
$A'=A_*+t\delta$ one has $z(A')=tz(A)$ exactly.

\begin{proposition}[Componentwise transfer of a coherent fit mode]
\label{prop:coherent-component-transfer}
Suppose the actual three C5 component readouts, on the realized fit
family, factor as $\rho_{5,\alpha}(A)=L_\alpha z(A)$ through fixed
linear maps, with fixed norms and reference normalization. Then
\[
 \rho_{5,\alpha}(A')=t\rho_{5,\alpha}(A),\qquad
 \Phi_5(A')=t^2\Phi_5(A).
\]
No sample partition or orthogonality between the $L_\alpha$ is needed.
\end{proposition}
\begin{proof}
Each linear map commutes with multiplication by $t$. Homogeneity of
each squared norm gives the score identity after summing the components.
\end{proof}
This does not recover the unspecified native maps. It shows that their
individual coefficients need not be known to transfer contraction once
the fixed linear factorization itself is proved. At least one map must
detect the common mode for the finite score to be positive.
Three differently typed readouts can thus share one adjustable amplitude
without being three independent dynamical degrees of freedom.

\subsection{A persistent offset cannot be hidden in the C5 fit}
An affine readout has the more general form
$R_5(\delta)=b+\delta a$. In a fixed positive fit metric, an exact
recurrence $\Phi_5(t\delta)=t^2\Phi_5(\delta)$ throughout a nonzero
orbit $\delta_n=t^n\delta_0$, $0<t<1$, forces $b=0$.
Indeed the recurrence makes the score tend to zero, while continuity
makes its limit $\|b\|^2$. Positive definiteness then gives $b=0$.
For a single step an accidental cancellation can occur, so one ledger
ratio alone is insufficient to exclude an offset.

In particular, if a persistent C4 datum $s_4$ contributes a constant
offset $b=L_4s_4$ to C5, the infinite exact C5 recurrence cannot hold
unless this contribution vanishes in the fit readout. This is compatible
with the source's separation rule: C4 remains a separate nonzero datum;
omitting it from C5 is not the same as cancelling or absorbing it.
The required check is that the actual C5 map annihilates that guarded
contribution, not that the full C4 residual vanishes.

\subsection{A moving C4 projector exposes a previously hidden coupling}
The actual numerical C4 interface monitor contains
$\|P_\partial S_{4,\mu}^{\rm mixker}\|^2$.
Keeping the residual fixed is insufficient if its projector depends on
the fit chart. For example let
\[
 P_0=\begin{pmatrix}1&0\\0&0\end{pmatrix},\qquad
 s_0=\begin{pmatrix}0\\d\end{pmatrix},\quad d>0,
\]
and let $Q(\theta)$ be a planar rotation. Set
$P(\theta)=Q(\theta)P_0Q(\theta)^T$ but keep $s_0$ fixed. Then
\[
 P_0s_0=0,\qquad
 \|P(\theta)s_0\|^2=d^2\sin^2\theta.
\]
Thus a chart-dependent projection can produce guard debt even without
changing the amplitude of the persistent residual.

\begin{proposition}[Covariant C4 guard preservation]
\label{prop:covariant-c4-guard}
Suppose a realized chart update supplies an invertible map $T$ with
\[
 P_\partial'=TP_\partial T^{-1},\qquad s_4'=Ts_4.
\]
If $P_\partial s_4=0$, then $P_\partial's_4'=0$. If its comparison
metrics satisfy $T^\dagger G'T=G$, the full residual norm is unchanged.
In particular $s_4\ne0$ remains nonzero.
\end{proposition}
\begin{proof}
Multiply the two transport identities to obtain
$P_\partial's_4'=TP_\partial s_4=0$.
The metric identity gives
$(s_4')^\dagger G's_4'=s_4^\dagger Gs_4$; invertibility preserves
nonvanishing.
\end{proof}
For the rotation example, co-transporting $s_4'=Q(\theta)s_0$
keeps the projection zero and the positive norm $d$ unchanged. This is
a checked cancellation of \emph{projection debt}, not cancellation of
the C4 residual. The positive ledger bound $\beta_4=1/256$ can remain
unspent. For a general invertible chart map, retaining the old metric
instead of transforming it need not preserve that bound.

The same conjugation preserves the algebraic MM/HC conditions if the
actual carrier operators also transform consistently:
\[
 \mathcal M'=T\mathcal M T^{-1},\quad J'=TJT^{-1},
 \quad C'=TCT^{-1}.
\]
Then $(J')^2=I$, $(C')^2=C'$, $C'=(I+J')/2$, and each commutator
is the conjugate of its original value. This is a representation-level
calculation, not a proof that the source's actual fit update uses such
a $T$ or preserves every native closure predicate. That map still has
to be matched to the real chart transition.

\subsection{The focal and relay ratios must not be conflated}
The source uses a focal ratio in $c=\kappa a_s$ and a relay coefficient
in $S=N+\kappa N(N-1)$. Denote them here by
$\kappa_{\rm lens}$ and $\kappa_{\rm rel}$. Their identical typography
does not prove that they are the same function of a fit package.
The lens formula fixes inherited $a_s$ while adjusting effective $A$;
in that reading fixing $c$ leaves $\kappa_{\rm lens}$ unchanged.

As a deliberate coupling stress test, instead impose
\[
 \kappa_{\rm rel}(A)=c/A,\qquad N_{\rm rel}=5,\qquad G=1.
\]
This is additional structure, not a source-derived identity. Under it
the preceding six-point example with $c=1$, $A_0=9/4$ and $A_*=2$
has $H_0=125/9<16$ and limiting $H_*=15<16$.
Its geometric guard passes but its relay guard fails. This quantifies
why that example alone did not establish full native admissibility.

\subsection{A coupled witness satisfying the numerical interface guards}
The same coupling admits a different finite test. Set $c=1$ and sample
the ellipse $x^2/(25/16)+y^2/(9/16)=1$ at
\[
 \begin{gathered}
 (5/4,0),\ (-5/4,0),\ (0,3/4),\ (0,-3/4),\\
 (5/8,3\sqrt3/8),\ (-5/8,-3\sqrt3/8).
 \end{gathered}
\]
All six focal sums equal $5/2$. With unit sample weights choose
\[
 A_*=5/4,\qquad A_0=4/3,\qquad
 A_n=\frac54+\frac1{12}\left(\frac{\sqrt3}{2}\right)^n.
\]
Then
\[
 E_0=1/6,\quad E_1=1/8,\qquad
 E_n/E_0=(3/4)^n>0\quad(n<\infty).
\]
Under the coherent component factorization, normalize the detected C5
score by its fixed input value; it has the same exact recurrence even
when the three component maps are not a partition of these samples.

The fit domain is $5/4<A\le4/3$, excluding the exact-fit endpoint for
finite witnesses. The update $A'=5/4+(\sqrt3/2)(A-5/4)$ maps that
domain into itself. Under the declared relay coupling,
\[
 \frac34\le\kappa_{\rm rel}(A)<\frac45,\qquad
 20\le H(A)=5+20/A<21.
\]
Set the boundary-source score $H_\partial(A)=H(A)$ in this test packet
and retain the clock readout $\mathcal C_\gamma(\tau_v)=G=1$.
The source's actual monitor formulas then give, at every finite stage,
\[
 E_{\rm rel}^{-}=E_{\rm src}^{-}=E_{\rm cmf}^{-}=0,
 \qquad E_{\rm clk}=\max(0,16/H(A)-1)=0.
\]
The relay surplus is not driven to zero; its comfortable margin is
preserved while $H$ increases. The number $N_{\rm rel}=5$ counts relay
packets here, not the six sampled lens points.

For $K=0$ and $\varepsilon=1/4$, the minimum sampled core reserve is
$1/2>0$ and remains fixed. Couple the mixed-sector chart frame to the
fit through $T(A)=Q(A_0-A)$, with $P_\partial(A)=T(A)P_0T(A)^{-1}$
and $s_4(A)=T(A)s_0$, choosing $0<d\le1/256$ in this Euclidean
test metric. Its projection is zero throughout, while $s_4$
remains nonzero with unchanged norm. Retain $\beta_4=1/256$ as its
separate bound record. Staying in one declared local chart with no new
transition seam and the same two-sided source type leaves their test
monitors zero; retain $\chi_{\rm all}=0$ without a promotion step.

Consequently the six numerical source interface monitors evaluated on
this \emph{declared compatible packet} satisfy
\[
 (R_{\rm bd},R_{\rm src},R_{\rm seam},R_{\rm mix},
   R_{\rm col},R_{\rm clk})=(0,0,0,0,0,0)
\]
at each finite update. The mixed component uses the covariant projector
proof, the relay and clock components use the coupled numerical formulas,
and the collar component uses the exact recurrence. This is stronger than
checking only the core distance. Source typing and absence of a seam are
declarations of this single-chart construction, not a coverage or overlap
theorem for the original atlas.

The construction is a concrete realization of the source's monitor-level
specification, conditional on the stated C5 readout factorization. The
source's r50 domain is explicitly a product of monitor domains, not a
physical phase space. Passing its numerical monitors therefore does not
identify the original r04 witness, establish halo/atlas invariance or
prove that a native carrier evolves by this update.

\subsection{If the actual C5 readout is nonlinear}
Fixed linear factorization is sufficient but not necessary for an upper
bound. Suppose the combined normalized native residual $R_5$ is twice
continuously differentiable along the fit family, with
\[
 R_5(A_*)=0,\quad a=R_5'(A_*),\quad m=\|a\|>0,
 \quad\|R_5''(A)\|\le L.
\]
For $\delta=A-A_*>0$, Taylor's estimate gives
\[
 \delta(m-L\delta/2)\le\|R_5(A)\|
 \le\delta(m+L\delta/2).
\]
At $A'=A_*+t\delta$, and for
$\epsilon=L\delta_{\max}/(2m)<1$, this yields the uniform bound
\[
 \frac{\Phi_5(A')}{\Phi_5(A)}
 \le t^2\left(\frac{1+t\epsilon}{1-\epsilon}\right)^2.
\]
In particular, as $\delta\to0$ the ratio tends to $t^2$, but need not
equal it at finite input. A three-quarter upper bound is guaranteed if
\[
 t(1+t\epsilon)\le\frac{\sqrt3}{2}(1-\epsilon).
\]
For a known positive curvature allowance this selects a smaller step
factor than $\sqrt3/2$. The input score stays positive by the lower
estimate. The coupled guard domain above is preserved for every
$0<t<1$, so this stronger calibration can retain the checked guards.
It would require the earlier inequality-monitor extension rather than
silently passing the old exact-equality collar monitor.

This estimate depends on actual values of $m,L$ and on a genuine zero
of the native residual at the reference fit. The lens zero alone does
not establish that native zero. Thus nonlinear component maps can be
handled quantitatively once supplied, without assuming that they are
exactly the constructed lens residuals.

\paragraph{Continuation result.}
The component issue has narrowed to coherent fixed readout factorization,
or verified local derivative bounds, together with exclusion of a
persistent C4 offset from the C5 fit. The guard issue has narrowed to
actual chart covariance and parameter dependencies. A fully evaluated
coupled monitor test now exists, with finite-positive exact C5 decay,
positive core reserve, comfortable relay margin, admissible clock and
nonzero unspent C4 remainder. All are properties of the declared local
implementation. Recovering the historical native witness still requires
its missing component and transition maps; physical fluid dynamics and
global OP5 closure do not follow from this monitor realization.

\paragraph{Verification.}
The coupled sample's focal sums, scores, relay endpoints and core reserve
were checked exactly. Symbolic rotation algebra verifies both the moving-
projector failure and covariant repair. The affine offset obstruction was
checked by polynomial expansion. The nonlinear bound follows from the
stated Taylor remainder estimate and uses a fixed fit norm.

\section{An explicit native-readout bridge and its overlap correction}
\label{sec:explicit-native-readout-bridge}
The attempted link can be made more concrete than an agreement of scalar
ledger values. This section constructs a carrier on which the source's
radial and HC halo formulas are actually evaluated, and connects its
readouts to the lens fit. The construction is a new finite realization
of the inspected formulas, not a recovery of the source's unevaluated
r04 witness or a derivation of its native MM update.

\subsection{What the original transition data do and do not provide}
The source's halo overlap layer near line 45132 declares maps
$\Phi_{jk}:U_{jk}\to U_k$ and distortion budgets, but does not give
their coordinate actions in the inspected symbolic certificate. Its
r26 transition audit, near line 57338, gives typed arrows
$\tau_{01},\tau_{12},\tau_{23}$ and a zero-debt ledger. Those arrows
are not numerical matrices or derivatives of the halo chart overlaps.
One cannot identify them with the previously chosen rotations merely
because both have zero monitor debt.

Similarly, the source's compressed OP theorem near line 21966 assumes
linearity of a closure-compatible observable $\mathcal I$. That does
not assert linearity of the three later C5 fit functions. The latter
still require their own readout map. The following bridge therefore
states its construction data and its remaining identification obligation.

\subsection{Overlap candidates forced by the two collar readouts}
Fix a common passive reference and common $r_c,\varepsilon_r,
\varepsilon_h$. Write
\[
 f_1(a)=a/4,\quad f_2(a)=3(1-a)/8,\quad
 f_3(a)=a(1-a)/2,\qquad p_j(a)=1+f_j(a)>0.
\]
The source symbolic readouts have the form
\[
 \rho_r=r_c+\varepsilon_r a,\qquad
 \rho_h=\varepsilon_h b p_j(a).
\]
On interior collar extensions, their common readout coordinates are
$y=(a,v)$ with $v=\rho_h/\varepsilon_h$. Thus
\[
 \chi_j(a,b)=(a,bp_j(a)),\qquad
 H_{kj}:=\chi_k^{-1}\chi_j:
 (a,b)\longmapsto
 \left(a,b\frac{p_j(a)}{p_k(a)}\right).
\]
These candidates preserve both readouts and satisfy the exact cocycle
$H_{\ell k}H_{kj}=H_{\ell j}$ wherever all maps are defined.

They are not an extraction of the actual native overlaps. The source
pieces are stated on the halo boundary; generic interior overlaps
require the previously identified collar extension. If references or
scales differ, the first coordinate instead satisfies
$a_k=(r_{c,j}+\varepsilon_{r,j}a_j-r_{c,k})/\varepsilon_{r,k}$,
and passive-coordinate transitions must also be supplied. Most
importantly, $b$ has an inherited cone-height meaning. The displayed
candidate is native only if that meaning transforms by this rule.
It cannot be justified by renaming a height coordinate after the fact.

\subsection{The earlier local score and flow do not glue naively}
Even under these candidate overlaps, the earlier Euclidean normal-form
score $\Psi_j=a^2+b^2+b^2f_j(a)^2$ is not invariant. For example
\[
 (a,b)=(1/2,1/2)\quad\hbox{in piece 1},\qquad
 H_{21}(a,b)=(1/2,9/19),
\]
gives
\[
 \Psi_1=129/256,\qquad \Psi_2\circ H_{21}=44569/92416.
\]
Both descriptions have the same radial and HC readouts. Therefore
their unequal scores cannot both be the same intrinsic C5 norm through
unchanged component identifications.

There is an independent update obstruction. For the earlier local
step $F_j(a,b)=(ta,tb)$, generally
\[
 H_{kj}F_j\ne F_kH_{kj}.
\]
At $a=b=t=1/2$, the target piece-2 height coordinates are respectively
$17/82$ and $9/38$, with difference $-23/779$.
This is a failure of the proposed interior gluing, not a counterexample
to a native overlap whose actual action has not been provided. It
shows why the valid fixed-piece estimates are insufficient for atlas
transfer.

\subsection{A common readout step supplies consistent chart dynamics}
Instead define the common step on $y=(a,v)$ by
$S_t(a,v)=(ta,tv)$, $0<t<1$. Its pullback is
\[
 F_j^{\rm cov}:=\chi_j^{-1}S_t\chi_j,
 \qquad
 F_j^{\rm cov}(a,b)
 =\left(ta,tb\frac{p_j(a)}{p_j(ta)}\right).
\]
Now $H_{kj}F_j^{\rm cov}=F_k^{\rm cov}H_{kj}$ follows by composition,
and $S_sS_t=S_{st}$. The radial and HC excess readouts both scale by
$t$, while the absolute radius remains at least $r_c>0$.
The intrinsic diagnostic
\[
 E_{\rm halo}(a,v)=a^2+v^2
\]
scales exactly by $t^2$. For $t=e^{-\tau}$ the pulled-back vector field is
\[
 \dot a=-a,\qquad
 \dot b=-b\left(1-\frac{a f_j'(a)}{p_j(a)}\right).
\]
This is different from the earlier chartwise $\dot b=-b$. On the three
symbolic pieces the bracket is respectively
\[
 \frac4{4+a},\qquad \frac{11}{11-3a},\qquad
 \frac{2+a^2}{2+a-a^2},
\]
and is positive on $0\le a\le1$.

One may retain the three-coordinate presentation
$R_j^{\rm NF}=(a,b,bf_j)$ with an explicit comparison metric. Set
\[
 B_j(a)=\begin{pmatrix}
 1&0&0\\0&1/p_j(a)&0\\0&f_j(a)/p_j(a)&1
 \end{pmatrix},\qquad
 G_j(a)=(B_j^{-1})^TB_j^{-1}.
\]
Then $R_j^{\rm NF}=B_j(a)(a,v,0)^T$ and
\[
 (R_j^{\rm NF})^TG_j(a)R_j^{\rm NF}=a^2+v^2.
\]
Residual-frame transitions $B_kB_j^{-1}$ preserve this comparison
through the corresponding metric transport. This repairs the
presentation-dependent diagnostic. It is \emph{not} permission to
silently mix native bounded/exterior/mixed channels or replace the
original C5 component norms. Such a change requires the source's
explicit comparison and component-typing data.

\subsection{A carrier that realizes the original halo formulas}
Use a declared dimensionless Euclidean active carrier $V=\mathbb R^4$,
with
\[
 J=\operatorname{diag}(1,1,-1,-1),\quad
 C=(I+J)/2,\quad\mathcal M=\operatorname{diag}(1,1,2,2).
\]
These matrices obey the inspected involution, HC idempotence and MM
intertwining identities, with scalar MM action on $\operatorname{im}C$.
Let $\Sigma_2=\left(\begin{smallmatrix}0&1\\1&0\end{smallmatrix}\right)$,
$U_1=\operatorname{diag}(\Sigma_2,\Sigma_2)$, and take the second
shift $U_\pi^{(2)}=U_1^2=I$. Supply the cone-to-radial projection
$\Pi_{\rm rad}=\operatorname{diag}(1,0,0,0)$.
Then the original tilt commutator is exactly zero. This is a concrete
finite two-cycle realization, not an identification of the source's
specific photon-frame shift or global $\mathcal M_{16}$.

Set $K=0$, $r_c=1/2$, $\varepsilon_r=1$,
$\varepsilon_h=1/2$, and the normalization floor $\delta_0=1$.
For $0\le a\le1/2$, $0\le v\le1$, define
\[
 r=\frac12+a,\qquad s=\frac v2,\qquad
 h=\sqrt{\frac{(1+s)r^2+s}{1-s}},\qquad
 x(a,v)=(r,h,0,0).
\]
The positive supplied first-shift collar radius is $r_c$; this
construction does not derive that radius from a physical first-pass law.
Direct evaluation of the source formulas near lines 44186--44211 gives
\[
 \rho_r(x)=\|\Pi_{\rm rad}U_\pi^{(2)}x\|=r,
 \qquad
 \rho_h(x)=\frac{h^2-r^2}{h^2+r^2+1}=s.
\]
The numerator is nonnegative, so this agrees with the absolute-value
definition. The tilt term vanishes for any nonnegative tilt weight.
The actual halo inequalities hold:
\[
 0<r_c\le\rho_r\le\varepsilon_r,
 \qquad0\le\rho_h\le\varepsilon_h,
 \qquad d(x,K)=\sqrt{r^2+h^2}\ge r_c.
\]
Thus the original operators and halo readouts, not just named output
variables, have been evaluated on a realized carrier.

On piece $j$, specifying the carrier height by the above formula with
$v=bp_j(a)$ gives
$\rho_h=\varepsilon_h b+\varepsilon_h bf_j(a)$ after evaluation.
Here the height parameter is supplied through $h_j(a,b)$ before the
residual is computed. The height-coordinate realization is regular:
\[
 \partial_b h_j
 =\frac{\varepsilon_h p_j(a)(2r^2+1)}
 {2h_j(1-s)^2}>0.
\]
Consequently $(a,b)$ are genuine independent local carrier coordinates
on the constructed collar interior. Matching them to the source's
\emph{inherited} height chart is still a separate identification; this
construction alone does not prove that earlier chart had this action.

\subsection{The fit-to-carrier bridge now has an explicit commuting diagram}
For the guarded lens sample of Section~\ref{sec:coherent-c5-coupled-guards},
put $\delta=A-5/4$, so $0<\delta\le1/12$. Define the synthesis map
\[
 \mathcal L(A)=x(\delta,\delta),\qquad
 \widehat F_t(A)=5/4+t(A-5/4).
\]
Its evaluated native readouts satisfy
\[
 a(\mathcal L(A))=v(\mathcal L(A))=\delta.
\]
Define $F_t^{\rm car}(x(a,v))=x(ta,tv)$ on the realized carrier surface.
Then the following identities are exact:
\[
 F_t^{\rm car}\mathcal L=\mathcal L\widehat F_t,\qquad
 \chi_j\,F_j^{\rm cov}=S_t\chi_j.
\]
Equivalently, using the source halo values $r=\rho_r$ and $s=\rho_h$,
the carrier step is
\[
 r'=r_c+t(r-r_c),\qquad s'=ts,\qquad
 h'=\sqrt{\frac{(1+s')(r')^2+s'}{1-s'}}.
\]
It preserves the constructed halo and the compensated sector at every
finite stage. The limit has $r=h=r_c>0$, not a through-kernel state.

The actual chart embeddings $X_j(a,b)=x(a,bp_j(a))$ also satisfy
\[
 X_k H_{kj}=X_j,\qquad
 F_t^{\rm car}X_j=X_jF_j^{\rm cov}.
\]
Thus the overlap and update diagrams commute on the realized carrier
itself, not only on assigned output coordinates.
For the dimensionless parameter $\tau$ with $t=e^{-\tau}$, its active
flow can be written directly in the carrier variables:
\[
 \dot r=-(r-r_c),\qquad \dot s=-s,
\]
\[
 \dot h=-\frac{r(1+s)(r-r_c)}{h(1-s)}
         -\frac{s(2r^2+1)}{2h(1-s)^2}.
\]
Differentiating the original expression
$s=(h^2-r^2)/(h^2+r^2+1)$ verifies $\dot s=-s$.
The denominators stay regular on the constructed halo because
$h\ge r_c>0$ and $s\le1/2$. This is an explicit internal relaxation
generator; no physical clock, viscosity or fluid conservation law has
been derived from it.

For the six signed lens residuals,
$z(A)=-2\delta\mathbf1_6$. They can now be read directly from the
realized native radial excess:
\[
 z(A)=-2\,\frac{\rho_r(\mathcal L(A))-r_c}{\varepsilon_r}\,
 \mathbf1_6.
\]
Hence any fixed typed C5 maps $\rho_{5,\alpha}=L_\alpha z$ on this
family factor through an explicitly evaluated halo channel and contract
by $t$. This is the promised readout-level link, conditional on those
maps being the original C5 maps. It is no longer necessary to assume an
unrealized pair of halo output numbers.

There is no unique off-family identification. For instance the maps
\[
 \rho_{5,\alpha}(a,v)
 =-2aL_\alpha\mathbf1_6+K_\alpha(v-a)
\]
agree on the fitted family $a=v$ for arbitrary fixed $K_\alpha$.
They are all homogeneous under $S_t$ but assign different off-family
component values. Their native channel meanings and zero sets cannot
be recovered from the one-dimensional score ledger alone.

\subsection{Guard transport can be included without changing the tilt}
In the same four-dimensional realization, let
\[
 P_{\partial,0}=\operatorname{diag}(1,1,1,0),\qquad
 s_{4,0}=d e_4,\quad 0<d\le1/256.
\]
Choose $T(A)=\operatorname{diag}(I_2,Q(A_0-A))$ on the complementary
two-plane, and co-transport $P_\partial$ and $s_4$ as in
Proposition~\ref{prop:covariant-c4-guard}. The positive C4 norm is
unchanged and its boundary projection remains zero. This $T$ commutes
with $\mathcal M,J,C,\Pi_{\rm rad}$ and $U_\pi^{(2)}$ and acts
trivially on $x(a,v)$. It therefore creates neither tilt nor a
noncommuting-MM perturbation. The C4 boundary projector is explicitly
different from the HC projector $C$.

Together with the previously evaluated relay coupling
$\kappa_{\rm rel}(A)=1/A$, boundary score and clock, this gives a
single declared packet with an evaluated halo carrier, shrinking fit
residual, positive unspent C4 remainder and passing numerical interface
guards. Neither the clock coupling nor the C5 masks have been extracted
from an original row-specific native realization.

\subsection{A passive chart change is not the carrier fit step}
For a passive isometric change of presentation, transport the state,
core, operators, metrics and scalar readout functions consistently.
Then $\Pi' U' x'=T\Pi Ux$, so its distance to the co-transported core
is unchanged. The cone scalars and tilt norm are likewise unchanged.
Such a presentation change cannot lower a score that factors only
through those same intrinsic halo values.

The present $F_t^{\rm car}$ instead changes $r$ and $h$ within the
compensated sector. It is an active, newly declared fit-coupled carrier
map, whereas $T(A)$ only co-transports the guarded complementary data.
In fact the supplied MM matrix obeys
$\mathcal MCx=x$ on this sector, and cannot by itself generate
$F_t^{\rm car}(x)\ne x$. Therefore this explicit bridge still does
not identify the fit step with the original MM action. A claimed native
derivation must supply that active update, not substitute passive
covariance for dissipation.

\subsection{Why scalar agreement cannot identify the old component maps}
There is a smooth ambiguity even after the fit-to-carrier bridge is
fixed. On the same positive fit amplitude $\delta$, consider
\[
 R_5^{(1)}(\delta)=(\delta,0,0),\qquad
 R_5^{(2)}(\delta)=(\delta\cos\delta,\delta\sin\delta,0).
\]
Both are smooth at zero and have squared norm $\delta^2$. Both therefore
give the same exact geometric score ledger under $\delta'=t\delta$.
But $R_5^{(2)}(t\delta)\ne tR_5^{(2)}(\delta)$ for generic input.
Fixed linear factorization is consequently sufficient for our transfer
proof, not a consequence of the ledger and not necessary for that
scalar recurrence. These are alternative mathematical component maps,
not assignments of their missing physical channel meanings.

Thus even an unlimited continuation of the existing score values cannot
select the original component map. The required additional information
is a component evaluation or defining rule tied to an actual witness
and chart, rather than another scalar recurrence value.

\paragraph{Bridge status.}
The lens fit, evaluated original halo formulas, compatible collar
overlaps and guarded carrier map now form an explicit commuting local
construction. The earlier naive atlas gluing fails and has been repaired
for this readout-based realization. What is proved is existence of a
consistent local completion of the inspected formulas. What remains
unproved is historical identity: the source's actual C5 component maps,
inherited cone-height transitions and native MM/HC update must agree
with the supplied completion. Its active relaxation is not yet a
physical Navier--Stokes law or a global OP5 closure theorem.

\paragraph{Verification.}
All three candidate overlap cocycles, corrected update intertwining and
comparison-metric identities were checked symbolically. The naive
score discrepancy and update seam were checked as exact rational values.
The carrier matrices satisfy the displayed MM/HC identities, their tilt
commutator vanishes, and substitution of the explicit height gives the
original normalized cone residual exactly. Guard positivity and the
fit-to-carrier commuting diagram follow from the displayed formulas.
The carrier height derivatives and the active-flow derivative of the
original cone residual were checked symbolically as well.

\section{A guarded Hilbert-shift memory lift for repeated C5 lowering}
\label{sec:hilbert-shift-memory-lift}
The preceding active fit step now receives a conservative enlarged-space
realization: the visible residual shrinks, while every emitted part is
retained in a separate history channel. This supplies an explicit
repetition rule missing from a reused front/back rotation. It is a new
operator completion compatible with stated HS gates when its embeddings
are licensed, not a reconstruction of the source's unspecified native
operator or a derivation of the numerical rate.

\subsection{Three shift meanings and the source HS(5) interface}
The QAM namespace shift $k\mapsto k+1$ leaves payloads unchanged and
reserves tag zero. It does not move amplitudes or produce decay. The
source's cyclic regime shift is a separate carrier operator. A third
construction is its HS tower, with carrier sizes $C_\ell=2^\ell$ and
Boolean mask counts $2^{C_\ell}$. Here $\ell$ denotes an architectural
level; $n$ below counts fit updates. These indices are not identified.

The source near lines 45774--45806 explicitly supplies
\[
 \mathcal H^{(5)}
 =\mathcal H^{(4)}_{\rm front}\oplus\mathcal H^{(4)}_{\rm back},
 \qquad Q_5=Q_{\rm front}\oplus Q_{\rm back},
\]
and the cross-side debt operator
\[
 \mathcal D_5(U)=(I-Q_5)UQ_5+Q_5U(I-Q_5).
\]
It licenses zero-debt front/back comparisons only through declared
selectors or comparison maps. The construction below provides such
maps as additional data. $Q$ is the meta-septum projector here, not
automatically the earlier HC projector $C$.

\subsection{A one-step coupling that leaves protected directions alone}
Let $\mathcal F$ be the fixed fit Hilbert space. If the three native
component spaces have been identified, it can be their orthogonal direct
sum with the declared fit norms. Three named components need not mean
three scalar dimensions. Let $0<t<1$, $s=\sqrt{1-t^2}$ and $q=t^2$.

Assume declared isometric embeddings
$\iota_f:\mathcal F\to\operatorname{im}Q_{\rm front}$ and
$\iota_b:\mathcal F\to\operatorname{im}Q_{\rm back}$ into orthogonal
Hilbert realizations of the source gates, with
$Q_{\rm front}=Q_{\rm front}^*$ and $Q_{\rm back}=Q_{\rm back}^*$.
Self-adjointness here is additional metric data, not a consequence of
the source's bare idempotence statement. Set
\[
 S_f=\iota_f\iota_f^*,\quad S_b=\iota_b\iota_b^*,
 \qquad W=\iota_b\iota_f^*.
\]
Then $W^*W=S_f$, $WW^*=S_b$. The explicit front/back coupling is
\[
 U_t=
 \begin{pmatrix}
 I+(t-1)S_f&-sW^*\\
 sW&I+(t-1)S_b
 \end{pmatrix}.
\]
On the two fit copies it is the rotation
\[
 (\xi,\eta)\longmapsto(t\xi-s\eta,s\xi+t\eta),
\]
and on the remaining orthogonal directions it is the identity.

\begin{proposition}[HS-gated unitary fit coupling]
\label{prop:hs-gated-unitary-coupling}
Under the stated embeddings, $U_t$ is unitary and
$\mathcal D_5(U_t)=0$. Any protected C4 vector orthogonal to both fit
copies is unchanged, including its positive norm.
\end{proposition}
\begin{proof}
The fit-copy rotation is unitary because $t^2+s^2=1$; its orthogonal
complement is fixed. The fit embeddings lie in the selected sectors,
so $Q_{\rm back}W=W=WQ_{\rm front}$ and the diagonal fit projections
commute with their respective gates. Consequently $[U_t,Q_5]=0$,
and both off-sector blocks vanish. The identity action proves the last
statement.
\end{proof}
Thus front-to-back transport within selected sectors is not by itself
cross-side leakage. This does not verify the original embedding data.
Compatibility with a separate HC projector or an existing MM converter
requires their own intertwinings on the selected fit copies. No equality
$Q=C$ or identity $U_t=\mathcal M$ is inferred.

With an empty incoming excitation amplitude $\eta=0$, the visible fit
vector becomes $t\xi$ and the outgoing vector is $s\xi$. Reusing that
same occupied slot gives a rotation through $n\theta$, where
$t=\cos\theta$. At $\theta=\pi/6$, the first three visible squared
norm ratios are $3/4,1/4,0$, not $(3/4)^n$. An outgoing history rule
is therefore indispensable for the intended recurrence.

\subsection{The explicit outgoing-history shift}
Let
\[
 \mathcal K_+=\mathcal F\oplus\ell^2(\mathbb N;\mathcal F),
\]
with visible coordinate $\xi$ and outgoing coordinates
$\eta_1,\eta_2,\ldots$, ordered newest first. Define
\[
 V_t(\xi,\eta_1,\eta_2,\ldots)
 =(t\xi,s\xi,\eta_1,\eta_2,\ldots).
\]
This is a weighted injection together with a history shift, not the
QAM tag shift. Existing outgoing coordinates are moved to older slots;
none is fed back into the active fit coordinate.

\begin{theorem}[Norm-preserving repeated residual lowering]
\label{thm:shift-memory-lowering}
$V_t$ is an isometry. For $X_0=(\xi_0,0,0,\ldots)$,
\[
 V_t^nX_0=
 (t^n\xi_0,st^{n-1}\xi_0,st^{n-2}\xi_0,\ldots,s\xi_0,0,\ldots).
\]
If $\xi_0\ne0$, the normalized visible fit score is $q^n>0$ for
every finite $n$, and the outgoing-history squared norm is
$(1-q^n)\|\xi_0\|^2$. The total squared norm is unchanged.
\end{theorem}
\begin{proof}
Orthogonality of the history copies gives
\[
 \|V_tX\|^2=(t^2+s^2)\|\xi\|^2+\sum_j\|\eta_j\|^2=\|X\|^2.
\]
Induction gives the displayed orbit. Its outgoing norm is
\[
 s^2\sum_{j=0}^{n-1}t^{2j}\|\xi_0\|^2
 =(1-q^n)\|\xi_0\|^2.
\]
Adding the visible norm $q^n\|\xi_0\|^2$ proves conservation.
\end{proof}
$V_t$ is not onto and must not be called a unitary on $\mathcal K_+$.
Norm preservation here concerns the chosen Hilbert norm; identifying it
with physical energy requires an independent observable and units.

Identification with the original C5 score also requires its three
component readouts to select the visible fit coordinate, rather than
the full memory. In a declared direct-sum fit realization this means
$\rho_{5,\alpha}(\mathsf W(X))=\pi_\alpha P_{\rm vis}X$, where
$\pi_\alpha$ is the corresponding component projection and
$\mathsf W$ realizes the native witness. If the old score instead
counts the whole fit-plus-history norm, that score remains constant
and the native lowering claim does not follow. The source's actual
readout selector therefore remains part of the required embedding.

The zero incoming amplitude denotes absence of an excitation in this
direct-sum description. A physical fresh quantum ancilla instead has
a normalized reference state, not the zero vector. The construction can
be interpreted in a suitable excitation sector of a chain model, but
that interpretation and the native C5 observable must be specified.
Likewise a QAM inactive tag is not itself a physical vacuum state.

\subsection{An explicit unitary completion with incoming and outgoing lanes}
To retain reversible full transport, introduce
\[
 \mathcal K=
 \mathcal F\oplus\ell^2(\mathbb N;\mathcal F)_{\rm in}
 \oplus\ell^2(\mathbb N;\mathcal F)_{\rm out}.
\]
Define the autonomous map
\[
 \begin{aligned}
 \mathcal U_t(\xi;\alpha_1,\alpha_2,\ldots;\eta_1,\eta_2,\ldots)
 ={}&(t\xi-s\alpha_1;\alpha_2,\alpha_3,\ldots;\\
     &s\xi+t\alpha_1,\eta_1,\eta_2,\ldots).
 \end{aligned}
\]
The local pair is rotated while the incoming lane advances and the
outgoing lane retains its emitted state. It is unitary: it preserves
norms, and its inverse reconstructs
\[
 \xi=t\xi'+s\eta_1',\qquad
 \alpha_1=-s\xi'+t\eta_1',
\]
prepends $\alpha_1$ to the new incoming tail and removes the first
outgoing slot. Thus no stored amplitude is erased.

For the embedding $E\xi=(\xi;0,0,\ldots;0,0,\ldots)$ and visible
projection $P_{\rm vis}$,
\[
 P_{\rm vis}\mathcal U_t^n E=t^nI_{\mathcal F}
 \qquad(n\ge0).
\]
Zero incoming excitations stay zero, so the previous isometric orbit
is the restriction of this unitary evolution. This is an explicit
power-dilation realization of the contraction $T=tI$, not merely a
one-step front/back dilation. Sz.-Nagy dilation theory supplies the
general operator-theoretic setting; see the definition and introductory
theorem discussion in Pal and Sahasrabuddhe,
\href{https://arxiv.org/html/2205.09093}{\emph{Minimal unitary dilations for commuting contractions}}.

\subsection{Finite registered histories and the actual capacity requirement}
For a declared horizon $N_{\rm hist}$, one can work entirely in
\[
 \mathcal K_N=\mathcal F_{\rm vis}
 \oplus\bigoplus_{j=1}^{N_{\rm hist}}\mathcal F_j.
\]
At step $j$, rotate only the visible fit coordinate and the still-empty
slot $j$ by the same two-copy unitary. Leave earlier slots unchanged.
The chronological stored vectors are then
\[
 \xi_n=t^n\xi_0,\qquad
 \eta_j=st^{j-1}\xi_0\quad(1\le j\le n\le N_{\rm hist}).
\]
Every step is unitary on this finite space, but the active slot changes
with the step index. No reset of an occupied slot is used. Continuing
beyond $N_{\rm hist}$ requires further licensed slots or a separately
specified reservoir operation.

If $d=\dim_{\mathbb R}\mathcal F$, this generic orthogonal-copy
realization uses $(N_{\rm hist}+1)d$ real coordinates before independent
protected guard coordinates are counted. At $q=3/4$ and visible
threshold $1/256$, the first detection is at $n=20$, so its fit/history
budget is $21d$. A three-scalar illustration uses $63$ coordinates;
a known coherent one-mode family can use $21$. The three native norm
slots do not establish $d=3$ or a one-mode realization.

The source's HS(6) carrier has $64$ real channels, not $2^{64}$ dimensions.
This raw count may accommodate particular finite constructions, but an
actual gate-compatible embedding needs enough \emph{licensed} fit
dimensions and room for protected data. A Boolean mask count does not
provide that capacity. The HS(0)--HS(6) declarations alone therefore do
not prove that this twenty-step memory has been realized in the native
Dixon support class.

More generally, a fixed unitary on one finite-dimensional carrier with
a fixed visible projection cannot realize nonzero strictly geometric
decay for all $n$. Compactness gives arbitrarily large powers tending
to the identity: choose increasingly distant pairs of powers in the
same increasingly small compactness ball. Their difference powers tend
to $I$. The visible norm then returns arbitrarily close to its positive
initial value, contradicting decay to zero. Thus an unlimited autonomous
unitary realization needs an enlarged memory, whereas a specified finite
prefix can use the finite chronological construction above.

\subsection{C4 separation and the existing fit-to-carrier map}
Extend the full memory evolution by the identity on a protected guard
space $\mathcal G_4$. Keep its actual positive C4 vector, projector and
ledger bound outside the fit/history coordinates. If initially
$P_\partial s_4=0$, this direct-sum extension preserves both the
projection and the full C4 norm exactly. This is an operator statement
on declared data, not a claim that the original projector has no
undisclosed fit/history dependence.

With licensed copies of the selected fit sector, the extended selection
projector acts copywise and commutes with the shift and local rotations.
Consequently there is no transport into its complementary sector.
All outgoing slots must nevertheless appear in the QAM lineage and
comparison record; an unregistered emitted fit contribution must not be renamed
as debt-free absorption or hidden in the C4 remainder. The original
guard tuple still requires identification of these declared sectors
with the source's actual leakage and absorption tests.

On the previously constructed coherent fit family, write
$R_5(A)=\kappa\delta u$ with $\delta=A-5/4>0$, fixed unit vector $u$
and fixed $\kappa>0$. The visible coordinate obeys
$\xi_n=\kappa t^n\delta_{\rm in}u$, where $\delta_{\rm in}$ is the
initial positive fit amplitude. Decode
\[
 \delta_n=\langle u,\xi_n\rangle/\kappa,
 \qquad A_n=5/4+\delta_n.
\]
The decoded update is exactly
$A_{n+1}=5/4+t(A_n-5/4)$, so its synthesis map to the realized halo
carrier commutes with the visible memory evolution. The corrected
chart maps $F_j^{\rm cov}$ therefore remain compatible. The previously
checked core reserve, relay and clock inequalities remain valid on
that declared family, while C4 remains unspent.

The visible score tends to zero, not the total stored state. In the
chronological infinite register, the states converge strongly to
\[
 (0;s\xi_0,st\xi_0,st^2\xi_0,\ldots),
\]
whose norm equals $\|\xi_0\|$. Indeed the squared difference from the
state after $n$ steps is $2q^n\|\xi_0\|^2$. In the autonomous
newest-first frame, $V_t^nX_0$ instead tends weakly to zero but not
strongly: every fixed coordinate tends to zero, while the total norm
is constant. The moving history labels explain the different limits.
A vanishing visible readout is consequently not a full-state zero,
and compact finite-witness limit arguments cannot be applied to the
enlarged history without checking its topology.

\subsection{What fixes the rate, and what the literature establishes}
The repeatability mechanism works for every $0<q<1$. At $q=3/4$,
$t=\sqrt3/2$, $s=1/2$, and the principal coupling angle is $\pi/6$.
Nothing in tag reindexing, carrier doubling or the memory shift selects
that angle. In the dimensionless readout flow $t=e^{-\Delta\tau}$,
the corresponding sampling interval is
$\Delta\tau=\tfrac12\log(4/3)$. This belongs to the new common-readout
flow, not the earlier fixed-piece upper-bound sampling convention.
A physical coupling strength and clock remain separate obligations.

Attal and Pautrat,
\href{https://arxiv.org/abs/math-ph/0311002}{\emph{From repeated to continuous quantum interactions}},
Sections II.1--II.3, explicitly use a sequence of new environment copies;
an already-used copy does not interact with the small system again.
Their full dynamics is unitary and yields reduced discrete evolutions.
Ziman and Buzek,
\href{https://arxiv.org/abs/1006.2794}{\emph{Open system dynamics of simple collision models}},
Sections 2.2 and 4.1, give partial-swap collisions and a surviving
system-excitation amplitude proportional to $\cos^n\eta$.
The corresponding excitation weight is $\cos^{2n}\eta$.
That is not a universal formula for every density-matrix residual norm,
nor an identification of C5 with a quantum excitation observable.

These sources support the storage and enlarged-unitary mechanism.
The explicit shift, HS-gate checks, finite capacity test and
fit-to-carrier decoding above are calculations for this note. Their
native application still needs the actual C5 readout and selector
embeddings. No Navier--Stokes law, viscosity or global all-readout
closure follows from the Hilbert-space construction alone.

\paragraph{Memory-lift result and verification.}
A specified C5 contraction now has a norm-preserving history lift,
an explicit autonomous unitary completion and a finite-horizon unitary
implementation. Fresh outgoing slots remove the previously demonstrated
backflow obstruction. Symbolic matrix multiplication checked unitarity,
zero HS(5) cross-side debt and unchanged protected directions on a
partial-fit realization. Twenty fresh-slot steps were checked for exact
geometric visible score, retained history and constant total norm; the
two-sided inverse was checked algebraically. The detector first passes
at step 20 for the stated normalized initial score.
What remains is the native fit/guard embedding and independent rate
selection, not a further assumption that the old occupied back slot
somehow resets itself.

\section{A Riemann-spectral test for the visible C5 contraction}
\label{sec:riemann-spectral-c5-test}

\paragraph{Scope and source-data audit.}
The source's Riemann/GUE block uses a claimed self-adjoint realization
$L_0$ with zero ordinates as spectral data (around lines 7220--7283).
Its r13 handoff block declares a finite clock $\{\tau_j\}$, a comparison
map to zero ordinates and a circular phase-mismatch score (around
lines 54460--54525). It explicitly calls this a diagnostic probe, not
an identification. Neither block supplies the spectral coordinates of
the three r04 C5 residual functions or a numerical handoff interval.
The r04 score and guard tuple around lines 52300--52365 remain separate
objects. Thus no unique native matrix can yet be evaluated from these
declarations. The construction below makes the missing data explicit
and tests one specified finite alternative without fitting its clock
to the desired ratio.

The spectral realization of $L_0$ is a premise here, not a theorem
proved by this working note. Finite computations with zero ordinates
do not establish that realization or the Riemann hypothesis. In
particular, the factor $i$ in $e^{-i\tau L_0}$ is essential: for a
self-adjoint $L_0$ this is unitary evolution, not the dissipative
semigroup $e^{-\tau L_0}$. The source's use of heat-kernel language
does not justify exchanging those two operators.

\subsection{The exact visible/history norm identity}
Let $\mathcal F$ be the declared fit Hilbert space with the actual
sum norm used by C5. Let $\mathcal H$ be a larger Hilbert space and
$E:\mathcal F\to\mathcal H$ an isometry. Set
\[
 P=EE^*,\qquad U_h=e^{-ihL_0},\qquad
 A_h=E^*U_hE,\qquad B_h=(I-P)U_hE.
\]
Here $h$ is a specified sampling interval; it is not the cone height.
The adjoints refer to the declared positive metrics. If the actual
score in raw coordinates $r$ is $r^\dagger G_5r$, with $G_5>0$,
use $\xi=G_5^{1/2}r$. In orthonormal ambient coordinates the raw
embedding matrix must satisfy $E_0^\dagger E_0=G_5$; then
$E=E_0G_5^{-1/2}$ satisfies $E^*E=I$. Here $\dagger$ denotes the
ordinary coordinate conjugate transpose. This metric identification
cannot be replaced by an arbitrary Euclidean norm on unnamed native slots.

\begin{proposition}[Quarter-loss criterion]
\label{prop:riemann-quarter-loss-criterion}
For the preceding orthonormal fit coordinates,
\[
 A_h^*A_h+B_h^*B_h=I_{\mathcal F}.
\]
Consequently, the uniform visible bound
\[
 \|A_h\xi\|^2\le\tfrac34\|\xi\|^2
 \quad(\xi\in\mathcal F)
\]
is equivalent to $B_h^*B_h\ge\tfrac14I_{\mathcal F}$.
Exact equality for every fit vector instead requires
$B_h^*B_h=\tfrac14I_{\mathcal F}$. On a restricted native residual
family it is enough to check the corresponding quadratic inequality
on that family; it need not hold on all of $\mathcal F$.
\end{proposition}
\begin{proof}
The summands $PU_hE\xi$ and $(I-P)U_hE\xi$ are orthogonal and
$U_hE$ is an isometry. Moreover $PU_hE=EA_h$. The Pythagorean
identity therefore gives the operator identity and both equivalences.
\end{proof}

The range of $B_h$ is the part invisible to this particular fit
selector. It is not automatically the source's forbidden new-leakage
channel. To keep it licensed, one sufficient declared condition is an
orthogonal gate $Q$ with $QE=E$ and $QU_h=U_hQ$; then
$QB_h=B_h$ because $P\le Q$. Both visible and stored contributions
stay inside $Q$. Identifying $Q$ with an actual HS/QAM gate and
checking the native C4 absorption and C6 scope monitors remain necessary.

If instead the visible space itself reduces $U_h$, so that
$PU_h=U_hP$, then $B_h=0$ and $A_h$ is unitary: the score ratio is
$1$, not $3/4$. Thus an invariant HC or cascade sector, when its
selector commutes with the spectral transport, cannot alone supply
the desired loss. A finer fit/history split inside a licensed sector
is required. The fit projector $P$, HC projector $C$ and HS gate $Q$
are not silently identified.

\subsection{A finite matrix that the missing native data would determine}
Assume a declared finite invariant spectral carrier with orthonormal
modes $e_k$, real frequencies $\omega_k$ and an $m\times d$ fit
embedding matrix $V$ satisfying $V^*V=I_d$. Then
\[
 D_h=\operatorname{diag}(e^{-i\omega_1h},\ldots,e^{-i\omega_mh}),
 \qquad A_h=V^*D_hV,
\]
\[
 K_h:=B_h^*B_h=I_d-A_h^*A_h.
\]
This gives an executable decision test: calculate
$\lambda_{\min}(K_h)$ and compare it with $1/4$, or equivalently
calculate $\lambda_{\max}(A_h^*A_h)$ and compare it with $3/4$.
Three component names or a finite zero list do not determine $V$.
The source's circular mismatch score likewise does not determine $h$
or the products $(\omega_j-\omega_k)h$ without an explicit timing law.

Finite truncation is not automatically valid for a native fit supported
on further modes. If a computed finite compression $\widetilde A_h$
approximates the actual $A_h$ with operator error at most $\varepsilon$,
then
\[
 \|A_h\|^2\le(\|\widetilde A_h\|+\varepsilon)^2.
\]
A finite certificate therefore needs enough margin to absorb a proved
tail/embedding error. No such native error bound is supplied by the
three residual-slot declarations.

\subsection{Two spectral modes: the angle is calculated, not postulated}
Choose two orthonormal spectral modes and, for $0\le w\le1$, define
the visible direction and its orthogonal history partner by
\[
 f=\sqrt w\,e_a+\sqrt{1-w}\,e_b,\qquad
 g=\sqrt{1-w}\,e_a-\sqrt w\,e_b.
\]
Writing $\varphi=(\omega_b-\omega_a)h$, the visible amplitude is
\[
 a_h=w e^{-i\omega_a h}+(1-w)e^{-i\omega_b h},
\]
and direct multiplication gives
\[
 q(w,\varphi)=|a_h|^2
 =1-4w(1-w)\sin^2(\varphi/2).
\]
Thus the quarter-loss condition is exactly
\[
 4w(1-w)\sin^2(\varphi/2)\ge\tfrac14.
\]
No universal ratio follows from the frequency gap alone: the visible
weights matter. A single eigenmode ($w=0$ or $1$) has ratio $1$
regardless of the same clock and spectrum.

For equal weights and $\bar\omega=(\omega_a+\omega_b)/2$,
the full matrix in the $(f,g)$ basis is
\[
 U_h=e^{-i\bar\omega h}
 \begin{pmatrix}
 \cos(\varphi/2)&i\sin(\varphi/2)\\
 i\sin(\varphi/2)&\cos(\varphi/2)
 \end{pmatrix},\qquad q=\cos^2(\varphi/2).
\]
The $60$-degree relative phase gives $q=3/4$ and a $30$-degree
visible/history mixing angle. A $120$-degree relative phase gives
$q=1/4$, so the source's triolic phase label is not itself a
$3/4$ derivation. For a positive gap, the first-period interval
that yields $q\le3/4$ is
\[
 \frac{\pi}{3(\omega_b-\omega_a)}
 \le h\le
 \frac{5\pi}{3(\omega_b-\omega_a)}.
\]
Other intervals repeat periodically. These are constraints to test
against an independent clock, not permission to define the clock
by the desired contraction.

\subsection{A fixed-clock experiment using six Riemann-zero ordinates}
For a diagnostic only, use the first six positive critical-line zero
ordinates computed numerically and the source's rounded frame input
$\gamma_L=1.864$. Declare
\[
 \omega_k=\gamma_k/\gamma_L,\qquad
 h_{\rm probe}=\gamma_L^{-3/2}
 \simeq0.392944430867119.
\]
The interval borrows the already present routed cutoff of the older
Theta-balance defect (source around line 26241). It is \emph{not}
identified with an r13 handoff interval or a physical time. In
particular, this reuse of a cutoff is an extra model choice, not a
deduction from the native clock. The frequencies and interval are
fixed before testing $3/4$.

Use three orthogonal equal-weight pairs $(1,2)$, $(3,4)$ and $(5,6)$,
giving three complex scalar fit directions in a six-mode carrier.
This is an illustrative component realization, not a claim about
the native component dimensions or their spectral supports. The result is
\[
\begin{array}{c|c|c|c|c}
\text{pair}&\omega_b-\omega_a&\varphi&1-q&q\\
\hline
(1,2)&3.694911211&1.451894783&0.440689211&0.559310789\\
(3,4)&2.904516387&1.141313539&0.291799691&0.708200309\\
(5,6)&2.495234212&0.980488387&0.221691562&0.778308438
\end{array}
\]
Here $K_h$ is diagonal with the three displayed loss fractions.
Its smallest eigenvalue is about $0.221691562<1/4$; the uniform
$3/4$ certificate therefore fails. The same data give a genuine
uniform contraction bound $q_{\max}\simeq0.778308438<1$ for this
declared carrier, but not the prescribed bound $3/4$.

The first pair alone has quarter-loss margin about $0.190689211$.
It still satisfies $q\le3/4$ for unequal visible weights in the
approximate interval
\[
 0.1710977413\le w\le0.8289022587.
\]
Thus exact equal weighting is unnecessary for that one-mode upper
bound. The last pair, by contrast, has maximal possible loss
$\sin^2(\varphi/2)\simeq0.221691562$; changing its two-mode weights
cannot repair the $3/4$ bound at the same clock.

For the three equal-weight pairs, the common first-period clock window is
\[
 0.4196790611\lesssim h\lesssim1.4170808058.
\]
An actual independently determined handoff interval in this window
would pass this particular three-pair test. Moving $h_{\rm probe}$
into the window merely to obtain a pass would be calibration,
not an explanation of the original rate.

The phase is invariant under a consistent change of frequency/time
units: $\omega_k=\gamma_k/\gamma_L$ requires
$h_{\rm photon}=\gamma_L h_{\rm matter}$ for the same experiment.
Changing only the frequency scale cannot create a physical loss.
The borrowed cutoff law above specifies a distinct dimensionless
experiment; it does not follow from that units transformation.

\subsection{Fresh history, direction-dependent rates and the scalar bridge}
Apply each two-mode unitary to the visible direction and a fresh
zero-excitation history slot, leaving old slots unchanged. This is
the chronological memory construction of
Section~\ref{sec:hilbert-shift-memory-lift}, with complex phases retained.
It uses successive registered copies of the chosen spectral pairs,
not powers of one fixed unitary on the original six-mode carrier.
Embedding those copies and their switching rule into the original
generator is an additional native obligation. For a horizon $N$,
three unrestricted complex fit coordinates require $6(N+1)$ real
fit/history coordinates in the generic orthogonal-copy construction,
before protected guard coordinates. The six-mode one-step experiment
does not establish this repeated-history capacity in HS(6).
For a diagonal visible compression with amplitude entries $a_\alpha$,
\[
 \xi_{n,\alpha}=a_\alpha^n\xi_{0,\alpha},\qquad
 \Phi_n=\sum_\alpha q_\alpha^n|\xi_{0,\alpha}|^2.
\]
The emitted squared norm at step $j$ in channel $\alpha$ is
$(1-q_\alpha)q_\alpha^{j-1}|\xi_{0,\alpha}|^2$; summing visible
and all registered histories preserves the initial full norm.
Consequently $\Phi_{n+1}\le q_{\max}\Phi_n$ on this model.
Reusing one occupied partner would instead return the oscillatory
ratio $\cos^2(n\varphi/2)$ for an equal-weight pair, not $q^n$.

For the three-pair probe, put $p_\alpha=|\xi_\alpha|^2$.
The state-dependent $3/4$ test is approximately
\[
 0.028308438\,p_3
 \le0.190689211\,p_1+0.041799691\,p_2.
\]
This allows many initial mixtures but is not forward invariant when
$p_3>0$: the slowest pair eventually dominates the score ratio.
For equal initial powers $p_\alpha=1/3$, the first ratio is about
$0.681939846$. Nevertheless the first failure occurs at the update
$n=9$ to $10$, with ratio about $0.750465631$; the preceding ratio
is about $0.746822300$. A successful initial score is therefore not
a uniform or repeatable $3/4$ certificate.

On the independently justified invariant subspace $\xi_3=0$, the
bound $3/4$ does hold. Removing the third component merely to get
this result would not be a valid full-C5 proof. For a one-mode native
fit independently supported on the first pair, the stronger rate
$q\simeq0.559310789$ is a candidate rather than an imposed $3/4$.

On that first-pair model, removing the known scalar phase by a
declared norm-preserving readout co-transport leaves positive
amplitude $t=\cos(\varphi/2)$. The existing constructed lens family
may then use
\[
 A_{n+1}=5/4+t(A_n-5/4).
\]
For $5/4<A_0\le4/3$, this remains in the already checked collar and
relay interval, and the fit-to-carrier synthesis still commutes with
this scalar update. This is a conditional spectral-rate realization
of that declared family. The native spectral embedding, phase
co-transport, licensed history capacity and actual guard evaluations
are not supplied by the scalar calculation.

The original exact-ratio monitor would reject these different rates.
The earlier explicitly declared inequality monitor accepts the first
pair but not every direction of the three-pair example. At normalized
initial score $1$ and detector $1/256$, the equal-power three-pair
memory orbit first passes at step $19$ (scores about $0.004340056$
at step $18$ and $0.003328877$ at step $19$). The separate single-pair
orbits first pass at steps $10$, $17$ and $23$. These are diagnostic
results for the specified model, not revisions of the native
twenty-step ledger.

\subsection{Why this does not yet derive continuum damping}
For a bounded finite spectral generator,
\[
 B_h=-ih(I-P)L_0E+O(h^2),\qquad
 K_h=h^2E^*L_0(I-P)L_0E+O(h^3).
\]
The per-step loss starts quadratically, not linearly, in $h$.
For one equal-weight pair of fixed gap $\Delta\omega$ and fresh
histories, holding total time $\tau$ fixed gives
\[
 \lim_{n\to\infty}
 \cos^{2n}\!\left(\frac{\Delta\omega\,\tau}{2n}\right)=1.
\]
Indeed the logarithm is
$-(\Delta\omega)^2\tau^2/(4n)+O(n^{-3})$. Thus arbitrarily frequent
sampling with unchanged finite coupling does not produce a nonzero
continuous relaxation rate. A nontrivial dissipative scaling would
need independently specified reservoir/coupling data; it is not
obtained by renaming spectral interference as viscosity.

\paragraph{Verification and native continuation.}
Symbolic two-mode calculations checked unitarity, the orthogonal
visible/history identity and the unequal-weight loss formula.
The first six zero ordinates were computed with
\texttt{mpmath.zetazero} at $40$ and $80$ decimal digits; their
maximum disagreement was about $3.7\times10^{-40}$. The matrix
eigenvalues, first-period clock window, fresh-memory score ratios
and detector indices were evaluated from those ordinates.
The full six-mode matrices and forty fresh-history steps were checked
at $65$ decimal digits, including the identity after a non-Euclidean
fit-coordinate change. Compression and total-norm discrepancies were
below $5\times10^{-66}$; the evaluated zeta residuals at the six
computed ordinates were below $6\times10^{-65}$.
Repeating the probe at rounded frame inputs $1.8635$ and $1.8645$
left the last-pair ratio above $3/4$ (about $0.778035$ and
$0.778581$). These are reproducible floating-point diagnostics,
not an interval-arithmetic certificate or a native-spectrum proof.

For a native result, supply the actual spectral fit embedding and
metric, an independent handoff timing law, and a licensed history
and guard embedding. Prove that the source residuals equal the
visible coordinates and that the native witness update intertwines
with their compression. Then apply
Proposition~\ref{prop:riemann-quarter-loss-criterion} with any necessary
tail bound. The Riemann part now gives an explicit operator test and
a non-calibrated finite example, but does not select a universal
$3/4$ or prove the separate routed Theta--R balance.

\section{A metric-mediated information-to-fit response candidate}
\label{sec:metric-mediated-gravity-candidate}

\paragraph{Construction status.}
This section constructs a conditional gravitational response interface
and an explicit finite norm-preserving coupling model. It does not
derive a new law of gravity, establish physical tachyons or identify
the three C5 components with three physical sectors. The label $T$
denotes the source's information-facing channel. A scalar phase factor,
a frame relabeling and a physical metric perturbation must remain
different objects.

The source has three relevant ingredients. The symmetric-frame block
places matter and tachyon readouts at opposite offsets from the photon
reference (around line 6410). The Einstein-reconstruction block proposes
the coframe and metric
\[
 E^a{}_{\mu}=\operatorname{Re}
 (\bar\Psi_\triangle\Gamma^a\mathcal D_\mu\Psi_\triangle),
 \qquad g_{\mu\nu}=\eta_{ab}E^a{}_{\mu}E^b{}_{\nu}
\]
under exact projected closure, nondegeneracy and coarse-graining
assumptions (around lines 8270--8395). These assumptions do not follow
from detector-level C5 terminality. The r82 channel supplies a typed
projection $\pi_{T\to M}$ through which information can appear as
matter-frame residual or phase data (around lines 62879--62918).
Neither a source stress tensor nor a numerical gravitational transfer
kernel is evaluated there. The following weak-field realization is
therefore an added effective-model assumption, not a promotion of
those declarations to a completed native Einstein theorem.

\subsection{A genuinely gravitational first leg: the metric response}
Take a weak-field, asymptotically flat test-field regime with
signature $(-,+,+,+)$ and
\[
 g_{\mu\nu}=\eta_{\mu\nu}+p_{\mu\nu}.
\]
Let $z_T$ be a declared input perturbation, and let
\[
 \delta T_{\mu\nu}=\mathcal S_T z_T
\]
be its independently specified conserved stress-energy source,
$\partial^\mu\delta T_{\mu\nu}=0$. If the $T$ channel exchanges
stress-energy with other fields, their contributions must be included
in the conserved total; a nonconserved isolated source cannot simply
be inserted into the gravitational equation.

In Lorenz gauge and with zero incoming perturbation, set
\[
 \bar p_{\mu\nu}=p_{\mu\nu}
 -\tfrac12\eta_{\mu\nu}\eta^{\alpha\beta}p_{\alpha\beta},
 \qquad\partial^\mu\bar p_{\mu\nu}=0,
\]
\[
 \Box\bar p_{\mu\nu}
 =-\frac{16\pi G_N}{c^4}\delta T_{\mu\nu},
 \qquad \Box=\nabla^2-c^{-2}\partial_t^2.
\]
The retarded solution is
\[
 \bar p_{\mu\nu}(t,x)
 =\frac{4G_N}{c^4}\int_{\mathbb R^3}
 \frac{\delta T_{\mu\nu}(t-|x-y|/c,y)}{|x-y|}\,d^3y.
\]
Trace reversal recovers $p$ from $\bar p$. Denote this entire causal
source-to-metric operation by $\mathcal R_E$.
The equations and retarded solution are the standard weak-field
construction in Flanagan and Hughes,
\href{https://arxiv.org/abs/gr-qc/0501041}{\emph{The basics of gravitational wave theory}},
Sections 2 and 4. A curved or self-gravitating background requires
its appropriate constrained linearized Einstein operator, not an
unqualified reuse of the flat formula. None of these equations
licenses faster-than-light gravitational information transport.

The map $\mathcal S_T$ is a substantive obligation. If it is the
derivative of a quadratic field stress tensor at a zero field,
then that derivative vanishes: the leading response is quadratic
in the field amplitude. A nonzero linear response needs a suitable
background or an input such as a density perturbation. Moreover a
constant global phase of a standard phase-invariant field does not
change its stress tensor. Phase gradients, physical interference or
an explicit interaction are needed before such a phase label becomes
a gravitational source.

\subsection{From the metric to a C5 residual}
Let $\bar\psi_M$ be a reference matter probe with a specified
metric-dependent field operator $\mathcal D_M[g]$. Define
\[
 \mathcal V_M[p]
 :=-(D_g\mathcal D_M)[p]\bar\psi_M,
 \qquad \delta\psi_M
 =\mathcal R_M\mathcal V_M[p],
\]
where $\mathcal R_M$ is its retarded linearized solution operator.
For an illustrative scalar probe
$\mathcal D_M=\Box_g-\mu_M^2$,
\[
 \begin{aligned}
 (D_g\Box_g)[p]\bar\psi_M
 ={}&-p^{\mu\nu}\partial_\mu\partial_\nu\bar\psi_M\\
 &-\left(\partial_\mu p^{\mu\nu}
 -\tfrac12\partial^\nu p\right)\partial_\nu\bar\psi_M,
 \end{aligned}
\]
where $p=\eta^{\mu\nu}p_{\mu\nu}$. The last bracket vanishes
in Lorenz gauge. The source's spinorial probe would require its
own coframe/spin-connection variation; this scalar example does
not replace that calculation.

Let $L_5$ be the derivative of the actual three-component C5 readout
on the chosen matter/witness perturbations, including any declared
witness reconstruction. With a licensed source projector $\Pi_T$,
the candidate is
\[
 \boxed{\mathcal G_{5\leftarrow T}
 =L_5\mathcal R_M\mathcal V_M\mathcal R_E
 \mathcal S_T\Pi_T.}
\]
This is a source-to-geometry-to-probe-to-fit response, not a
species-changing vertex introduced merely by naming a matrix entry
``gravity''. It is linear only about the specified background.
Minimal coupling to a common metric alone does not specify a direct
linear conversion between two independent field species.

If its factors are bounded between the declared source, field and
residual norms on a finite time window, their operator norms give
a product bound for $\mathcal G_{5\leftarrow T}$. Regularity and
boundary/initial data must be chosen so that these mappings actually
exist. A local-in-time readout of retarded factors gives a causal
kernel, schematically
\[
 f_5(t)=\int_{t_0}^t\mathcal K_{5T}(t,s)z_T(s)\,ds.
\]
Such an eliminated response is generally a memory operator; it need
not itself be self-adjoint. A self-adjoint closed-system generator
and a retarded reduced response are different descriptions.

\subsection{A concrete weak-field phase operator}
There is a particularly simple gravitational specialization. Take
two identical stationary matter-clock modes at distinct locations
$x_a,x_b$, with common proper angular frequency $\omega_0$.
In the weak, quasi-static metric
$g_{00}=-(1+2\Phi/c^2)$, their phase-rate corrections in the specified
coordinate time are
\[
 \delta\omega_a=\omega_0\Phi(x_a)/c^2,
 \qquad\delta\omega_b=\omega_0\Phi(x_b)/c^2.
\]
These follow from $d\tau/dt=\sqrt{-g_{00}}
=1+\Phi/c^2+O(\Phi^2/c^4)$ for stationary clocks.
In a slow-source near-zone approximation, declare a source density
perturbation $\delta\rho_T$ and use
\[
 \Phi_T(x)=-G_N\int
 \frac{\delta\rho_T(y)}{|x-y|}\,d^3y.
\]
The density is an effective stress-energy input, not a claim that
the information-facing channel already has an identified material
mass density. For rapidly changing sources one must return to the
retarded metric response rather than use this quasi-static formula.

The differential gravitational frequency is therefore explicitly
\[
 \Delta\omega_g
 =-\frac{\omega_0G_N}{c^2}\int\delta\rho_T(y)
 \left(\frac1{|x_b-y|}-\frac1{|x_a-y|}\right)d^3y.
\]
It has frequency units and is fixed by the source and geometry,
not by a desired C5 ratio. Equal potentials give no differential
effect for these identical clocks. For a point-source perturbation
$\delta M$ at distances $R_a,R_b$ from the two clock locations,
\[
 \Delta\omega_g
 =-\frac{\omega_0G_N\delta M}{c^2}
 \left(\frac1{R_b}-\frac1{R_a}\right).
\]
Swapping the locations reverses its sign, so it can increase or
decrease an existing positive phase separation.

In the visible/history basis
$f=(e_a+e_b)/\sqrt2$, $g=(e_a-e_b)/\sqrt2$, the gravitational
frequency generator is
\[
 H_g=
 \begin{pmatrix}
 \overline{\delta\omega}&-\Delta\omega_g/2\\
 -\Delta\omega_g/2&\overline{\delta\omega}
 \end{pmatrix},
 \qquad
 \overline{\delta\omega}
 =(\delta\omega_a+\delta\omega_b)/2.
\]
It is Hermitian. The common part is a scalar phase, while its
differential part mixes the specified visible and history directions.
For a compatible diagonal background generator and constant source,
the empty-history visible ratio is
\[
 q_g=\cos^2\!\left(
 \frac{(\Delta\omega_0+\Delta\omega_g)\Delta t}{2}\right).
\]
Here $\Delta\omega_0$ is the independently specified background gap.
This is an explicit example of gravitational phase mediation; it
does not identify these two clock modes with the Riemann fit modes.
The source's photon-frame clock is an internal phase/readout reference,
not a physical inertial rest frame or proper-time clock of a photon.
An operational calibration is needed to compare $\Delta t$ with
the dimensionless spectral handoff intervals.

\subsection{A reciprocal finite coupling and the effect of genuine drift}
For a declared finite-mode approximation, let the closed generator
on a visible coordinate $r$ and a licensed projected information or
history coordinate $z$ be
\[
 H_{\rm pair}=\bar\omega I+
 \begin{pmatrix}
 \Delta/2&\kappa e^{i\chi}\\
 \kappa e^{-i\chi}&-\Delta/2
 \end{pmatrix},
 \qquad i\partial_t\begin{pmatrix}r\\z\end{pmatrix}
 =H_{\rm pair}\begin{pmatrix}r\\z\end{pmatrix},
\]
with real $\Delta,\chi$ and $\kappa\ge0$. Entries are angular
frequencies. Coordinates are orthonormal for declared positive
comparison norms; the opposite Lorentzian interval sign of the
information-facing channel does not make its Hilbert norm negative.
The conjugate reverse entry makes the full model
norm-preserving; a one-way transfer with no registered reaction
would not have this property. The coefficients must ultimately be
matched to a mode reduction of the metric/probe response. This
finite reciprocal completion is not automatically identical to
the retarded operator above. Conservation of this comparison norm is
not a derivation of the physical gravitational energy budget.

For constant coefficients, put
$\Omega=\sqrt{\kappa^2+\Delta^2/4}$. Direct exponentiation gives
\[
 U(t)=e^{-i\bar\omega t}
 \left[\cos(\Omega t)I
 -i\frac{\sin(\Omega t)}{\Omega}
 (H_{\rm pair}-\bar\omega I)\right].
\]
With $z_0=0$, the surviving squared visible amplitude is
\[
 q(\Delta t)=1-
 \frac{\kappa^2}{\kappa^2+\Delta^2/4}
 \sin^2\!\left(\sqrt{\kappa^2+\Delta^2/4}\,\Delta t\right).
\]
A constant coupling phase $\chi$ does not change this ratio.
Nonzero incoming data instead produce
\[
 r_1=a r_0+b z_0,\qquad
 |r_1|^2=|a r_0|^2+|b z_0|^2
 +2\operatorname{Re}(\overline{a r_0}\,b z_0),
\]
so a coherent relative phase can enhance or suppress the visible
score. A maintained incoming source needs its own supply/history
account; it is not a cost-free fixed vector in a closed pair.

For a genuine time-dependent coupling phase, use
$W(t)=\operatorname{diag}(e^{i\chi(t)/2},e^{-i\chi(t)/2})$ and
$\widetilde\psi=W^*\psi$. The transformed generator is
\[
 \widetilde H=W^*H_{\rm pair}W-iW^*\dot W,
 \qquad \Delta_{\rm eff}=\Delta+\dot\chi.
\]
Thus a constant physical drift rate can change detuning; the
constant-coefficient formula applies with $\Delta_{\rm eff}$.
For general drift the time-ordered evolution is needed, not this
formula with a casually substituted average phase. A pure change
of coordinates, with the Hamiltonian, readout, clock and metric
consistently co-transported, cannot change the measured score.

\begin{proposition}[Detuning obstruction to quarter transfer]
For a nonzero constant coupling and constant effective detuning,
there exists a sampling interval with $q\le3/4$ if and only if
\[
 |\Delta_{\rm eff}|\le2\sqrt3\,\kappa.
\]
\end{proposition}
\begin{proof}
The largest possible transferred fraction is
$\kappa^2/(\kappa^2+\Delta_{\rm eff}^2/4)$, since the sine squared
can attain $1$. It is at least $1/4$ exactly under the displayed
condition. The condition permits a suitable interval but does not
guarantee that an independently fixed clock uses one.
\end{proof}

For comparison with the existing third-pair probe only, retain its
half-gap $\kappa=1.247617105978133$ and interval
$\Delta t=0.392944430867119$. These are still dimensionless model
units, not gravitationally calculated parameters. The resulting
test is
\[
\begin{array}{c|c|c}
\Delta_{\rm eff}/\kappa&q&\text{maximal transfer}\\
\hline
0&0.778308438&1\\
1&0.782784889&0.8\\
2&0.795784219&0.5\\
4&0.841762894&0.2
\end{array}
\]
The last case cannot meet the quarter-transfer condition at any
interval. The first three can in principle do so at other times,
but do not at the stated probe interval. This proves a mechanism
for slower visible lowering under detuning, not that the source's
third C5 component has the required detuning or gravity coupling.
In general the finite-time loss oscillates; increasing detuning
need not monotonically reduce loss at every chosen time.

\subsection{The HC/QAM guard test is part of the construction}
If the pair is literally the raw $J$-exchanged matter/tachyon pair
and $J$ is the fixed linear swap matrix, the traceless generator obeys
\[
 \|[H_{\rm pair},J]\|_{\rm F}^2
 =2\Delta^2+8\kappa^2\sin^2\chi.
\]
For $\kappa>0$, exact commutation requires $\Delta=0$ and
$\sin\chi=0$. Arbitrary independent shifts of that raw pair
therefore leave the strict fixed-$J$ converter class. A changing
$J$ would need its own covariant transport law and native scope
authorization. An anti-linear realization of $J$ also requires
its own intertwining calculation; the linear-swap test must not
be transferred to it without change.

A different declared realization is possible after a licensed
projection of information into the same HC-even support as the fit.
For example, on $\mathbb C^4$ set
\[
 J=\operatorname{diag}(1,1,-1,-1),\qquad
 C=\operatorname{diag}(1,1,0,0),
\]
\[
 Q=\operatorname{diag}(1,1,1,0),\qquad
 P_4=\operatorname{diag}(0,0,1,0),\qquad
 H=H_{\rm pair}\oplus0_2.
\]
Then $[H,J]=[H,Q]=0$ and $HP_4=P_4H=0$. The pair evolves inside
a licensed declared sector, while a positive protected vector
$s_4=d e_3$, $0<d\le\beta_4$, is fixed. Here $P_4$ denotes the protected
register support, not the source's boundary projector $P_\partial$.
Also $Q\ne C$, and the projected
information coordinate is not the raw kernel $K_2=JK_1$.
The example proves that a guarded projected coupling is
mathematically possible; it does not identify these projectors
with the source's actual ones.

For a native response, no-absorption requires its actual C4
projection to annihilate the matter response to the permitted
source, not merely a copied zero flag. Linearized guard conditions
such as $D\Gamma_5\,\delta W=0$ establish only first-order
tangency. Exact finite updates require domain, collar, relay,
scope and guard invariance on the full orbit. All incoming and
outgoing information must be registered in the QAM lineage.
Fresh history remains necessary for a repeatable constant
visible contraction; the isolated closed pair has backflow.

\paragraph{Construction result and verification.}
There is now a specified causal gravitational response template,
an explicit weak-field source/geometry phase kernel, and a
Hermitian finite coupling with a derived detuning obstruction.
Symbolic calculations checked Hermiticity, the squared traceless
generator, the rotating-frame connection and the fixed-$J$
commutator defect. Finite-mode numerical tests checked the
displayed ratios and conservation with arbitrary complex
incoming data. The separate protected-sector realization was
checked for zero gate commutators and unchanged positive C4 data.
Twelve parameter combinations and $144$ complex input vectors
were checked at $70$ decimal digits; the largest norm discrepancy
was below $2\times10^{-70}$. A direct first-order variation of the
metric wave operator in general $1+1$ dimensions also checked the
scalar response formula, and constant-phase cancellation was
verified for a representative stress bilinear.

The missing native inputs are the admissible stress source and
metric reconstruction, the probe response and actual C5 derivative,
the licensed information projection, the physical/spectral clock
calibration and the guard-compatible mode reduction. Gravity
does not independently force $3/4$. It supplies a concrete way
to test whether physical relative drift or returning information
weakens the fit lowering without hiding the persistent C4 burden.

\section{Testing the existing triadic coupling against phase, metric and C5}
\label{sec:native-triadic-phase-metric-test}

\paragraph{Native starting point and modal scope.}
The source already supplies the triadic two-probe equation
\[
 (i\Gamma^\mu\mathcal D_\mu-\Omega_\triangle)\Psi_\triangle=0,
 \qquad \Psi_\triangle=(\Psi_1,\Psi_2,\Psi_\perp)^T,
\]
with direct $g_Y\mathcal T_{12}$ coupling and the two third-channel
legs $\chi\mathcal T_{1\perp}$ and $\chi\mathcal T_{2\perp}$
(source around lines 7970--8005). In the probe-lensing specialization
the first two fields are ART/QFT regime readouts, not an established
identification with raw matter/tachyon kernels. The orthogonal
component is likewise not identified with the third C5 residual or
with the third Riemann-zero pair.

The matrix $\Omega_\triangle$ is a resonance term in a Dirac-type
equation, not by itself the time-evolution Hamiltonian. A native
mode reduction must specify the derivative term, domain, time
parameter and positive comparison metric before extracting a
generator. Reciprocity and self-adjointness of its modal entries
are additional checks. The following family retains the existing
block pattern but does not claim that this reduction has already
been evaluated in the source. The source's $\chi$ is a coupling
coefficient here, not the phase angle used in the preceding section.

\subsection{Relative phase opens a difference channel}
In orthonormal modal coordinates, declare the Hermitian matrix
\[
 H=\begin{pmatrix}
 \omega_1&g&c_1\\
 \bar g&\omega_2&c_2\\
 \bar c_1&\bar c_2&\omega_\perp
 \end{pmatrix},\qquad \omega_1,\omega_2,\omega_\perp\in\mathbb R.
\]
The entries $c_1,c_2$ represent the reduced third-channel legs.
Set $|+\rangle=(1,1,0)^T/\sqrt2$,
$|-\rangle=(1,-1,0)^T/\sqrt2$ and
$|\perp\rangle=(0,0,1)^T$. Their couplings to the third direction are
\[
 c_+=(c_1+c_2)/\sqrt2,\qquad
 c_-=(c_1-c_2)/\sqrt2.
\]
For equal magnitudes with relative phase $\delta$,
\[
 c_1=k e^{i\delta/2},\qquad c_2=k e^{-i\delta/2},\qquad k>0,
\]
\[
 c_+=\sqrt2 k\cos(\delta/2),\qquad
 c_-=i\sqrt2 k\sin(\delta/2).
\]
Thus $|c_-|^2=2k^2\sin^2(\delta/2)$ is an explicit phase-sensitive
difference strength. This is a coupling calculation, not yet a
gravitational response or a score-loss formula.

\subsection{The fixed-reversal test cannot be omitted}
If, as an additional identification, reversal exchanges the first
two modes and fixes the third, use
\[
 J_s=\begin{pmatrix}0&1&0\\1&0&0\\0&0&1\end{pmatrix},
 \qquad C_s=(I+J_s)/2.
\]
Direct multiplication gives the complete finite-dimensional defect
\[
 \|[H,J_s]\|_{\rm F}^2
 =2(\omega_1-\omega_2)^2
 +8(\operatorname{Im}g)^2+4|c_1-c_2|^2.
\]
With equal diagonal frequencies, real $g$ and the equal-magnitude
phase family, this becomes $16k^2\sin^2(\delta/2)$.
Consequently an active nonzero relative phase generally leaves
the strict fixed-$J_s$ commuting converter class. It must be
registered as a response/debt or realized through a separately
licensed projected-information interface; it cannot be silently
declared a zero-leakage C5 update. This test is conditional on
the specified reversal identification and does not assert that
the ART/QFT readout indices already carry this $J_s$.

A passive basis change is different. If
$W_\delta=\operatorname{diag}(e^{i\delta/2},e^{-i\delta/2},1)$,
co-transporting $H$, $J_s$, the initial state and the visible
projector by $W_\delta$ preserves the score exactly. For example,
changing only $c_1,c_2$ while keeping a nonzero real $g$ fixed is
not that co-transport: its direct-coupling phase would also transform.
Changing the reversal alone to remove a diagnostic while retaining
incompatible state/readout data is not a repair of the native guard.

\subsection{An exact phase-stability test at a fixed benchmark clock}
For a solvable test, set the direct coupling to zero and take equal
diagonal frequencies, subtracting their common scalar phase.
Then $H_\delta^3=K^2H_\delta$ with $K=\sqrt2 k$, and
\[
 U_\delta(t)=I-i\frac{\sin(Kt)}K H_\delta
 +\frac{\cos(Kt)-1}{K^2}H_\delta^2.
\]
Starting at $|+\rangle$ with no initial excitation in the other
two modes, write $a=\delta/2$ and $c=\cos(Kt)$. The three squared
readouts are exactly
\[
 \begin{aligned}
 p_+&=[\sin^2a+\cos^2a\,c]^2,\\
 p_-&=\sin^2a\cos^2a\,(1-c)^2,\\
 p_\perp&=\cos^2a\,\sin^2(Kt),
 \end{aligned}\qquad p_++p_-+p_\perp=1.
\]
The additional difference direction weakens transfer from this
specified visible $|+\rangle$ state. For $0<c<1$,
\[
 p_+=c^2+\tfrac12c(1-c)\delta^2+O(\delta^4).
\]
There is no first-order score change at the symmetric point.

To test stability of an already specified $3/4$ benchmark, fix
$Kt=\pi/6$ once and do not retune it during the phase sweep:
\[
\begin{array}{c|c|c|c}
\delta&p_+&p_-&p_\perp\\
\hline
0&0.750000000&0&0.250000000\\
\pi/12&0.753958680&0.000300592&0.245740728\\
\pi/6&0.765625000&0.001121825&0.233253175\\
\pi/3&0.809134526&0.003365474&0.187500000\\
\pi/2&0.870512702&0.004487298&0.125000000
\end{array}
\]
In particular $p_+(\pi/6)=49/64>3/4$ exactly. This establishes
a sensitivity mechanism, not an independent derivation of the
benchmark interval or a native gravitational explanation of
$0.778308438$ in the earlier zero-pair example. Here $p_+$ is
one declared visible readout, not the source's total C5 score.
The isolated three-mode evolution also has backflow; repeatable
lowering would still need the licensed history construction.

Extending $H_\delta$ by zero on a separate positive C4 register
preserves that register and its norm, but does not erase the new
$J_s$-odd signal $p_-$. Its status must be evaluated independently.
If a meta-gate licenses all three modal directions, meta-gate
commutation can hold while fixed-HC commutation fails. These are
different tests, not interchangeable zero flags.

\subsection{The native coframe supplies a metric susceptibility test}
On a declared closed projected background with nondegenerate coframe,
use the actual source definition
\[
 E^a{}_{\mu}=\operatorname{Re}
 (\bar\Psi\Gamma^a\mathcal D_\mu\Psi),\qquad
 g_{\mu\nu}=\eta_{ab}E^a{}_{\mu}E^b{}_{\nu}.
\]
Let $\zeta=\delta\Psi$, keep the source's fixed Clifford frame,
and allow a connection variation $\delta\mathcal D_\mu$.
Its first variation is
\[
 \begin{aligned}
 \delta E^a{}_{\mu}=\operatorname{Re}\bigl(&
 \bar\zeta\Gamma^a\mathcal D_\mu\Psi
 +\bar\Psi\Gamma^a\mathcal D_\mu\zeta\\
 &+\bar\Psi\Gamma^a(\delta\mathcal D_\mu)\Psi\bigr),
 \end{aligned}
\]
\[
 p_{\mu\nu}:=\delta g_{\mu\nu}
 =\eta_{ab}\bigl(\delta E^a{}_{\mu}E^b{}_{\nu}
 +E^a{}_{\mu}\delta E^b{}_{\nu}\bigr).
\]
This is a concrete derivative of the existing metric map, not a
new freely chosen gravity matrix. The selected closure and
nondegeneracy assumptions remain to be checked on the actual witness;
neither is supplied by the phase sweep or the C5 detector.

For an internal Hermitian phase generator $A(x)$ compatible with
the Dirac adjoint, take $\zeta=iA\Psi$ and write
$\mathcal D_\mu=\partial_\mu+\mathcal B_\mu$. Then
\[
 \begin{aligned}
 \delta E^a{}_{\mu}=\operatorname{Re}\bigl(&
 i\bar\Psi[\Gamma^a,A]\mathcal D_\mu\Psi\\
 &+i\bar\Psi\Gamma^a
 (\partial_\mu A+[\mathcal B_\mu,A])\Psi\\
 &+\bar\Psi\Gamma^a(\delta\mathcal B_\mu)\Psi\bigr).
 \end{aligned}
\]
The source must therefore provide these mixed bilinears and
connection changes to decide whether a relative phase is
gravitationally visible. For a common phase $A=\theta I$,
fixed connection and real current $\bar\Psi\Gamma^a\Psi$, the
displayed coframe change is zero, even for a varying $\theta$.
This is a property of this real-part coframe under the stated
field variation, not a general prohibition of gravitational
effects from physical phase gradients. Such gradients may force
additional field or connection responses through the dynamics.

\begin{proposition}[First-order metric blindness on a symmetric stratum]
Assume a fixed compatible reversal $J$, an even reference field
$J\Psi=\Psi$, even Clifford operators and reference derivative,
and an odd tangent field $J\zeta=-\zeta$. Assume also that
$(\delta\mathcal D_\mu)\Psi$ is odd and that the Dirac pairing
respects the even/odd split. Then $\delta E^a{}_{\mu}=0$ and
$\delta g_{\mu\nu}=0$.
\end{proposition}
\begin{proof}
Every first-variation summand pairs an even field with an odd
field, in one order or the other. Such pairings vanish under
the stated compatibility. The metric derivative then vanishes.
\end{proof}

The equal-diagonal, real-direct-coupling phase family has an odd
generator derivative at $\delta=0$. An initially even field
therefore has an odd first-order dynamical response there.
For a coframe obeying the preceding parity assumptions, the
metric response begins no earlier than second order. Equivalently,
if $H_{-\delta}=J_sH_\delta J_s$ and the complete coframe data
respect that covariance, then $g(-\delta)=g(\delta)$.
An asymmetric reference, mixed Clifford/readout matrices or
additional even connection responses can invalidate this
first-order blindness. The source does not evaluate these
conditions for the actual C5 witness.

The metric-dependent connection also introduces an implicit
consistency problem. Write the derivative abstractly as
$\delta E=\mathcal A_E\zeta+\mathcal B_Ep$ and let
$\mathcal K_E$ be the derivative of $g=\eta EE$. Then
\[
 (I-\mathcal K_E\mathcal B_E)p
 =\mathcal K_E\mathcal A_E\zeta.
\]
A bounded inverse, for example under a proved
$\|\mathcal K_E\mathcal B_E\|<1$ in suitable fixed-gauge norms,
would determine $p$. It is not supplied by the coframe notation.
If the induced metric is also required to obey the effective
Einstein equation, this $p$ must agree with its constrained
retarded response. The two metric constructions cannot be
assigned independently and both called the native gravity map.

\subsection{The last arrow is the actual C5 derivative}
On a specified witness reconstruction, the gravitational branch
would give
\[
 \delta R_5^{\rm grav}
 =L_5\mathcal R_M\mathcal V_M[p],
\]
with the meanings of Section~\ref{sec:metric-mediated-gravity-candidate}.
The source's earlier first-window probe mask called $L_5$ is not
already this three-component derivative. The actual component
functions and their metric/witness dependence remain necessary.

For a twice differentiable native family
$R_5(\varepsilon)=R_0+\varepsilon R_1+
\tfrac12\varepsilon^2R_2+o(\varepsilon^2)$,
\[
 \begin{aligned}
 \Phi_5(\varepsilon)={}&\|R_0\|^2
 +2\varepsilon\operatorname{Re}\langle R_0,R_1\rangle\\
 &+\varepsilon^2\bigl(\|R_1\|^2
 +\operatorname{Re}\langle R_0,R_2\rangle\bigr)
 +o(\varepsilon^2).
 \end{aligned}
\]
If the first-order metric branch is zero, its first-order residual
response is also zero for bounded linear probe/readout factors.
At a positive reference score, a second-order response can still
raise or lower the score through its overlap with $R_0$; its sign
is not fixed by the word ``leakage''. At $R_0=0$ with $R_1=0$,
a nonzero $R_2$ instead gives a fourth-order leading score. A direct phase
response outside the metric branch must be counted separately.

Native execution thus needs the modal generator and reversal
assignment, background probe and connection, coframe compatibility,
clock calibration and actual C5 derivatives. The source's clock
comparison and compressed $\Delta_\gamma<1/5$ selector condition
do not themselves bound $\dot\delta$ or fix the benchmark interval.
The persistent C4 datum remains a separate unspent register.

\paragraph{Executed test and outcome.}
Symbolic calculations checked the symmetric/difference couplings,
complete reversal defect, cubic generator identity and exact
$49/64$ visible benchmark. A phase sweep checked the three output
norms and their unit sum. The coframe susceptibility was checked
directly with a varying relative phase and noncommuting connection
and comparison matrices. The even/odd first-variation cancellation
and passive co-transport invariance were also checked separately.
Six phase cases were evaluated at $65$ decimal digits for the
three-output norm partition, and at $70$ digits for passive
co-transport and a separate positive protected register. The
largest passive-score discrepancy was below $6\times10^{-71}$.

The existing triadic block supports a concrete mechanism for weaker
visible lowering. Under the raw-mirror identification it also opens
a fixed-$J$ compatibility defect. Under a symmetric compatible
coframe it has no first-order gravitational metric response.
Consequently it is not yet a guard-admissible gravitational C5
explanation. The remaining alternatives are a justified asymmetric
or higher-order metric response, or a licensed projected-readout
coupling distinct from unrestricted independent raw frame shifts.

\section{Quadratic memory and a compatibility test for the reconstructed background}
\label{sec:quadratic-memory-background-test}

\paragraph{Scope of the continuation.}
The preceding phase test suggests a second-order response on a symmetric
stratum. There are two distinct questions: whether the existing modal
coupling generates such a response, and whether a background realizing the
source's coframe is admissible for a gravitational calculation. We can
answer the first in the declared modal model and obtain a conditional
compatibility criterion for the second. Neither calculation identifies the
actual C5 residual functions or supplies a solved native field background.
The source specifies the triadic block, the coframe and the effective
action, but does not evaluate the required mode reduction, background
fields, Dirac-adjoint convention or three-component C5 derivatives.

\subsection{Eliminating the difference mode gives quadratic memory}
Use the phase family of Section~\ref{sec:native-triadic-phase-metric-test}
with equal first-two diagonal frequencies $\omega$, real direct coupling
$g$, and constant coefficients. Reorder the parity basis into the even
subspace $\mathcal E=\operatorname{span}\{|+\rangle,|\perp\rangle\}$
and the odd direction $|-\rangle$. Then
\[
 H_\delta=\begin{pmatrix}A_\delta&B_\delta\\B_\delta^*&d\end{pmatrix},
 \quad d=\omega-g,
\]
\[
 A_\delta=\begin{pmatrix}
 \omega+g&\sqrt2 k\cos(\delta/2)\\
 \sqrt2 k\cos(\delta/2)&\omega_\perp
 \end{pmatrix},\qquad
 B_\delta=\begin{pmatrix}0\\-i\sqrt2 k\sin(\delta/2)\end{pmatrix}.
\]
Let $P_\perp=|\perp\rangle\langle\perp|$ on $\mathcal E$. For a
spectral parameter $z\ne d$ away from the full spectrum, block elimination
gives the exact compressed resolvent
\[
 P_{\mathcal E}(z-H_\delta)^{-1}|_{\mathcal E}
 =\bigl(z-A_\delta-\Sigma_\delta(z)\bigr)^{-1},
\]
\[
 \Sigma_\delta(z)
 :=B_\delta(z-d)^{-1}B_\delta^*
 =\frac{2k^2\sin^2(\delta/2)}{z-\omega+g}P_\perp
 =\frac{k^2\delta^2}{2(z-\omega+g)}P_\perp+O(\delta^4).
\]
Thus the first feedback through the eliminated odd direction is
quadratic in the relative phase. There is also a direct quadratic change
of the even coupling:
\[
 A_\delta-A_0
 =-\frac{\sqrt2 k}{8}\delta^2
 \begin{pmatrix}0&1\\1&0\end{pmatrix}+O(\delta^4).
\]
Both contributions belong in a second-order comparison; the self-energy
alone is not the complete phase response. If $|z-d|\ge d_{\rm gap}>0$,
\[
 \|\Sigma_\delta(z)\|
 \le\frac{2k^2\sin^2(\delta/2)}{d_{\rm gap}}.
\]
For real $z$ outside the pole, $\Sigma_\delta(z)$ is Hermitian and its
sign depends on $z-d$. A complex retarded parameter instead describes
a response function, not an extra self-adjoint Hamiltonian. No spectral
gap or modal coefficients have yet been evaluated for the native witness.

\begin{proposition}[Exact causal memory in the declared triad]
Write the modal state as $(u,v)\in\mathcal E\oplus\mathbb C$ and impose
$i\dot u=A_\delta u+B_\delta v$,
$i\dot v=dv+B_\delta^*u$, with initial odd datum $v_0$. Then
\[
 v(t)=e^{-idt}v_0-i\int_0^t e^{-id(t-s)}B_\delta^*u(s)\,ds,
\]
\[
 i\dot u(t)=A_\delta u(t)+B_\delta e^{-idt}v_0
 -i\int_0^t K_\delta(t,s)u(s)\,ds,
\]
where, for $0\le s\le t$,
\[
 K_\delta(t,s)=B_\delta e^{-id(t-s)}B_\delta^*
 =2k^2\sin^2(\delta/2)e^{-id(t-s)}P_\perp.
\]
\end{proposition}
\begin{proof}
Integrating the odd equation with the scalar propagator $e^{-idt}$ gives
the first formula. Substitution into the even equation gives the second
and the displayed kernel. This keeps the initial odd source explicitly.
\end{proof}

The temporal delay here is exact modal memory, not yet propagation of a
spacetime metric. In this finite closed system the kernel is oscillatory,
not a proved irreversible damping term; stored amplitude can return.
Replacing its spectral dependence by a constant local generator would
require a separate approximation with an error bound. At $v_0=0$ the
memory term starts at order $\delta^2$; a fixed nonzero $v_0$ also produces
an order-$\delta$ driving term.

For independently varying phases the same elimination remains available.
With smooth time-dependent modal coefficients in this fixed parity basis,
replace $e^{-id(t-s)}$ by
\[
 U_o(t,s)=\exp\left(-i\int_s^t d(r)\,dr\right).
\]
The exact causal kernel is now
\[
 K(t,s)=2k(t)k(s)\sin\!\frac{\delta(t)}2
 \sin\!\frac{\delta(s)}2\,U_o(t,s)P_\perp,
\]
with initial source $B(t)U_o(t,0)v_0$. For
$\delta(t)=\varepsilon r(t)$, $|r(t)|\le r_{\max}$ and
$|k(t)|\le k_{\max}$ on $[0,T]$, its finite-time bound is
\[
 \left\|\int_0^t K(t,s)u(s)\,ds\right\|
 \le\tfrac12 k_{\max}^2\varepsilon^2r_{\max}^2 T
 \sup_{0\le s\le T}\|u(s)\|.
\]
This captures phase history in the declared model. It is not a
uniform infinite-time lowering estimate, and the phase trajectories
are still input data rather than solutions of a native gravitational law.

The reduced equation acts within the even subspace, where $J_s=I$, but
this does not repair full-space fixed-$J_s$ commutation. The recovered
$v(t)$ still records odd excitation. It cannot be erased from the guard
ledger or identified with absorption of the separate protected C4 datum.
Consequently $\Sigma_\delta$ is a derived modal feedback operator, not
already the native gravitational mediator or a certified C5 update.

\subsection{The real-part coframe has a restrictive standard-spinor interpretation}
Return to the source definition, not an alternative bilinear:
\[
 E^a{}_{\mu}=\operatorname{Re}
 (\bar\Psi\Gamma^a\mathcal D_\mu\Psi),\qquad
 g_{\mu\nu}=\eta_{ab}E^a{}_{\mu}E^b{}_{\nu}.
\]
Here the triadic index may be included in $\Psi$. To test a standard
geometric interpretation, assume all of the following on a smooth,
pointwise closed background: the vector current
$j^a=\bar\Psi\Gamma^a\Psi$ is real; the Dirac pairing is compatible with
the Lorentz spin connection
$\mathcal B^L_\mu=\tfrac14\omega_{\mu bc}\Sigma^{bc}$; the remaining
internal connection is compatible with that pairing and commutes with
the Clifford frame; and $E^a{}_{\mu}$ is exactly the geometric coframe
$e^a{}_{\mu}$ whose torsion-free spin connection is $\omega$.
The corresponding spacetime current is defined by
$j^\mu=e_a{}^\mu j^a$, with $e_a{}^\mu$ the inverse coframe. These
identifications and compatibility assumptions must be verified, not
inferred merely from the names ``probe current'' and ``Einstein limit''.
In particular, identifying an exact microscopic coframe with a
coarse-grained geometric coframe is an additional assumption.

\begin{proposition}[Conditional coframe--current compatibility]
Under the preceding assumptions in dimension $n$, the native real-part
coframe satisfies
\[
 2E^a{}_{\mu}=D^L_\mu j^a
 :=\partial_\mu j^a+\omega_\mu{}^a{}_b j^b.
\]
If $E=e$ is nondegenerate, then
\[
 \nabla_\mu j^\nu=2\delta_\mu{}^\nu,\qquad
 \nabla_\mu j_\nu=2g_{\mu\nu},\qquad
 \nabla_\mu j^\mu=2n.
\]
In four dimensions this divergence is $8$, so these assumptions cannot
also hold for a conserved current $\nabla_\mu j^\mu=0$.
\end{proposition}
\begin{proof}
Reality and compatibility of the vector-current bilinear give
\[
 2\operatorname{Re}(\bar\Psi\Gamma^a\partial_\mu\Psi)
 =\partial_\mu j^a.
\]
For the compatible Lorentz connection,
\[
 2\operatorname{Re}(\bar\Psi\Gamma^a\mathcal B^L_\mu\Psi)
 =\bar\Psi[\Gamma^a,\mathcal B^L_\mu]\Psi
 =\omega_\mu{}^a{}_b j^b.
\]
A compatible Clifford-commuting internal connection has zero real
contribution. Adding the terms proves the first identity. The tetrad
postulate then gives
$\nabla_\mu j^\nu=e_a{}^\nu D^L_\mu j^a=2\delta_\mu{}^\nu$.
Lowering and contracting the index proves the remaining claims.
\end{proof}

The number $8$ refers to the unscaled coframe definition and four
coordinates, not a measured creation rate in physical units. The
obstruction is conditional: the source does not explicitly impose this
spacetime-current identification, a compatible adjoint and conservation
of this particular current. It therefore does not by itself refute the
paper. It does rule out simultaneously assigning all of these usual
properties without further structure. If the bilinear is instead an
effective nonconserved readout, its relation to physical conserved
charge, probability and the complete source balance must be specified.

There is also a restriction on curvature. The identity implies
\[
 \mathcal L_jg=4g,\qquad R^a{}_b\,j^b=0,
\]
where the second equation is an equality of curvature two-forms applied
to $j$. Indeed $D^L(D^Lj)=2D^Le=0$ under vanishing torsion.
Thus $j$ generates a local metric dilation, rather than an arbitrary
conserved fluid flow; the curvature must annihilate this direction.
Curvature need not vanish. In the simpler case with no Lorentz spin
connection and only a compatible Clifford-commuting internal connection,
however,
\[
 E^a=\tfrac12 dj^a=dY^a,\qquad Y^a=j^a/2.
\]
If this coframe has rank four, $Y^a$ are local coordinates and
$g=\eta_{ab}\,dY^a dY^b$ is locally flat. A phase/profile choice alone
cannot generate nonzero curvature in that restricted subcase.

\subsection{An explicit curved realization, but not a solved field background}
The conditional conservation obstruction does not prohibit all
nondegenerate curved coframes. Here is a direct geometric check in
dimensionless coordinates $(t,x^1,x^2,x^3)$ with $t>0$.
Let $\gamma_{\rm std}^a$ be the usual Dirac matrices with Clifford
signature $(+---)$, set $\Gamma^a=i\gamma_{\rm std}^a$, and choose
$\bar\Psi=-\Psi^\dagger\Gamma^0$. Then $\{\Gamma^a,\Gamma^b\}
=2\eta^{ab}I$ with signature $(-+++)$, and $j^a$ is real. Declare
\[
 \Psi=\sqrt{2t}\begin{pmatrix}1\\0\\0\\0\end{pmatrix},\qquad
 \mathcal B_t=0,\qquad
 \mathcal B_i=-\tfrac12\Gamma^0\Gamma^i\quad(i=1,2,3).
\]
These matrices constitute the torsion-free compatible Lorentz spin
connection of $e^0=dt$, $e^i=t\,dx^i$, with zero internal connection.
Direct substitution into the source's bilinear gives
\[
 (E^a{}_{\mu})=\operatorname{diag}(1,t,t,t),\qquad
 g_{\mu\nu}=\operatorname{diag}(-1,t^2,t^2,t^2),\qquad
 j^a=(2t,0,0,0).
\]
It has rank four and nonzero curvature:
\[
 R_{\mu\nu}=\operatorname{diag}(0,2,2,2),\qquad
 R=6/t^2,\qquad
 \nabla_\mu j^\mu=t^{-3}\partial_t(2t^4)=8.
\]
This is a realization of the coframe--connection geometry, with a
nonconserved bilinear current, not a solution of the native probe,
resonance, closure and effective Einstein equations. In particular,
using $\Gamma^\mu=e_a{}^\mu\Gamma^a$ gives
\[
 i\Gamma^\mu\mathcal D_\mu\Psi
 =\frac{2i}{t}\Gamma^0\Psi.
\]
The master equation would require the independently specified native
$\Omega$ to have this action on $\Psi$. Assigning $\Omega$ after the
fact to force equality would not verify that requirement. Likewise,
assigning $T=G/\kappa_{\rm eff}$ would not derive the microscopic source.
No C5 evaluation or full guard audit is provided by this example.

Its lapse is constant: identical comoving clocks have $d\tau=dt$.
Nonzero curvature therefore does not by itself supply the differential
lapse used in the earlier two-clock phase kernel. A gravitational phase
comparison needs specified worldlines, trajectories and clock readouts,
not merely a nonzero scalar curvature.

The geometric identities also admit spatially asymmetric cone metrics
\[
 g=-dt^2+t^2h_{ij}(x)\,dx^i dx^j,\qquad j=2t\partial_t,
\]
for a smooth positive spatial metric $h$. Locally choose an orthonormal
spatial coframe $\vartheta^i$ for $h$ and set $e^i=t\vartheta^i$.
Its Levi--Civita connection has $\omega^i{}_0=\vartheta^i$, so
$D^Lj=2e$. The same spinor profile with that compatible connection
therefore realizes the real-part coframe. This demonstrates geometric
room for asymmetry within the nonconserved-current interpretation,
not a native asymmetric probe/Einstein/C5 solution. Its constant lapse
again does not determine a two-clock detuning.

\subsection{What a gravitational response must preserve}
On a background satisfying the conditional geometric identity, metric
and current perturbations cannot be chosen independently. Write
$p_{\mu\nu}=\delta g_{\mu\nu}$, $v^\mu=\delta j^\mu$, and
$\xi_\nu=\delta j_\nu=p_{\nu\rho}j^\rho+g_{\nu\rho}v^\rho$.
Linearizing $\nabla_\mu j_\nu=2g_{\mu\nu}$ gives
\[
 \nabla_\mu\xi_\nu
 -\tfrac12j^\sigma\bigl(
 \nabla_\mu p_{\nu\sigma}+\nabla_\nu p_{\mu\sigma}
 -\nabla_\sigma p_{\mu\nu}\bigr)=2p_{\mu\nu}.
\]
In particular, with $p=g^{\mu\nu}p_{\mu\nu}$,
\[
 \nabla_\mu v^\mu+\tfrac12 j^\rho\nabla_\rho p=0.
\]
This follows by varying the constant divergence $2n$. A generic
retarded Einstein response need not obey these coframe/current tangent
conditions. The response must solve them together with the linearized
probe equation, metric reconstruction, Einstein constraints and chosen
boundary/gauge conditions. An unconstrained gravitational Green
operator cannot automatically be used as their common solution.

If a genuine non-Clifford-commuting or non-pairing-compatible internal
connection is retained, define its coframe contribution
\[
 Z^a{}_{\mu}:=\operatorname{Re}
 (\bar\Psi\Gamma^a\mathbb A_\mu\Psi),\qquad
 e^a{}_{\mu}=\tfrac12D^L_\mu j^a+Z^a{}_{\mu}.
\]
The same contraction now yields
\[
 \nabla_\mu j^\mu=8-2e_a{}^\mu Z^a{}_{\mu}.
\]
A conserved current would require $e_a{}^\mu Z^a{}_{\mu}=4$.
This is a necessary condition, not a license to insert an arbitrary
repair tensor. Its $Z$ must come from the actual connection and solve
the native dynamics. Under the explicit exact-closure conditions
$[\mathcal M_{16},Q]=0$, $\mathfrak R_C=0$ and $(I-Q)\Psi=0$,
the positive compensator formula gives $\mathcal H_C=0$ and the
off-sector penalty annihilates $\Psi$. Those two stabilization terms
cannot simultaneously be invoked as a nonzero repair on that same
exactly closed background. The projected connection, collective
fields and coarse-graining may need further analysis; none has been
evaluated here to supply $Z$.

Once an admissible joint response $p$ is actually determined, the
gravitational C5 arrow remains
\[
 \delta R_5^{\rm grav}=L_5\mathcal R_M\mathcal V_M[p],
\]
with the actual three-component derivative $L_5$, not just an older
mask sharing its name. If the symmetric first variation is zero,
one must also evaluate the second variations of this entire chain;
the modal $\Sigma_\delta$ alone does not determine them. Their overlap
with the background residual decides the sign of the score change.
Neither curvature, quadratic memory nor current compatibility fixes a
native $3/4$ contraction constant.

\paragraph{Verification and result.}
Exact symbolic calculations checked the modal Schur complement,
Clifford relations, pairing compatibility, displayed coframe, current,
metric, master-equation residual and Ricci/scalar curvature.
The linearized coframe/current equation and its trace were also checked
on a coordinate-deformation family and a spatially asymmetric cone
perturbation; both residuals vanish exactly. These are algebraic
compatibility tests, not evaluations of native C5 or field dynamics.
The causal odd-state reconstruction and the reduced even equation
were independently checked at $55$ decimal digits at three times
with nonzero complex initial odd data; the largest discrepancy was
below $1.1\times10^{-56}$. The numerical coefficients in that check
were diagnostic choices, not recovered native field data.

The existing triadic pattern therefore does provide a derived
quadratic memory operator. The native real-part coframe also admits
curved geometric realizations, but under the stated standard
compatibilities it forces a dilation current incompatible with
current conservation. Selecting the correct current/connection and
coarse-grained interpretation, solving a joint background and
evaluating C5 are still required before calling this a gravitational
OP5 closure mechanism. The source coframe and resonance operator
have not been changed to bypass those requirements.

\section{Can the conserved and geometric currents belong to different frame readouts?}
\label{sec:sector-current-readout-bridge}

\paragraph{Native support and the missing identification.}
The source's triolic-realization corollary explicitly labels
$\Psi_1,\Psi_2,\Psi_3$ as matter, dual tachyon and collapse/light sectors
(around lines 17817--17835). Its three-readout principle identifies the
tachyon branch as the $J$-dual reading of the same branch source
(around lines 18213--18246). The semi-concrete structural master equation
also includes inter-sector $\Lambda_{ab}$ couplings. Finally r82 specifies
a typed $\pi_{T\to M}$ information/coarse-graining interface. These are
genuine starting points for distinct current roles, not evidence that a
change of coordinates changes current conservation.

Two qualifications are essential. The triadic probe-lensing equation
uses ART/QFT regime probes and an orthogonal channel; it is not already
the evaluated matter/tachyon/light field decomposition. The source's
coframe uses the entire closed probe field $\Psi_\triangle$, without an
explicit tachyon-only mask. A map identifying these two field realizations
must therefore be constructed before importing the sector names into
the metric calculation. No native gravitational transfer law or C5
derivative is supplied by these labels.

\subsection{Sector projectors are not the HC parity projectors}
Write the global carrier as
$\Upsilon=\Upsilon_M\oplus\Upsilon_T\oplus\Upsilon_\gamma$, renaming
the source's global field to distinguish it from the auxiliary probe.
A direct-sum realization has canonical coordinate masks $P_a$. In an
additional positive orthogonal comparison realization they obey
\[
 P_a^*=P_a,\qquad P_aP_b=\delta_{ab}P_a,\qquad
 P_M+P_T+P_\gamma=I.
\]
On a licensed projected carrier the last identity is instead the
identity of that carrier, not an assertion about its external complement.
The masks require their own compatibility with the closure domain.

In the three-dimensional diagnostic realization take
\[
 P_M=\operatorname{diag}(1,0,0),\quad
 P_T=\operatorname{diag}(0,1,0),\quad
 P_\gamma=\operatorname{diag}(0,0,1),\quad
 J_s=\begin{pmatrix}0&1&0\\1&0&0\\0&0&1\end{pmatrix}.
\]
Then $J_sP_MJ_s=P_T$, whereas $C_s=(I+J_s)/2$ selects the
symmetric matter/tachyon combination together with the third direction.
Neither $P_M$ nor $P_T$ is $C_s$ or $I-C_s$. In particular, the
tachyon mask must not be identified with forbidden off-$C_7$ leakage
or with the $J_s$-odd difference mode merely because it is a dual label.
The global light sector and the probe's orthogonal mode are identified
only in this diagnostic realization, not in the native equations.

\subsection{An exact projected-current balance}
The adjoint and kinetic term matter as much as the sector labels. On a
common geometric background, suppose
$\bar\Upsilon=\Upsilon^\dagger A_D$ and the current matrices
$K^\mu=A_D\Gamma^\mu$ are Hermitian. Assume the connection obeys the
covariant current Leibniz identity, and take Hermitian internal masks
commuting with $K^\mu$. Multiplication of the probe-type equation by
$A_D$ gives
\[
 iK^\mu\mathcal D_\mu\Upsilon=\mathcal V\Upsilon,
 \qquad \mathcal V=A_D\Omega.
\]
Here this equation is a conditional kinetic realization of the global
carrier, not a claimed evaluation of its structural master equation.
Define real sector currents and the projector derivative by
\[
 j_a^\mu=\Upsilon^\dagger K^\mu P_a\Upsilon,\qquad
 \mathcal D_\mu P_a=\partial_\mu P_a+[\mathcal B_\mu,P_a],
 \quad\mathcal D_\mu=\partial_\mu+\mathcal B_\mu.
\]

\begin{proposition}[Coupling and readout sources]
Under these kinetic and pairing assumptions,
\[
 \begin{aligned}
 \nabla_\mu j_a^\mu={}&
 i\Upsilon^\dagger(\mathcal V^\dagger P_a-P_a\mathcal V)\Upsilon\\
 &+\Upsilon^\dagger K^\mu(\mathcal D_\mu P_a)\Upsilon.
 \end{aligned}
\]
Consequently
\[
 \nabla_\mu j_{\rm tot}^\mu
 =i\Upsilon^\dagger(\mathcal V^\dagger-\mathcal V)\Upsilon,
 \qquad j_{\rm tot}=j_M+j_T+j_\gamma.
\]
If $\mathcal V=\mathcal V^\dagger$, the total current is conserved
and each sector source is
\[
 S_a=i\Upsilon^\dagger[\mathcal V,P_a]\Upsilon
 +\Upsilon^\dagger K^\mu(\mathcal D_\mu P_a)\Upsilon,
 \qquad S_M+S_T+S_\gamma=0.
\]
\end{proposition}
\begin{proof}
Apply the compatible covariant Leibniz rule to
$\Upsilon^\dagger K^\mu(P_a\Upsilon)$. It gives one differentiated
bra, one differentiated ket and the displayed projector-derivative
term. The equation and its Hermitian conjugate replace the first two
terms by $i\Upsilon^\dagger\mathcal V^\dagger P_a\Upsilon$ and
$-i\Upsilon^\dagger P_a\mathcal V\Upsilon$. Summing uses
$\sum_aP_a=I$ and $\mathcal D_\mu\sum_aP_a=0$.
\end{proof}

This derives an exchange source, rather than postulating cancellation
by frame names. It also isolates an assumption still needing a native
check: Hermiticity of $\Omega$ in an unrelated positive Hilbert metric
does not automatically imply $A_D\Omega$ is Hermitian in the chosen
current convention. An incompatible connection adds a further current
defect and cannot silently be treated as compatible. The source must
fix the adjoint, geometric kinetic operator and complete source balance
before this is a physical conservation theorem.

With a parallel mask, a diagonal phase or gravitational potential that
commutes with $P_a$ produces no sector-charge exchange through this
identity. It can still change relative phases or other readouts.
Gravitational information transfer is therefore not automatically a
transfer of matter charge into the tachyon sector.

\subsection{The existing modal coupling realizes the exchange mechanism}
For a Hermitian reduced generator $H$ and a modal state satisfying
$i\dot y=Hy$, define $n_a=y^*P_a(t)y$. Its exact balance is
\[
 \dot n_a=y^*\bigl(\dot P_a+i[H,P_a]\bigr)y.
\]
For fixed coordinate masks and scalar modal legs,
\[
 F_{a\leftarrow b}:=2\operatorname{Im}(\bar y_aH_{ab}y_b),
 \qquad F_{a\leftarrow b}=-F_{b\leftarrow a},
\]
\[
 \begin{aligned}
 \dot n_M&=F_{M\leftarrow T}+F_{M\leftarrow\gamma},\\
 \dot n_T&=-F_{M\leftarrow T}+F_{T\leftarrow\gamma},\\
 \dot n_\gamma&=-F_{M\leftarrow\gamma}-F_{T\leftarrow\gamma}.
 \end{aligned}
\]
Thus the third direction must be counted when it exchanges amplitude;
the two-current cancellation is valid only if the third contribution
has vanished or is independently included in the budget. For
$H_{MT}=k e^{i\theta_g}$ and
$y_a=\sqrt{n_a}e^{i\theta_a}$,
\[
 F_{M\leftarrow T}
 =2k\sqrt{n_Mn_T}\sin(\theta_g+\theta_T-\theta_M).
\]
Relative phase controls a concrete exchange rate, not a universal
constant source or a predetermined contraction factor.

For an explicit isolated pair,
\[
 H=\begin{pmatrix}0&k\\k&0\end{pmatrix},\quad
 y(t)=\begin{pmatrix}\cos(kt)\\-i\sin(kt)\end{pmatrix},
\]
\[
 n_M=\cos^2(kt),\quad n_T=\sin^2(kt),\quad
 \dot n_M=-k\sin(2kt)=-\dot n_T.
\]
The total norm is conserved while the tachyon-labelled readout is not
constant. The generator commutes with the two-slot exchange involution.
An initial matter-only vector is nevertheless not an HC-even state;
generator commutation alone does not certify that initial witness as
native closure-stable. Starting in a strictly HC-even isolated pair
instead gives equal fixed sector populations. The displayed example
is an exchange calculation, not a completed native witness lift.

A passive time-dependent basis change must transport all data. With
$y'=W^*y$, use $P_a'=W^*P_aW$ and
\[
 H'=W^*HW-iW^*\dot W.
\]
Then $n_a'=n_a$ and their physical rates agree. Omitting the connection
term or transporting only the coupling phases can create a spurious
source. Moving projectors are useful for genuine changing readouts,
but cannot turn a nonzero scalar divergence of the same current into
zero by relabeling. A separate positive C4 register remains unspent;
these populations and exchange sources are not the native C5 residuals.

\subsection{Why merely splitting the old coframe does not repair it}
Suppose the geometric probe is the whole global carrier and all the
standard coframe compatibility assumptions of
Section~\ref{sec:quadratic-memory-background-test} hold. Orthogonal
Clifford-compatible masks give
\[
 j_{\rm geom}^a=j_M^a+j_T^a+j_\gamma^a,
 \qquad 2E^a=D^Lj_{\rm geom}^a.
\]
If $e=E$ is nondegenerate, its spacetime current still satisfies
\[
 \nabla_\mu j_{\rm geom}^\mu=8.
\]
Thus a conserved total carrier current cannot also be this same current
under those assumptions. Off-diagonal Hermitian sector couplings can
redistribute the sector divergences, but their sum is zero and cannot
supply the required total $8$. Assigning $8$ to one sector and $-8$ to
another changes the total identity to zero, contradicting the old
full-field coframe identity rather than resolving it. The shared
source must genuinely have two different current roles, or a justified
connection/coarse-graining correction must invalidate the stated
identity. Neither an opposite frame offset nor a signed tachyon scale
provides a negative probability or an automatic compensating source.

\subsection{A conditional bridge with genuinely different current roles}
A precise candidate is to keep the native coframe formula but distinguish
its auxiliary probe input from the conserved global carrier:
\[
 \Psi_\triangle=\mathcal W_T\Upsilon,\qquad
 \mathcal W_T=\mathcal W_TP_T,\qquad
 Q\mathcal W_T=\mathcal W_T,
\]
where $Q=\hat P_{C_7}$ is the probe closure projector, not the sector
mask $P_T$. This proposes a tachyon-facing probe reconstruction, not
a silent replacement of $E$ by a tachyon-only coframe on the original
full field. Its operator $\mathcal W_T$, domain and native provenance
are not supplied by the r82 map. The latter is a typed information
interface and has not been proved linear, current-preserving or
intertwining with either field equation.

If, on the common background, a reconstruction satisfies the stronger
current-intertwining condition
\[
 \mathcal W_T^*K_{\rm probe}^\mu\mathcal W_T
 =K_{\rm global}^\mu P_T,
\]
then $j_{\rm geom}^\mu=j_T^\mu$, whereas the conserved carrier current
is $j_{\rm tot}^\mu$. Under the standard nondegenerate coframe
conditions the required exchange balance is now
\[
 S_T=8,\qquad S_M+S_\gamma=-8,
 \qquad \nabla_\mu j_{\rm tot}^\mu=0.
\]
These are necessary background equations for the independently derived
sources $S_a$, not values one may assign after the calculation. They
are no longer algebraically contradictory because the geometric probe
current is not the total conserved current. In particular, the isolated
matter current is not separately conserved if it alone supplies this
exchange. A conserved effective matter-accessible current would have
to include the induced hidden-sector/interface contribution, or be a
different conserved quantity with its own derivation.

The reduction must keep its source/history. For the modal example,
let $L_T$ be a constant isometric embedding of the tachyon slot and
$z_T=L_TP_Ty$. Then
\[
 i\dot z_T=L_TP_THP_TL_T^*z_T
 +L_TP_TH(I-P_T)y.
\]
The second term is the exact supply from the other slots; a moving
embedding adds $i\partial_t(L_TP_T)y$. Deleting these terms does not
give the induced probe equation. Eliminating hidden slots generates
retarded memory, as in the preceding section. The prepared history
construction can store those degrees of freedom, but does not identify
them with native tachyon data or certify the full QAM guard set. They
must not be taken from the protected C4 register.

\paragraph{A local constructive bridge on the earlier cone geometry.}
The distinction can be realized explicitly at the level of geometry,
compatible carrier kinetics and projected sources. This is an inverse
model construction, not a recovery of the native resonance coefficients.
Keep $g=-dt^2+t^2\sum_i(dx^i)^2$, the Clifford/adjoint convention and
spin connection of Section~\ref{sec:quadratic-memory-background-test}.
Let $u=(1,0,0,0)^T$, choose a positive finite charge budget $N_0$, and
restrict to
\[
 0<t<(N_0/4)^{1/4}.
\]
Define a three-slot carrier by
\[
 \Upsilon_M=\Upsilon_T=-i\sqrt{2t}\,u,\qquad
 \Upsilon_\gamma=t^{-3/2}\sqrt{N_0-4t^4}\,u.
\]
All three sector densities are nonnegative and their currents have no
spatial components. They satisfy exactly
\[
 \begin{aligned}
 j_M^0=j_T^0&=2t,& j_\gamma^0&=N_0/t^3-4t,\\
 j_{\rm tot}^0&=N_0/t^3,&
 (\nabla j_M,\nabla j_T,\nabla j_\gamma)&=(8,8,-16).
 \end{aligned}
\]
Here $\nabla j_a$ abbreviates $\nabla_\mu j_a^\mu$. Since
$\sqrt{|g|}=t^3$, the conserved comoving charge is $N_0$ per unit
coordinate volume. The increasing projected densities are supplied by
the third slot, not by current creation or a negative probability.
The numbers refer to the dimensionless normalization of this example.

For the kinetic realization set
\[
 \theta(t)=\arcsin\frac{2t^2}{\sqrt{N_0}},\qquad
 h(t)=\frac{\dot\theta(t)}{\sqrt2}
 =\frac{2\sqrt2\,t}{\sqrt{N_0-4t^4}},
\]
\[
 H(t)=h(t)\begin{pmatrix}0&0&1\\0&0&1\\1&1&0\end{pmatrix},
 \qquad \Omega_{\rm toy}=iH(t)\otimes I_4.
\]
The volume-rescaled amplitudes
\[
 y(t)=\frac{t^{3/2}}{\sqrt{N_0}}
 \begin{pmatrix}-i\sqrt{2t}\\-i\sqrt{2t}\\
 t^{-3/2}\sqrt{N_0-4t^4}\end{pmatrix}
 =\begin{pmatrix}-i\sin\theta/\sqrt2\\
 -i\sin\theta/\sqrt2\\\cos\theta\end{pmatrix}
\]
obey $i\dot y=Hy$ and $\|y\|=1$. In the stated convention the
compatible spinor equation is also satisfied:
\[
 (i\Gamma^\mu\mathcal D_\mu-\Omega_{\rm toy})\Upsilon=0,
\]
with the Clifford matrices and connection acting identically on each
slot. Indeed on these homogeneous upper spinors its kinetic term is
$-(\partial_t+3/(2t))\Upsilon=iH\Upsilon$.
Moreover $A_D\Omega_{\rm toy}=-iH\otimes\Gamma^0$ is Hermitian.
An ordinary-Hilbert-space check on $\Omega_{\rm toy}$ alone would
miss this convention-dependent reality condition.

An HC-even probe embedding is available using the diagnostic masks:
\[
 \mathcal W_T=\sqrt2\,C_sP_T,\qquad
 \mathcal W_M=\sqrt2\,C_sP_M.
\]
With the spinor identity factor understood, these satisfy
\[
 \mathcal W_T^*\mathcal W_T=P_T,\qquad
 C_s\mathcal W_T=\mathcal W_T,\qquad
 \mathcal W_TJ_s=\mathcal W_M.
\]
Thus the single-sector map is isometric on its input slot and has
HC-even image. It does not intertwine $J_s$ as a full-carrier converter:
$J_s\mathcal W_T=\mathcal W_T$ but
$\mathcal W_TJ_s=\mathcal W_M\ne\mathcal W_T$.
The matter/tachyon readout label must therefore be retained; summing or
forgetting the two input labels is not a reversible two-sector transport.
These are paired readouts, not newly certified native MM operators.

For the carrier above the auxiliary probe is
\[
 \Psi_\triangle=\mathcal W_T\Upsilon
 =(\Upsilon_T/\sqrt2,\Upsilon_T/\sqrt2,0)^T.
\]
It is HC-even and direct evaluation gives
\[
 E_{\rm probe}=\operatorname{diag}(1,t,t,t).
\]
Thus the earlier curved coframe arises from a sourced readout of an
exactly conserved carrier. The source must be retained: the unembedded
sector field obeys
\[
 i\Gamma^\mu\mathcal D_\mu\Upsilon_T
 =ih(t)\Upsilon_\gamma,
\]
not an independently closed homogeneous single-slot equation. Using
the entire $\Upsilon$ instead in the coframe formula would give a
different coframe, not the declared $e$. This example therefore does
not solve the original full-field reconstruction problem by relabeling.

Both $[H,J_s]=0$ and $J_s\Upsilon=\Upsilon$ hold exactly.
However this symmetry also forces the two raw sector currents to be
equal. More generally, with a $J_s$-even carrier, compatible current
matrices and $P_T=J_sP_MJ_s$, one has $j_M=j_T$ and $S_M=S_T$.
It is then impossible for the raw matter current to be conserved while
this same-type raw tachyon current has divergence $8$. In the example
the third sector supplies $-16$. The conserved quantity is the common
carrier current, not the isolated matter-labelled current.

The finite interval matters: $h(t)$ diverges at its upper endpoint as
the third-slot budget is exhausted. No regular all-time continuation
is established. The third slot is a scalar internal amplitude with a
spinor profile in this construction, not a derived Maxwell photon
field. Identity of a model closure projector on its finite carrier,
even parity and a separately protected C4 register do not certify
the native $C_7$ realization, all-readout guards, Einstein/resonance
sources or C5 contraction. Those obligations remain independent.

If the geometry instead uses an information readout that is not
current-intertwining, it need not have the units or conservation law of
physical charge. For example a scalar readout gain $w(x)$ gives
$j_{\rm read}^\mu=|w|^2j_T^\mu$ and
\[
 \nabla_\mu j_{\rm read}^\mu
 =|w|^2S_T+j_T^\mu\partial_\mu|w|^2.
\]
This permits a nonconserved readout without physical charge production,
but a changing gain is not itself a physical gravitational effect.
Normalization covariance and nondegenerate metric reconstruction still
need proof. The gain formula does not establish a rank-four coframe.

There is also a capacity condition if $j_T$ carries a positive part of
the conserved charge. On a no-flux spatial cell,
\[
 \frac{dN_T}{dt}=\int\sqrt{|g|}\,S_T\,d^3x,
 \qquad 0\le N_T\le N_{\rm tot}.
\]
For the earlier cone background $j_T^0=2t$, $\sqrt{|g|}=t^3$ and
constant comoving volume $V$, this would give $N_T=2Vt^4$.
A finite conserved positive budget cannot support that assignment for
all $t>0$ without boundary supply, a changed background or a different
readout interpretation. Local exchange is not an unlimited hidden
reservoir. This bound does not apply to an uncalibrated information
readout as though it were already a physical charge.

\paragraph{Executed checks and outcome.}
Symbolic calculations verified the projected coupling sources,
antisymmetric three-slot flows, their zero sum, the isolated-pair
solution and the distinction between sector masks and HC parity masks.
As a reality-condition diagnostic, replacing $H$ by
$\mathcal V=H-i\lambda I$ gives total source
$-2\lambda\|y\|^2$, rather than sector exchange canceling it.
Fixed and moving projector balances, complete time-dependent phase
co-transport and the Hermiticity of its transformed generator were
checked at $55$ decimal digits at three times. The largest discrepancy
was below $1.8\times10^{-56}$; the symmetric diagnostic triad commutes
with $J_s$ exactly. Its coefficients and initial state were model inputs,
not reconstructed native backgrounds or an HC-stable witness certificate.
Separate exact checks verified the constructive cone example: its
compatible sourced spinor equation, Hermitian current potential,
$J_s$-even carrier, projected coframe and sector divergence partition
$(8,8,-16)$. That example uses a different, explicitly even initial
family and deliberately constructed time-dependent coefficients.
The partial-isometry identity for $\sqrt2 C_sP_T$, its HC-even image,
paired-readout covariance and the two-copy projected coframe were
also verified exactly. Full-carrier $J_s$ intertwining of that readout
fails, as recorded above; it has not been assumed as a guard.

The proposed distinction is therefore structurally possible as a
conserved common-carrier current versus a sourced projected probe or
information readout, with a concrete local geometric/kinetic realization
in the declared model. It is not justified as ``conserved in matter,
nonconserved in tachyon'' for the same covariant current. An explicit
$\mathcal W_T$ derived from the native data, its sourced dynamics,
background compatibility and
physical-current calibration are the next native obligations. Until
they are evaluated, this construction neither repairs the full-field
coframe by relabeling nor proves the gravitational C5 contraction.

\section{Native-data continuation: normalized readout and triadic supply}
\label{sec:native-normalized-readout}

\paragraph{Source audit and level of identification.}
The active source supplies three relevant pieces: the Meta-Septum
projector $Q:=\hat P_{C_7}$ and the probe resonance in
\texttt{thm:master-field-equation}; the operator-valued three-mode
matrix in \texttt{prop:triadic-two-probe}; and the conditional spectral
identity $\mathcal M C=C\mathcal M=\lambda_0C$ in
\texttt{cor:M-HC-fixpoint}, under the hypotheses of
\texttt{thm:hc-spectral-projector}. The first two are in the active
physics-baseline part; the spectral statement is in the structural-core
part. The global master-equation part calls the first two global
branches matter and its $J$-dual tachyon branch, but does not evaluate
their map to the two probe modes.

In particular, the labels ART/QFT/$\perp$ of the probe architecture
do not by themselves identify its mode fibers with the global
matter/tachyon/photon fibers. No evaluated $16\times16$ Millennium
matrix, joint projector realization or branch-to-probe embedding has
been located in these definitions. The following results identify a
normalized candidate from the named projectors and the existing
triadic source leg, conditional on an explicit common realization.
They are not a numerical extraction of a native background.

\subsection{A canonical candidate on the common retained sector}

Assume a specified positive Hilbert carrier, a self-adjoint unitary
involution $J$, $C=(I+J)/2$, and orthogonal projectors $Q$ and $P_T$.
Assume also that $J P_T J=P_M$ with $P_MP_T=0$. These are realization
hypotheses, not consequences of a sector name. Write
\[
 E:=\Pi_{\operatorname{Ran}Q\cap\operatorname{Ran}C},
 \qquad A_T:=EP_T,\qquad G_T:=A_T^\dagger A_T=P_TEP_T,
 \qquad R_T:=\operatorname{supp}G_T.
\]
The common-sector projector is available without assuming commuting
masks:
\[
 E=\mathbf1_{\{1\}}(CQC)
   =\mathop{\mathrm{s\!-\!lim}}_{n\to\infty}(CQC)^n.
\]
Indeed $CQC$ is a positive contraction. Its eigenvalue-one space is
exactly the intersection: $CQCv=v$ implies $Cv=v$ and
$\|Qv\|=\|v\|$, hence $Qv=v$. The spectral theorem then gives the
strong limit. Only when $[Q,C]=0$ can this be replaced by $E=QC$.
Commutation of both masks with $\mathcal M$ does not imply their
commutation with each other.

\begin{proposition}[Normalized common-sector readout]
The polar partial isometry
\[
 W_T:=\mathop{\mathrm{s\!-\!lim}}_{\epsilon\downarrow0}
       A_T(G_T+\epsilon I)^{-1/2}
\]
satisfies
\[
 W_T^\dagger W_T=R_T,\qquad QW_T=CW_T=W_T.
\]
It is isometric on the retained input $R_T\mathcal H$, not necessarily
on the entire tachyon-labelled input $P_T\mathcal H$.
\end{proposition}
\begin{proof}
Use the polar decomposition $A_T=W_TG_T^{1/2}$, with $W_T$ zero on
$\ker G_T$. The regularized positive factor tends strongly to the
support projector. The image of $A_T$, and therefore its closure,
lies in the closed common sector $E\mathcal H$. This gives both
image constraints and the initial projection $R_T$.
\end{proof}

In finite dimensions the formula is $W_T=A_TG_T^{-1/2}$ with the
inverse restricted to the support and extended by zero. For a field
family, pointwise existence does not imply differentiability:
$E(x)$ and $A_T(x)$ must be smooth of locally constant rank, and a
uniform positive bound $G_T\ge g_*R_T$, $g_*>0$, supplies stable
inversion on that support. Rank changes require a separate analysis.

\begin{proposition}[Origin and limits of the square-root factor]
If additionally $[Q,J]=[Q,P_T]=0$, then
\[
 G_T=\tfrac12P_TQ,\qquad R_T=P_TQ,
 \qquad W_T=\sqrt2\,QC P_T.
\]
With $W_M:=\sqrt2\,QC P_M$, one has $W_TJ=W_M$.
\end{proposition}
\begin{proof}
Orthogonality of the exchanged sectors gives $P_TJP_T=0$, so
$P_TCP_T=P_T/2$. The two added commutation conditions allow $Q$
to be moved through this expression. Thus every retained singular
value of $A_T$ is $1/\sqrt2$, and its polar normalization is $\sqrt2$.
Finally $P_TJ=JP_M$ and $CJ=C$ give the paired-readout identity.
\end{proof}

The factor is therefore derived in this compatible realization,
not a freely chosen coupling or an OP5 contraction constant. Without
these hypotheses a fixed $\sqrt2$ need not normalize the map.
Nor is $W_TJ=W_M$ a full-carrier intertwining $W_TJ=JW_T$.
If the masks and $W_T$ act only on internal indices and commute with
the common Dirac-current matrices $K^\mu$, then
\[
 W_T^\dagger K^\mu W_T=K^\mu R_T.
\]
This follows by commuting $K^\mu$ through $W_T$ and using its initial
projection. A positive-Hilbert partial isometry without this kinetic
compatibility does not establish physical-current preservation.

\paragraph{Nonzero overlap is an independent obligation.}
Even a rank-seven orthogonal $Q$ need not retain a common sector.
For an algebraic diagnostic on $\mathbb C^{7}\oplus\mathbb C^{7}
\oplus\mathbb C^2$, take
\[
 J=\begin{pmatrix}0&I_7&0\\I_7&0&0\\0&0&I_2\end{pmatrix},
 \quad C=(I+J)/2,\quad Q=(I-J)/2,\quad \mathcal M=I_{16}.
\]
Then $\operatorname{rank}Q=7$, both closure commutators with
$\mathcal M$ vanish, and $QC=0$. This disproves the implication from
those limited algebraic constraints to a nonzero $W_T$; it is not
asserted to be an admissible native Meta-Septum realization.
Likewise, if a particular spectral reading makes $P_T$ orthogonal
to the compensated photon sector $C$, then $CP_T=0$ and this direct
candidate vanishes. The source's photon-fix-sector language does not
license replacing that reading by an exchanged-pair mask silently.
A nonzero branch embedding or an explicitly different readout would
then have to be supplied.

\subsection{What MM closure does and does not transport}

On a bounded common realization, $[\mathcal M,Q]=[\mathcal M,C]=0$
implies $[\mathcal M,E]=0$. The unnormalized map therefore has the
exact defect
\[
 A_T\mathcal M-\mathcal M A_T=E[P_T,\mathcal M].
\]
On a fixed-rank, gapped support put $N_T=G_T^{-1/2}$, extended by
zero. The normalized defect is
\[
 W_T\mathcal M-\mathcal M W_T
 =E[P_T,\mathcal M]N_T+A_T[N_T,\mathcal M].
\]
These identities follow by expanding the products. An additional
$[\mathcal M,P_T]=0$ makes both terms vanish; MM/HC commutation
alone does not. For unbounded operators, compatible domains are
required before these products or conclusions can be used. A full
operator defect may also vanish on a specified smaller witness
domain; that restriction must be checked rather than inferred.

Under the stronger spectral hypothesis actually stated for
\texttt{cor:M-HC-fixpoint}, namely
$\mathcal M C=C\mathcal M=\lambda_0C$, one obtains
\[
 E Q\mathcal M Q E=\lambda_0 E.
\]
The proof is $QE=CE=E$ and $\mathcal ME=\lambda_0E$.
Consequently the MM term $\alpha Q\mathcal M Q$ has no mixed
channel blocks inside this common sector: it is scalar there.
This does not imply that the whole resonance or kinetic operator is
scalar, nor that it preserves $E$ without further conditions.
The spectral hypothesis is stronger than merely imposing
$[\mathcal M,C]=0$. Its scalar readout cannot on its own supply the
three-channel exchange needed by the preceding conserved-carrier
example; the triadic and connection terms remain available.

\subsection{A missing full-field closure test}

Use the actual probe form $L=i\Gamma^\mu\mathcal D_\mu-\Omega$ and
assume $[\Gamma^\mu,Q]=0$. For a section satisfying $Q\Psi=\Psi$,
the complementary part of its field equation is
\[
 (I-Q)L\Psi=\mathcal K_Q\Psi,
 \qquad
 \mathcal K_Q:=i(I-Q)\Gamma^\mu(\mathcal D_\mu Q)Q
                 -(I-Q)\Omega Q.
\]
Here, if $\mathcal D_\mu=\partial_\mu+B_\mu$, its action on an
endomorphism is $\mathcal D_\mu Q=\partial_\mu Q+[B_\mu,Q]$.
Differentiate $Q\Psi=\Psi$ to obtain
$\mathcal D_\mu\Psi=(\mathcal D_\mu Q)\Psi+
Q\mathcal D_\mu\Psi$, which proves the identity.

In the minimal exact-MM-closure sector, $\mathcal H_C=0$, so the
source resonance gives
\[
 \boxed{\displaystyle
 \mathcal K_Q=
 i(I-Q)\Gamma^\mu(\mathcal D_\mu Q)Q
 -\beta g_Y(I-Q)\mathcal T Q.}
\]
The projected MM term contributes no complementary block because
it is $Q\mathcal M Q$. The scalar mass term contributes none either.
For the displayed native connection, exact closure also reduces
$\mathcal D_\mu Q$ to
$\partial_\mu Q+i[\mathcal A_\mu,Q]$: the penalty
$\eta(I-Q)n_\mu$ commutes with $Q$.

A constrained solution must satisfy $\mathcal K_Q\Psi=0$ in
addition to the projected equation. An autonomous equation on the
entire $Q$-domain requires this compatibility on every admissible
section, together with well-posedness of the restricted evolution.
Strong sufficient algebraic conditions are
$\mathcal D_\mu Q=0$ and $(I-Q)\mathcal T Q=0$. Neither is separately
necessary, because their contributions may cancel. Such cancellation
needs an evaluated connection and coupling, not just a closure label.
Thus $[\mathcal M,Q]=0$ certifies MM block closure, but not by itself
full-field dynamical invariance. This is an additional audit of the
legacy autonomy claim, not a claim that no constrained solution exists.

For example, a constant $Q=\operatorname{diag}(1,1,0)$ and a scalar
$\mathcal M$ commute, while the symmetric coupling with nonzero
entries $\mathcal T_{13}=\mathcal T_{23}=
\mathcal T_{31}=\mathcal T_{32}=1$ has
$(I-Q)\mathcal T Q\ne0$. With $\mathcal DQ=0$ and nonzero
$\beta g_Y$, it fails this compatibility on generic retained states.
Supply from an unobserved mode \emph{inside} $Q$ is different:
it may source a narrower readout without violating $C_7$ closure.

\subsection{The exact field-map defect and the existing supply leg}

For distinct carrier and probe connections define
\[
 \nabla_\mu W:=\partial_\mu W+B^p_\mu W-WB^c_\mu,
 \quad L_c=i\Gamma_c^\mu\mathcal D^c_\mu-\Omega_c,
 \quad L_p=i\Gamma_p^\mu\mathcal D^p_\mu-\Omega_p.
\]
Expanding the covariant derivative gives the exact formula
\begin{align*}
 L_p(W\Upsilon)-WL_c\Upsilon
 &=i(\Gamma_p^\mu W-W\Gamma_c^\mu)
                   \mathcal D^c_\mu\Upsilon\\
 &\quad+
 \bigl(i\Gamma_p^\mu\nabla_\mu W+W\Omega_c-\Omega_pW\bigr)
                   \Upsilon.
\end{align*}
The first term is a principal-symbol defect, not a potential term.
It vanishes only with kinetic intertwining. Even then the second
term sources the probe unless it vanishes on the specified carrier
solutions. For the \emph{constructed} compression
$\Omega_p=W\Omega_cW^\dagger$ on the image and $R=W^\dagger W$,
\[
 W\Omega_c-\Omega_pW=W\Omega_c(I-R).
\]
Compression therefore retains hidden-mode supply; it does not
identify this compressed operator with the independently written
native probe resonance automatically.

To locate that supply in the existing source, suppose the two probe
fibers have been explicitly identified and their exchange acts as
$J_s(\Psi_1,\Psi_2,\Psi_\perp)=(\Psi_2,\Psi_1,\Psi_\perp)$.
Require a common connection preserving this exchange. In this local
three-mode realization put
\[
 G_{12}=g_Y\mathcal T_{12},\quad G_{21}=g_Y\mathcal T_{21},
 \quad V_a=\chi\mathcal T_{a\perp},\quad
 U_a=\chi\mathcal T_{\perp a},
 \quad a=1,2.
\]
The already displayed resonance, not a newly introduced coupling,
is
\[
 \Omega_\triangle=
 \begin{pmatrix}\Omega_1&G_{12}&V_1\\
 G_{21}&\Omega_2&V_2\\U_1&U_2&\Omega_\perp\end{pmatrix}.
\]
In the exchange basis
$a=(\Psi_1+\Psi_2)/\sqrt2$,
$d=(\Psi_1-\Psi_2)/\sqrt2$, $b=\Psi_\perp$, its blocks are
\[
 \Omega'_\triangle=
 \begin{pmatrix}
 \frac{\Omega_1+\Omega_2+G_{12}+G_{21}}2&
 \frac{\Omega_1-\Omega_2-G_{12}+G_{21}}2&
 \frac{V_1+V_2}{\sqrt2}\\[2pt]
 \frac{\Omega_1-\Omega_2+G_{12}-G_{21}}2&
 \frac{\Omega_1+\Omega_2-G_{12}-G_{21}}2&
 \frac{V_1-V_2}{\sqrt2}\\[2pt]
 \frac{U_1+U_2}{\sqrt2}&\frac{U_1-U_2}{\sqrt2}&\Omega_\perp
 \end{pmatrix}.
\]
This basis change does not require the resonance blocks to commute.
Exchange invariance of this matrix requires
\[
 \Omega_1=\Omega_2,\quad G_{12}=G_{21},\quad
 V_1=V_2,\quad U_1=U_2.
\]
On the even subspace $d=0$ the matrix then reduces to
\[
 \Omega_+=
 \begin{pmatrix}
  \Omega_1+g_Y\mathcal T_{12}&\sqrt2\chi\mathcal T_{1\perp}\\
  \sqrt2\chi\mathcal T_{\perp1}&\Omega_\perp
 \end{pmatrix}.
\]
Restrict to a compatible retained fiber where $Q$ acts as identity,
and identify the tachyon-facing mask with the second mode only
as an explicit additional branch-to-probe hypothesis. The normalized
map $\sqrt2 C_sP_2$ sends an even carrier $(a,0,b)$ in exchange
coordinates to $(z,0,0)$, where $z=a/\sqrt2$. Dividing the first
even field equation by $\sqrt2$ proves
\[
 \boxed{\displaystyle
 (i\Gamma^\mu\mathcal D_\mu-
       \Omega_1-g_Y\mathcal T_{12})z
       =\chi\mathcal T_{1\perp}b.}
\]
This is the source-native \emph{symbolic identification} of the
feeding leg that was freely constructed in the earlier cone model.
The twin probe uses the corresponding $\chi\mathcal T_{2\perp}$;
exchange symmetry makes these legs equal, not opposite sources.
It neither evaluates their coefficients nor proves that the global
branch embedding satisfies the assumptions above.

If $b$ belongs to $Q$, this is exchange within a closed carrier and
may yield a nonautonomous projected probe. If $b$ lies outside $Q$,
a nonzero mixed leg requires the preceding $\mathcal K_Q$ test;
with a parallel $Q$ it cannot simply be ignored. Homogeneous
autonomous readout for arbitrary hidden $b$ requires the feeding
leg to vanish, or a justified enlarged state/history description.
The third slot here remains the source's orthogonal probe component,
not a derived physical photon field. A resonance matrix also is not
automatically a self-adjoint time generator: kinetic inversion,
the current pairing and domains must be specified before invoking
conservation or damping.

\paragraph{Executed checks and remaining native obligations.}
Exact symbolic matrix calculations verified the normalized map and
paired covariance on a six-dimensional internal carrier with a
nontrivial commuting $Q$; a noncommuting-mask example verified that
the naive product need not have HC image and that intersection can
lose a tachyon-labelled direction. The rank-seven zero-overlap
example and the closed-MM/nonclosed-triad example were checked
exactly. The full exchange-basis transformation was checked with
nine independent $2\times2$ matrix blocks, without assuming their
commutation. These are algebraic diagnostics, not native measurements.

The advance is a derived normalization, an explicit full-field
compatibility defect, and a supply leg already named in the source.
To upgrade it to a native gravitational OP5 construction still needs
an evaluated joint carrier and branch embedding, nonzero retained
overlap, compatible kinetic/current data, an actual coupled solution
including connection and total stress conservation, and the C5
residual and guard estimates. In particular, neither the scalar
$3/4$ bound nor its failed third benchmark follows from this map.

\paragraph{Transport reference.}
Andre Massing, NTNU, \emph{Weak formulation and well-posedness of the
advection-reaction-diffusion problem},
\href{https://www.math.ntnu.no/emner/MA8502/2024h/lectures/html/chapter_03/weak_form_adr.html}{lecture notes}.
The multichannel time-dependent energy identity and the source and sampling
estimates above are derived explicitly in this note.

\paragraph{Fluid references and verification.}
NASA Glenn Research Center,
\href{https://www1.grc.nasa.gov/beginners-guide-to-aeronautics/navier-strokes-equation/}{Navier--Stokes equations}
and
\href{https://www.grc.nasa.gov/WWW/K-12/airplane/reynolds.html}{Reynolds number}.
The divergence, advection, equation residual and mean squared velocity for
the displayed shear field were independently checked by symbolic
differentiation and integration. The resulting divergence, advection and
equation residual are identically zero. This is an analytic solution check,
not a numerical simulation or a validation of the proposed native fluid map.

\paragraph{Bundle reference.}
Ko Honda, \emph{Differential Geometry}, vector bundles, local trivializations
and sections,
\href{https://www.math.ucla.edu/~honda/math535a/notes.pdf}{author's lecture notes}.
The contraction construction and derivative counterexample above are
explicit calculations for this working note.

\paragraph{Sources for the nonstandard framework.}
H. Jerome Keisler, \emph{Foundations of Infinitesimal Calculus},
Sections 1B, 3C, 9A and 15C,
\href{https://people.math.wisc.edu/~hkeisler/foundations.pdf}{author's monograph};
Isaac Davis, \emph{An Introduction to Nonstandard Analysis},
Corollary 2.12 and Theorem 5.6,
\href{https://www.math.uchicago.edu/~may/VIGRE/VIGRE2009/REUPapers/Davis.pdf}{University of Chicago paper}.
The geometric-sequence and residual-model arguments above are given explicitly
for this working note.

\medskip
\noindent\textbf{Native continuation.} Specify the MM/HC witness domain and
component residual functions; construct its witness update; prove residual
intertwining, domain and guard invariance; then prove native closure-class
preservation before using QAM admissible transport. A native limit requires
its own state topology and admissible limit witness. The scalar or model
limit alone does not establish global all-readout closure.
\end{document}
